\sectiontitle{Section VII:}

                                                               Section VII:
                                                       Sanity and Survival
        In the four chapters of this concluding section, themes of the previous section
are carried further and brought into contact with common social dilemmas and,
eventually, the current world situation. On a small scale, we are constantly faced with
dilemmas like the Prisoner's Dilemma, where personal greed conflicts with social
gain. For any two persons, the dilemma is virtually identical. What would be sane
behavior in such situations? For true sanity, the key element is that each individual
must be able to recognize both that the dilemma is symmetric and that the other
individuals facing it are equally able. Such individuals-individuals who will cooperate
with one another despite all temptations toward crude egoism-are more than just
rational; they are superrational, or for short, sane. But there are dilemmas and "egos"
on a suprahuman level as well. We live in a world filled with opposing belief systems
so similar as to be nearly interchangeable, yet whose adherents are blind to that
symmetry. This description applies not only to myriad small ,conflicts in the world
but also to the colossally blockheaded opposition of the United States and the Soviet
Union. Yet the recognition of symmetry-in short, the sanity-has not yet come. In fact,
the insanity seems only to grow, rather than be supplanted by sanity. What has an
intelligent species like our own done to get itself into this horrible dilemma? What can
it do to get itself out? Are we all helpless as we watch this spectacle unfold, or does
the answer lie, for each one of us, in recognition of our own typicality, and in small
steps taken on an individual level toward sanity?




Dilemmas for Superrational Thinkers,                                                738


                           Dilemmas for
                      Superrational Thinkers,
                   Leading Up to a Luring Lottery
                                        June, 1983


AND then one fine day, out of the blue, you get a letter from S. N. Platonia, well-
known Oklahoma oil trillionaire, mentioning that twenty leading rational thinkers
have been selected to participate in a little game. "You are one of them!" it says.
"Each of you has a chance at winning one billion dollars, put up by the Platonia
Institute for the Study of Human Irrationality. Here's how. If you wish, you may send
a telegram with just your name on it to the Platonia Institute in downtown Frogville,
Oklahoma (pop. 2). You may reverse the charges. If you reply within 48 hours, the
billion is yours-unless there are two or more replies, in which case the prize is
awarded to no one. And if no one replies, nothing will be awarded to anyone."
         You have no way of knowing who the other nineteen participants are; indeed,
in its letter, the Platonia Institute states that the entire offer will be rescinded if it is
detected that any attempt whatsoever has been made by any participant to discover the
identity of, or to establish contact with, any other participant. Moreover, it is a
condition that the winner (if there is one) must agree in writing not to share the prize
money with any other participant at any time in the future. This is to squelch any
thoughts of cooperation, either before or after the prize is given out.
         The brutal fact is that no one will know what anyone else is doing. Clearly,
everyone will want that billion. Clearly, everyone will realize that if their name is not
submitted, they have no chance at all. Does this mean that twenty telegrams will
arrive in Frogville, showing that even possessing transcendent levels of rationality-as
you of course do-is of no help in such an excruciating situation?
         This is the "Platonia Dilemma", a little scenario I thought up recently in trying
to get a better handle on the Prisoner's Dilemma, of which I wrote




Dilemmas for Superrational Thinkers,                                                    739

last month. The Prisoner's Dilemma can be formulated in terms resembling this
dilemma, as follows. Imagine that you receive a letter from the Platonia Institute
telling you that you and just one other anonymous leading rational thinker have been
selected for a modest cash giveaway. As before, both of you are requested to reply by
telegram within 48 hours to the Platonia Institute, charges reversed. Your telegram is
to contain, aside from your name, just the word "cooperate" or the word "defect". If
two "cooperate"s are received, both of you will get $3. If two "defect"s are received,
you both will get $1. If one of each is received, then the cooperator gets nothing and
the defector gets $5.
        What choice would you make? It would be nice if you both cooperated, so
you'd each get $3, but doesn't it seem a little unlikely? After all, who wants to get
suckered by a nasty, low-down, rotten defector who gets $5 for being sneaky?
Certainly not you! So you'd probably decide not to cooperate.
        It seems a regrettable but necessary choice. Of course, both of you, reasoning
alike, come to the same conclusion. So you'll both defect, and that way get a mere
dollar apiece. And yet-if you'd just both been willing to risk a bit, you could have
gotten $3 apiece. What a pity!

                                       *   *   *

       It was my discomfort with this seemingly logical analysis of the "one-round
Prisoner's Dilemma" that led me to formulate the following letter, which I sent out to
twenty friends after having cleared it with Scientific American

  Dear X:
         I am sending this letter out via Special Delivery to twenty of 'you'
  (namely, various friends of mine around the country). I am proposing to all of
  you a one-round Prisoner's Dilemma game, the payoffs to be monetary
  (provided by Scientific American). It's very simple. Here is how it goes.
         Each of you is to give me a single letter: 'C' or 'D', standing for
  'cooperate' or 'defect'. This will be used as your move in a Prisoner's Dilemma
  with each of the nineteen other players. Thepayoff matrix I am using for the
  Prisoner's Dilemma is given in the diagram [see Figure 29-1c].
         Thus if everyone sends in 'C', everyone will get $57, while if everyone
  sends in 'D', everyone will get $19. You can't lose! And of course, anyone who
  sends in 'D' will get at least as much as everyone else will. If, for example, I 1
  people send in 'C' and 9 send in 'D', then the 11 C-ers will get $3 apiece from
  each of the other C-ers (making $30), and zero from the D-ers. So C-ers will get
  $30 each. The D-ers, by contrast, will pick up $5 apiece from each of the C-ers,
  making $55, and $1 from each of the other D-ers, making $8, for a grand total
  of $63. No matter what the distribution is, D-ers always do better than C-ers. Of
  course, the more C-ers there are, the better everyone will do!
         By the way, I should make it clear that in making your choice, you
  should not aim to be the winner, but simply to get as much money for yourself
  as possible. Thus you should be happier to get $30 (say, as a result of saying 'C'
  along with 10 others, even though the 9 D-sayers get more than you) than to get
  $19 (by




Dilemmas for Superrational Thinkers,                                               740

  saying 'D' along with everybody else, so nobody `beats' you). Furthermore, you
  are not supposed to think that at some subsequent time you will meet with and
  be able to share the goods with your co-participants. You are not aiming at
  maximizing the total number of dollars Scientific American shells out, only at
  maximizing the number that come to you!
          Of course, your hope is to be the unique defector, thus really cleaning
  up: with 19 C-ers, you'll get $95 and they'll each get 18 times $3, namely $541
  But why am I doing the multiplication or any of this figuring for you? You're
  very bright. So are all of you! All about equally bright, I'd say, in fact. So all
  you need to do is tell me your choice. I want all answers by telephone (call
  collect, please) the day you receive this letter.
          It is to be understood (it almost goes without saying, but not quite) that
  you are not to try to get in touch with and consult with others who you guess
  have been asked to participate. In fact, please consult with no one at all. The
  purpose is to see what people will do on their own, in isolation. Finally, I would
  very much appreciate a short statement to go along with your choice, telling me
  why you made this particular choice.
                                              Yours....

           P. S. -By the way, it may be helpful for you to imagine a related
  situation, the same as the present one except that you are told that all the other
  players have already submitted their choice (say, a week ago), and so you are
  the last. Now what do you do? Do you submit 'D', knowing full well that their
  answers are already committed to paper? Now suppose that, immediately after
  having submitted your 'D' (or your 'C') in that circumstance, you are informed
  that, in fact, the others really haven't submitted their answers yet, but that they
  are all doing it today. Would you retract your answer? Or what if you knew (or
  at least were told) that you were the first person being asked for an answer?
           And-one last thing to ponder-what would you do if the payoff matrix
  looked as shown in Figure 30-la ?

        FIGURE 30-1. In (a), a modification of Figure 29-1(c). Here, the incentive to
defect seems considerably stronger. In (b), the payoff matrix for a Wolf s-Dilemma
situation involving just two participants. Compare it to that in Figure 29-1(c).




Dilemmas for Superrational Thinkers,                                                741

                                    *      *       *

         I wish to stress that this situation is not an iterated Prisoner's Dilemma
(discussed in last month's column). It is a one-shot, multi-person Prisoner's Dilemma.
There is no possibility of learning, over time, anything about how the others are
inclined to play. Therefore all lessons described last month are inapplicable here,
since they depend on the situation's being iterated. All that each recipient of my letter
could go on was the thought, "There are nineteen people out there, somewhat like me,
all in the same boat, all grappling with the same issues as I am." In other words, there
was nothing to rely on except pure reason.
         I had much fun preparing this letter, deciding who to send it out to,
anticipating the responses, and then receiving them. It was amusing to me, for
instance, to send Special Delivery letters to two friends I was seeing every day,
without forewarning them. It was also amusing to send identical letters to a wife and
husband at the same address.
         Before I reveal the results, I invite you to think how you would play in such a
contest. I would particularly like you to take seriously the assertion "everyone is very
bright". In fact, let me expand on that idea, since I felt that people perhaps did not
really understand what I meant by it. Thus please consider the letter to contain the
following clarifying paragraph:

           All of you are very rational people. Therefore, I hardly need to tell you
   that you are to make what you consider to be your maximally rational choice. In
   particular, feelings of morality, guilt, vague malaise, and so on, are to be
   disregarded. Reasoning alone (of course including reasoning about the others'
   reasoning) should be the basis of your decision. And please always remember
   that everyone is being told this (including this!)!

I was hoping for-and expecting-a particular outcome to this experiment. As I received
the replies by phone over the next several days, I jotted down notes so that I had a
record of what impelled various people to choose as they did. The result was not what
I had expected-in fact, my friends "faked me out" considerably. We got into heated
arguments about the "rational" thing to do, and everyone expressed much interest in
the whole question.
I would like to quote to you some of the feelings expressed by my friends caught in
this deliciously tricky situation. David Policansky opened his call tersely by saying,
"Okay, Hofstadter, give me the $19!" Then he presented this argument for defecting:
"What you're asking us to do, in effect, is to press one of two buttons, knowing
nothing except that if we press button D, we'll get more than if we press button C.
Therefore D is better. That is the essence of my argument. I defect."
Martin Gardner (yes, I asked Martin to participate) vividly expressed the emotional
turmoil he and many others went through. "Horrible dilemma", he said. "I really don't
know what to do about it. If I wanted to maximize




Dilemmas for Superrational Thinkers,                                                 742

my money, I would choose D and expect that others would also; to maximize my
satisfactions, I'd choose C, and hope other people would do the same (by the Kantian
imperative). I don't know, though, how one should behave rationally. You get into
endless regresses: 'If they all do X, then I should do Y, but then they'll anticipate that
and do Z, and so . . .' You get trapped in an endless whirlpool. It's like Newcomb's
paradox." So saying, Martin defected, with a sigh of regret.
        In a way echoing Martin's feelings of confusion, Chris Morgan said, "More by
intuition than by anything else, I'm coming to the conclusion that there's no way to
deal with the paradoxes inherent in this situation. So I've decided to flip a coin,
because I can't anticipate what the others are going to do. I think-but can't know-that
they're all going to negate each other." So, while on the phone, Chris flipped a coin
and "chose" to cooperate.
        Sidney Nagel was very displeased with his conclusion. He expressed great
regret: "I actually couldn't sleep last night because I was thinking about it. I wanted to
be a cooperator, but I couldn't find any way of justifying it. The way I figured it, what
I do isn't going to affect what anybody else does. I might as well consider that
everything else is already fixed, in which case the best I can do for myself is to play a
D."
        Bob Axelrod, whose work proves the superiority of cooperative strategies in
the iterated Prisoner's Dilemma, saw no reason whatsoever to cooperate in a one-shot
game, and defected without any compunctions.
        Dorothy Denning was brief: "I figure, if I defect, then I always do at least as
well as I would have if I had cooperated. So I defect." She was one of the people who
faked me out. Her husband, Peter, cooperated. I had predicted the reverse.

                                        *   *   *

        By now, you have probably been counting. So far, I've mentioned five D's and
two C's. Suppose you had been me, and you'd gotten roughly a third of the calls, and
they were 5-2 in favor of defection. Would you dare to extrapolate these statistics to
roughly 14-6? How in the world can seven individuals' choices have anything to do
with thirteen other individuals' choices? As Sidney Nagel said, certainly one choice
can't influence another (unless you believe in some kind of telepathic transmission, a
possibility we shall discount here). So what justification might there be for
extrapolating these results?
        Clearly, any such justification would rely on the idea that people are "like"
each other in some sense. It would rely on the idea that in complex and tricky
decisions like this, people will resort to a cluster of reasons, images, prejudices, and
vague notions, some of which will tend to push them one way, others the other way,
but whose overall impact will be to push a certain percentage of people toward one
alternative, and another percentage of people toward the other. In advance, you can't
hope to predict what those percentages will be, but given a sample of people in the
situation, you can




Dilemmas for Superrational Thinkers,                                                  743

hope that their decisions will be "typical". Thus the notion that early returns running
5-2 in favor of defection can be extrapolated to a final result of 14-6 (or so) would be
based on assuming that the seven people are acting "typically" for people confronted
with these conflicting mental pressures.
        The snag is that the mental pressures are not completely explicit; they are
evoked by, but not totally spelled out by, the wording of the letter. Each person brings
a unique set of images and associations to each word and concept, and it is the set of
those images and associations that will collectively create, in that person's mind, a set
of mental pressures like the set of pressures inside the earth in an earthquake zone.
When people decide, you find out how all those pressures pushing in different
directions add up, like a set of force vectors pushing in various directions and with
strengths influenced by private or unmeasurable factors. The assumption that it is
valid to extrapolate has to be based on the idea that everybody is alike inside, only
with somewhat different weights attached to certain notions.
        This way, each person's decision can be likened to a "geophysics experiment"
whose goal is to predict where an earthquake will appear. You set up a model of the
earth's crust and you put in data representing your best understanding of the internal
pressures. You know that there unfortunately are large uncertainties in your
knowledge, so you just have to choose what seem to be "reasonable" values for
various variables. Therefore no single run of your simulation will have strong
predictive power, but that's all right. You run it and you get a fault line telling you
where the simulated earth shifts. Then you go back and choose other values in the
ranges of those variables, and rerun the whole thing. If you do this repeatedly,
eventually a pattern will emerge revealing where and how the earth is likely to shift
and where it is rock-solid.
        This kind of simulation depends on an essential principle of statistics: the idea
that when you let variables take on a few sample random values in their ranges, the
overall outcome determined by a cluster of such variables will start to emerge after a
few trials and soon will give you an accurate model. You don't need to run your
simulation millions of times to see valid trends emerging.
        This is clearly the kind of assumption that TV networks make when they
predict national election results on the basis of early returns from a few select towns
in the East. Certainly they don't think that free will is any "freer" in the East than in
the West-that whatever the East chooses to do, the West will follow suit. It is just that
the cluster of emotional and intellectual pressures on voters is much the same all over-
-the nation. Obviously, no individual can be taken as representing the whole nation,
but a well-selected group of residents of the East Coast can be assumed to be
representative of the whole nation in terms of how much they are "pushed" by the
various pressures of the election, so that their choices are likely to show general
trends of the larger electorate.
        Suppose it turned out that New Hampshire's Belknap County and




Dilemmas for Superrational Thinkers,                                                 744

California's Modoc County had produced, over many national elections, very similar
results. Would it follow that one of the two counties had been exerting some sort of
causal influence on the other? Would they have had to be in some sort of eerie cosmic
resonance mediated by "sympathetic magic" for this to happen? Certainly not. All it
takes is for the electorates of the two counties to be similar; then the pressures that
determine how people vote will take over and automatically make the results come
out similar. It is no more mysterious than the observation that a Belknap County
schoolgirl and a Modoc County schoolboy will get the same answer when asked to
divide 507 by 13: the laws of arithmetic are the same the world over, and they operate
the same in remote minds without any need for "sympathetic magic".
        This is all elementary common sense; it should be the kind of thing that any
well-educated person should understand clearly. And yet emotionally it cannot help
but feel a little peculiar since it flies in the face of free will and regards people's
decisions as caused simply by combinations of pressures with unknown values. On
the other hand, perhaps that is a better way to look at decisions than to attribute them
to "free will", a philosophically murky notion at best.

                                        *   *   *

        This may have seemed like a digression about statistics and the question of
individual actions versus group predictability, but as a matter of fact it has plenty to
do with the "correct action" to take in the dilemma of my letter. The question we were
considering is: To what extent can what a few people do be taken as an indication of
what all the people will do? We can sharpen it: To what extent can what one person
does be taken as an indication of what all the people will do? The ultimate version of
this question, stated in the first person, has a funny twist to it: To what extent does my
choice inform me about the choices of the other participants?
        You might feel that each person is completely unique and therefore that no one
can be relied on as a predictor of how other people will act, especially in an intensely
dilemmatic situation. There is more to the story, however. I tried to engineer the
situation so that everyone would have the same image of the situation. In the dead
center of that image was supposed to be the notion that everyone in the situation was
using reasoning alone-including reasoning about the reasoning-to come to an answer.
        Now, if reasoning dictates an answer, then everyone should independently
come to that answer (just as the Belknap County schoolgirl and the Modoc County
schoolboy would independently get 39 as their answer to the division problem).
Seeing this fact is itself the critical step in the reasoning toward the correct answer,
but unfortunately it eluded nearly everyone to whom I sent the letter. (That is why I
came to wish I had included in the letter a paragraph stressing the rationality of the
players.) Once you realize




Dilemmas for Superrational Thinkers,                                                  745

this fact, then it dawns on you that either all rational players will choose D or all
rational players will choose C. This is the crux.
        Any number of ideal rational thinkers faced with the same situation and
undergoing similar throes of reasoning agony will necessarily come up with the
identical answer eventually, so long as reasoning alone is the ultimate justification for
their conclusion. Otherwise reasoning would be subjective, not objective as arithmetic
is. A conclusion reached by reasoning would be a matter of preference, not of
necessity. Now some people may believe this of reasoning, but rational thinkers
understand that a valid argument must be universally compelling, otherwise it is
simply not a valid argument.
        If you'll grant this, then you are 90 percent of the way. All you need ask now
is, "Since we are all going to submit the same letter, which one would be more
logical? That is, which world is better for the individual rational thinker: one with all
C's or one with all D's?" The answer is immediate: "I get $57 if we all cooperate, $19
if we all defect. Clearly I prefer $57, hence cooperating is preferred by this particular
rational thinker. Since I am typical, cooperating must be preferred by all rational
thinkers. So I'll cooperate." Another way of stating it, making it sound weirder, is this:
"If I choose C, then everyone will choose C, so I'll get $57. If I choose D, then
everyone will choose D, so I'll get $19. I'd rather have $57 than $19, so I'll choose C.
Then everyone will, and I'll get $57."

                                        *       *      *

         To many people, this sounds like a belief in voodoo or sympathetic magic, a
vision of a universe permeated by tenuous threads of synchronicity, conveying
thoughts from mind to mind like pneumatic tubes carrying messages across Paris, and
making people resonate to a secret harmony. Nothing could be further from the truth.
This solution depends in no way on telepathy or bizarre forms of causality. It's just
that the statement "I'll choose C and then everyone will", though entirely correct, is
somewhat misleadingly phrased. It involves the word "choice", which is incompatible
with the compelling quality of logic. Schoolchildren do not choose what 507 divided
by 13 is; they figure it out. Analogously, my letter really did not allow choice; it
demanded reasoning. Thus, a better way to phrase the "voodoo" statement would be
this: "If reasoning guides me to say C, then, as I am no different from anyone else as
far as rational thinking is concerned, it will- guide everyone to say C."
         The corresponding foray into the opposite world ("If I choose D, then
everyone will choose D") can be understood more clearly by likening it to a musing
done by the Belknap County schoolgirl before she divides: "Hmm, I'd guess that 13
into 507 is about 49-maybe 39. I see I'll have to calculate it out. But I know in
advance that if I find out that it's 49, then sure as shootin', that Modoc County kid will
write down 49 on his paper as well; and if I get 39 as my answer, then so will he." No
secret transmissions are involved; all that is needed is the universality and uniformity
of arithmetic.




Dilemmas for Superrational Thinkers,                                                  746

Likewise, the argument "Whatever I do, so will everyone else do" is simply a
statement of faith that reasoning is universal, at least among rational thinkers, not an
endorsement of any mystical kind of causality.
        This analysis shows why you should cooperate even when the opaque
envelopes containing the other players' answers are right there on the table in front of
you. Faced so concretely with this unalterable set of C's and D's, you might think,
"Whatever they have done, I am better off playing D than playing C-for certainly
what I now choose can have no retroactive effect on .what they chose. So I defect."
Such a thought, however, assumes that the logic that now drives you to playing D has
no connection or relation to the logic that earlier drove them to their decisions. But if
you accept what was stated in the letter, then you must conclude that the decision you
now make will be mirrored by the plays in the envelopes before you. If logic now
coerces you to play D, it has already coerced the others to do the same, and for the
same reasons; and conversely, if logic coerces you to play C, it has also already
coerced the others to do that.
        Imagine a pile of envelopes on your desk, all containing other people's
answers to the arithmetic problem, "What is 507 divided by 13?" Having hurriedly
calculated your answer, you are about to seal a sheet saying "49" inside your
envelope, when at the last moment you decide to check it. You discover your error,
and change the `4' to a '3'. Do you at that moment envision all the answers inside the
other envelopes suddenly pivoting on their heels and switching from "49" to "39"? Of
course not! You simply recognize that what is changing is your image of the contents
of those envelopes, not the contents themselves. You used to think there were many
"49"s. You now think there are many "39"s. However, it doesn't follow that there was
a moment in between, at which you thought, "They're all switching from `49' to '39'!"
In fact, you'd be crazy to think that.
        It's similar with D's and C's. If at first you're inclined to play one way but on
careful consideration you switch to the other way, the other players obviously won't
retroactively or synchronistically follow you-but if you give them credit for being able
to see the logic you've seen, you have to assume that their answers are what yours is.
In short, you aren't going to be able to undercut them; you are simply "in cahoots"
with them, like it or not! Either all D's, or all C's. Take your pick.
        Actually, saying "Take your pick" is 100 percent misleading. It's not as if you
could merely "pick", and then other people-even in the past-would magically follow
suit! The point is that since you are going to be "choosing" by using what you believe
to be compelling logic, if you truly respect your logic's compelling quality, you would
have to believe that others would buy it as well, which means that you are certainly
not "just picking". In fact, the more convinced you are of what you are playing, the
more certain you should be that others will also play (or have already played) the
same way, and for the same reasons. This holds whether you play C or D, and it is the
real core of the solution. Instead of being a paradox, it's a self-reinforcing solution: a
benign circle of logic.




Dilemmas for Superrational Thinkers,                                                  747

                                         *   *   *

        If this still sounds like retrograde causality to you, consider this little tale,
which may help make it all make more sense. Suppose you and Jane are classical
music lovers. Over the years, you have discovered that you have incredibly similar
tastes in music-a remarkable coincidence! Now one day you find out that two concerts
are being given simultaneously in the town where you live. Both of them sound
excellent to you, but Concert A simply cannot be missed, whereas Concert B is a
strong temptation that you'll have to resist. Still, you're extremely curious about
Concert B, because it features Zilenko Buznani, a violinist you've always heard
amazing things about.
        At first, you're disappointed, but then a flash crosses your mind: "Maybe I can
at least get a first-hand report about Zilenko Buznani's playing from Jane. Since she
and I hear everything through. virtually the same ears, it would be almost as good as
my going if she would go." This is comforting for a moment, until it occurs to you
that something is wrong here. For the same reasons as you do, Jane will insist on
hearing Concert A. After all, she loves music in the same way as you do-that's
precisely why you wish she would tell you about Concert B! The more you feel Jane's
taste is the same as yours, the more you wish she would go to the other concert, so
that you could know what it was like to have gone to it. But the more her taste is the
same is yours, the less she will want to go to it!
        The two of you are tied together by a bond of common taste. And if it turns
out that you are different enough in taste to disagree about which concert is better,
then that will tend to make you lose interest in what she might report, since you no
longer can trust her opinion as that of someone who hears music "through your ears".
In other words, hoping she'll choose Concert B is pointless, since it undermines your
reasons for caring which concert she chooses!
        The analogy is clear, I hope. Choosing D undermines your reasons for doing
so. To the extent that all of you really are rational thinkers, you really will think in the
same tracks. And my letter was supposed to establish beyond doubt the notion that
you are all "in synch"; that is, to ensure that you can depend on the others' thoughts to
be rational, which is all you need.
        Well, not quite. You need to depend not just on their being rational, but on
their depending on everyone else to be rational, and on their depending on everyone
to depend on everyone to be rational-and so on. A group of reasoners in this
relationship to each other I call superrational. Superrational thinkers, by recursive
definition, include in their calculations the fact that they are in a group of
superrational thinkers. In this way, they resemble elementary particles that are
renormalized.
        A renormalized electron's style of interacting with, say, a renormalized photon
takes into account that the photon's quantum-mechanical structure includes "virtual
electrons" and that the electron's quantum-mechanical structure includes "virtual
photons"; moreover it takes into account that all




Dilemmas for Superrational Thinkers,                                                   748

these virtual particles (themselves renormalized) also interact with one another. An
infinite cascade of possibilities ensues but is taken into account in one fell swoop by
nature. Similarly, superrationality, or renormalized reasoning, involves seeing all the
consequences of the fact that other renormalized reasoners are involved in the same
situation-and doing so in a finite swoop rather than succumbing to an infinite regress
of reasoning about reasoning about reasoning ...

                                       *   *   *

        `C' is the answer I was hoping to receive from everyone. I was not so
optimistic as to believe that literally everyone would arrive at this conclusion, but I
expected a majority would-thus my dismay when the early returns strongly favored
defecting. As more phone calls came in, I did receive some C's, but for the wrong
reasons. Dan Dennett cooperated, saying, "I would rather be the person who bought
the Brooklyn Bridge than the person who sold it. Similarly, I'd feel better spending $3
gained by cooperating than $10 gained by defecting."
        Charles Brenner, who I'd figured to be a sure-fire D, took me by surprise and
C'd. When I asked him why, he candidly replied, "Because I don't want to go on
record in an international journal as a defector." Very well. Know, World, that
Charles Brenner is a cooperator!
        Many people flirted with the idea that everybody would think "about the
same", but did not take it seriously enough. Scott Buresh confided to me: "It was not
an easy choice. I found myself in an oscillation mode: back and forth. I made an
assumption: that everybody went through the same mental processes I went through.
Now I personally found myself wanting to cooperate roughly one third of the time.
Based on that figure and the assumption that I was typical, I figured about one third of
the people would cooperate. So I computed how much I stood to make in a field
where six or seven people cooperate. It came out that if I were a D, I'd get about three
times as much as if I were a C. So I'd have to defect. Water seeks out its own level,
and I sank to the lower righthand corner of the matrix." At this point, I told Scott that
so far, a substantial majority had defected. He reacted swiftly: "Those rats-how can
they all defect? It makes me so mad! I'm really disappointed in your friends, Doug."
        So was I, when the final results were in: Fourteen people had defected and six
had cooperated-exactly what the networks would have predicted! Defectors thus
received $43 while cooperators got $15. I wonder what Dorothy's saying to Peter
about now? I bet she's chuckling and saying, "I told you I'd do better this way, didn't
l?" Ah, me ... What can you do with people like that?
        A striking aspect of Scott Buresh's answer is that, in effect, he treated his own
brain as a simulation of other people's brains and ran the simulation enough to get a
sense of what a "typical person" would do. This is very




Dilemmas for Superrational Thinkers,                                                 749

much in the spirit of my letter. Having assessed what the statistics are likely to be,
Scott then did a cool-headed calculation to maximize his profit, based on the
assumption of six or seven cooperators. Of course, it came out in favor of defecting.
In fact, it would have, no matter what the number of cooperators was! Any such
calculation will always come out in favor of defecting. As long as you feel your
decision is independent of others' decisions, you should defect. What Scott failed to
tike into account was that cool-headed calculating people should take into account
that cool-headed calculating people should take into account that cool-headed
calculating people should take into account that ...
        This sounds awfully hard to take into account in a finite way, but actually it's
the easiest thing in the world. All it means is that all these heavy-duty rational
thinkers are going to see that they are in a symmetric situation, so that whatever
reason dictates to one, it will dictate to all. From that point on, the process is very
simple. Which is better for an individual if it is a universal choice: C or D? That's all.

                                        *   *   *

        Actually, it's not quite all, for I've swept one possibility under the rug: maybe
throwing a die could be better than making a deterministic choice. Like Chris
Morgan, you might think the best thing to do is to choose C with probability p and D
with probability 1-p. Chris arbitrarily let p be 1/2, but it could be any number between
0 and 1, where the two extremes represent Ding and C'ing respectively. What value of
p would be chosen by superrational players? It is easy to figure out in a two-person
Prisoner's Dilemma, where you assume that both players use the same value of p. The
expected earnings for each, as a function of p, come out to be I+ 3p -p2, which grows
monotonically as p increases from 0 to 1. Therefore, the optimum value ofp is 1,
meaning certain cooperation. In the case of more players, the computations get more
complex but the answer doesn't change: the expectation is always maximal when p
equals 1. Thus this approach confirms the earlier one, which didn't entertain
probabilistic strategies. - Rolling a die to determine what you'll do didn't add anything
new to the standard Prisoner's Dilemma, but what about the modified-matrix version I
gave in the P. S. to my letter? I'll let you figure that one out for yourself. And what
about the Platonia Dilemma? There, two things are very clear: (1) if you decide not to
send a telegram, your chances of winning are zero; (2) if everyone sends a telegram,
your chances of winning are zero. If you believe that what you choose will be the
same as what everyone else chooses because you are all superrational, then neither of
these alternatives is very appealing. With dice, however, a new option presents itself
to roll a die with probability p of coming up "good" and then to send in your name if
and only if "good" comes up.
        Now imagine twenty people all doing this, and figure out what value of




Dilemmas for Superrational Thinkers,                                                  750

p maximizes the likelihood of exactly one person getting the go-ahead. It turns out
that it is p = 1/20, or more generally, p=1/N where N is the number of participants. In
the limit where N approaches infinity, the chance that exactly one person will get the
go-ahead is 1/e, which is just-under 37 percent. With twenty superrational players all
throwing icosahedral dice, the chance that you will come up the big winner is very
close to 1/(20e), which is a little below two percent. That's not at all bad! Certainly it's
a lot better than zero percent.
         The objection many people raise is: "What if my roll comes up bad? Then why
shouldn't I send in my name anyway? After all, if I fail to, I'll have no chance
whatsoever of winning. I'm no better off than if I had never rolled my die and had just
voluntarily withdrawn!" This objection seems overwhelming at first, but actually it is
fallacious, being based on a misrepresentation of the meaning of "making a decision".
A genuine decision to abide by the throw of a die means that you really must abide by
the throw of the die; if under certain circumstances you ignore the die and do
something else, then you never made the decision you claimed to have made. Your
decision is revealed by your actions, not by your words before acting!
         If you like the idea of rolling a die but fear that your will power may not be up
to resisting the temptation to defect, imagine a third "Policansky button": this one says
`R' for "Roll", and if you press it, it rolls a die (perhaps simulated) and then instantly
and irrevocably either sends your name or does not, depending on which way the die
came up. This way you are never allowed to go back on your decision after the die is
cast. Pushing that button is making a genuine decision to abide by the roll of a die. It
would be easier on any ordinary human to be thus shielded from the temptation, but
any superrational player would have no trouble holding back after a bad roll.

                                         *       *       *

        This talk of holding back in the face of strong temptation brings me to the
climax of this column: the announcement of a Luring Lottery open to all readers and
nonreaders of Scientific American. The prize of this lottery is $1,000,000/N, where N
is the number of entries submitted. Just think: If you are the only entrant (and if you
submit only one entry), a cool million is yours! Perhaps, though, you doubt this will
come about. It does seem a trifle iffy. If you'd like to increase your chances of
winning, you are encouraged to send in multiple entries-no limit! Just send in one
postcard per entry. If you send in 100 entries, you'll have 100 times the chance of
some poor slob who sends in just one. Come to think of it, why should you have to
send in multiple entries separately? Just send one postcard with your name and
address and a positive integer (telling how many entries you're making) to:




Dilemmas for Superrational Thinkers,                                                   751

       Luring Lottery
        c/o Scientific American
           415 Madison Avenue
             New York, N.Y. 10017

You will be given the same chance of winning as if you had sent in that number of
postcards with `1' written on them. Illegible, incoherent, ill-specified, or
incomprehensible entries will be disqualified. Only entries received by midnight June
30, 1983 will be considered. Good luck to you (but certainly not to any-other reader
of this column)!


Post Scriptum.
        The emotions churned up by the Prisoner's Dilemma are among the strongest I
have ever encountered, and for good reason. Not only is it a wonderful intellectual
puzzle, akin to some of the most famous paradoxes of all time, but also it captures in a
powerful and pithy way the essence of a myriad deep and disturbing situations that we
are familiar with from life. Some are choices we make every day; others are the kind
of agonizing choices that we all occasionally muse about but hope the world will
never make us face.
        My friend Bob Wolf, a mathematician whose specialty is logic, adamantly
advocated choosing D in the case of the letters I sent out. To defend his choice, he
began by saying that it was clearly "a paradox with no rational solution", and thus
there was no way to know what people would do. Then he said, "Therefore, I will
choose D. I do better that way than any other way." I protested strenuously: "How
dare you say `therefore' when you've just gotten through describing this situation as a
paradox and claiming there is no rational answer? How dare you say logic is forcing
an answer down your throat, when the premise of your `logic' is that there is no
logical answer?" I never got what I considered a satisfactory answer from Bob,
although neither of us could budge the other. However, I did finally get some insight
into Bob's vision when he, pushed hard by my probing, invented a situation with a
new twist to it, which I call "Wolf's Dilemma".
        Imagine that twenty people are selected from your high school graduation
class, you among them. You don't know which others have been selected, and you are
told they are scattered all over the country. All you know is that they are all connected
to a central computer. Each of you is in a little cubicle, seated on a chair and facing
one button on an otherwise blank wall. You are given ten minutes to decide whether
or not to push your button. At the end of that time, a light will go on for ten seconds,
and while it is on, you may




Dilemmas for Superrational Thinkers,                                                 752

either push or refrain from pushing. All the responses will then go to the central
computer, and one minute later, they will result in consequences. Fortunately, the
consequences can only be good. If you pushed your button, you will get $100, no
strings attached, emerging from a small slot below the button. If nobody pushed their
button, then everybody will get $1,000. But if there was even a single button-pusher,
the refrainers will get nothing at all.
        Bob asked me what I would do. Unhesitatingly, I said, "Of course I would not
push the button. It's obvious!" To my amazement, though, Bob said he'd push the
button with no qualms. I said, "What if you knew your co-players were all logicians?"
He said that would make no difference to him. Whereas I gave credit to everybody for
being able to see that it was to everyone's advantage to refrain, Bob did not. Or at
least he expected that there is enough "flakiness" in people that he would prefer not to
rely on the rationality of nineteen other people. But of course in assuming the
flakiness of others, he would be his own best example-ruining everyone else's chances
of getting $1,000.
        What bothered me about Wolf's Dilemma was what I have come to call
reverberant doubt. Suppose you are wondering what to do. At first it's obvious that
everybody should avoid pushing their button. But you do realize that among twenty
people, there might be one who is slightly hesitant and who might waver a bit. This
fact is enough to worry you a tiny bit, and thus to make you waver, ever so slightly.
But suddenly you realize that if you are wavering, even just a tiny bit, then most likely
everyone is wavering a tiny bit. And that's considerably worse than what you'd
thought at first-namely, that just one person might be wavering. Uh-oh! Now that you
can imagine that everybody is at least contemplating pushing their button, the
situation seems a lot more serious. In fact, now it seems quite probable that at least
one person will push their button. But if that's the case, then pushing your own button
seems the only sensible thing to do. As you catch yourself thinking this thought, you
realize it must be the same as everyone else's thought. At this point, it becomes
plausible that the majority of participants -possibly even all-will push their button!
This clinches it for you, and so you decide to push yours.
        Isn't this an amazing and disturbing slide from certain restraint to certain
pushing? It is a cascade, a stampede, in which the tiniest flicker of a doubt has
become amplified into the gravest avalanche of doubt. That's what I mean by
"reverberant doubt". And one of the annoying things about it is that the brighter you
are, the more quickly and clearly you see what there is to fear. A bunch of amiable
slowpokes might well be more likely to unanimously refrain and get the big payoff
than a bunch of razor-sharp logicians who all think perversely recursively
reverberantly. It's that "smartness" to see that initial flicker of a doubt that triggers the
whole avalanche and sends rationality a-tumblin' into-the abyss. So, dear reader . . . if
you push that button in front of you, do you thereby lose $900 or do you thereby gain
$100?




Dilemmas for Superrational Thinkers,                                                    753

                                       *       *      *

        Wolf's Dilemma is not the same as the Prisoner's Dilemma. In the Prisoner's
Dilemma, pressure towards defection springs from hope for asymmetry (i.e., hope that
the other player might be dumber than you and thus make the opposite choice)
whereas in Wolf's Dilemma, pressure towards button-pushing springs from fear of
asymmetry (i.e., fear that the other player might be dumber than you and thus make
the opposite choice). This difference shows up clearly in the games' payoff matrices
for the two-person case (compare Figure 30-lb with Figure 29-1c). In the Prisoner's
Dilemma, the temptation T is greater than the reward R (5 > 3), whereas in Wolf's
Dilemma, R is greater than T (1,000 > 100).
        Bob Wolf's choice in his own dilemma revealed to me something about his
basic assessment of people and their reliability (or lack thereof). Since his adamant
decision to be a button-pusher even in this case stunned me, I decided to explore that
cynicism a bit more, and came up with this modified Wolf's Dilemma.
        Imagine, as before, that twenty people have been selected from your high
school graduation class, and are escorted to small cubicles with one button on the
wall. This time, however, each of you is strapped into a chair, and a device containing
a revolver is attached to your head. Like it or not, you are now going to play Russian
roulette, the odds of your death to be determined by your choice. For anybody who
pushes their button, the odds of survival will be set at 90 percent-only one chance in
ten of dying. Not too bad, but given that there are twenty of you, it means that almost
certainly one or two of you will die, possibly more. And what happens to the
refrainers? It all depends on how many of them there are. Let's say there are N
refrainers. For each one of them, their chance of being shot will be one in N2. For
instance, if five people don't push, each of them will have only a 1/25 chance of
dying. If ten people refrain, they will each get a 99 percent chance of survival. The
bad cases are, of course, when nearly everybody pushes their button ("playing it safe",
so to speak), leaving the refrainers in a tiny minority of three, two, or even one. If
you're the sole refrainer, it's curtains for you-one chance in one of your death. Bye-
bye! For two refrainers, it's one chance in four for each one. That means there's nearly
a 50 percent chance that at least one of the two will perish.
        Clearly the crossover line is between three and four refrainers. If you have a
reasonable degree of confidence that at least three other people will hold back, you
should definitely do so yourself. The only problem is, they're all making their
decisions on the basis of trying to guess how many people will refrain, too! It's
terribly circular, and you hardly know where to start. Many people, sensing this, just
give up, and decide to push their button. (Actually, of course, how do I know? I've
never seen people in such a situation-but it seems that way from evidence of real-life
situations resembling this, and of course from how people respond to a mere
description of this situation,




Dilemmas for Superrational Thinkers,                                                754

where they aren't really faced with any dire consequences at all. Still, I tend to believe
them, by and large.) Calling such a decision "playing it safe" is quite ironic, because if
only everybody "played it dangerous", they'd have a chance of only one in 400 of
dying! So I ask you: Which way is safe, and which way dangerous? It seems to me
that this Wolf Trap epitomizes the phrase "We have nothing to fear but fear itself."
        Variations on Wolf's Dilemma include some even more frightening and
unstable scenarios. For instance, suppose the conditions are that each button-pusher
has a 50 percent chance of survival, but if there is unanimous refraining from pushing
the button, everyone's life will be spared-and as before, if anyone pushes their button,
all refrainers will die. You can play around with the number of participants, the
survival chance, and so on. Each such variation reveals a new facet of grimness.
These visions are truly horrific, yet all are just allegorical renditions of ordinary life's
decisions, day in, day out.

                                         *   *   *

       I had originally intended to close the column with the following paragraph, but
was dissuaded from it by friends and editors:

          I am sorry to say that I am simply inundated with letters from well-
   meaning readers, and I have discovered, to my regret, that I can barely find time
   to read all those letters, let alone answer them. I have been racking my brains
   for months trying to come up with some strategy for dealing with all this
   correspondence, but frankly I have not found a good solution yet. Therefore, I
   thought I would appeal to the collective genius of you-all out there. If you can
   think of some way for me to ease the burden of my correspondence, please send
   your idea to me. I shall be most grateful.




Dilemmas for Superrational Thinkers,                                                   755


                           Irrationality Is the
                         Square Root of All Evil
                                    September, 1983


THE Luring Lottery, proposed in my June column, created quite a stir. Let me
remind you that it was open to anyone; all you had to do was submit a postcard with a
clearly specified positive integer on it telling how many entries you wished to make.
This integer was to be, in effect, your "weight" in the final drawing, so that if you
wrote "100", your name would be 100 times more likely to be drawn than that of
someone who wrote 'I'. The only catch was that the cash value of the prize was
inversely proportional to the sum of all the weights received by June 30. Specifically,
the prize to be awarded was $1,000,0001W, where W is the sum of all the weights
sent in.
         The Luring Lottery was set up as an exercise in cooperation versus defection.
The basic question for each potential entrant was: "Should I restrain myself and
submit a small number of entries, or should I `go for it' and submit a large number?
That is, should I cooperate, or should I defect?" Whereas in previous examples of
cooperation versus defection there was a clear-cut dividing line between cooperators
and defectors, here it seems there is a continuum of possible answers, hence of
"degree of cooperation". Clearly one can be an extreme cooperator and voluntarily
submit nothing, thus in effect cutting off one's nose to spite one's face. Equally
clearly, one can be an extreme defector and submit a giant number of entries, hoping
to swamp everyone else out but destroying the prize in so doing. However, there
remains a lot of middle ground between these two extremes. What about someone
who submits two entries, or one? What about someone who throws a six-sided die to
decide whether or not to send in a single entry? Or a million-sided die?
         Before I go further, it would be good for me to present my generalized and
nonmathematical sense of these terms "cooperation" and "defection". As a child, you
undoubtedly often encountered adults who admonished you




Irrationality Is the Square Root of All Evil                                       756

for walking on the grass or for making noise, saying "Tut, tut, tut just think if
everyone did that!" This is the quintessential argument used against the defector, and
serves to define the concept:

   A defection is an action such that, if everyone did it, things would clearly be
   worse (for everyone) than if everyone refrained from doing it, and yet which
   tempts everyone, since if only one individual (or a sufficiently small number)
   did it while others refrained, life would be sweeter for that individual (or select
   group).

Cooperation, of course, is the other side of the coin: the act of resisting temptation.
However, it need not be the case that cooperation is passive while defection is active;
often it is the exact opposite: The cooperative option may be to participate
industriously in some activity, while defection is to lay back and accept the sweet
things that result for everybody from the cooperators' hard work. Typical examples of
defection are:

   • loudly wafting your music through the entire neighborhood on a fine
     summer's day;
   • not worrying about speeding through a four-way stop sign, figuring that the
     people going in the crosswise direction will stop anyway;
   • not being concerned about driving a car everywhere, figuring that there's no
     point in making a sacrifice when other people will just continue to guzzle gas
     anyway;
   • not worrying about conserving water in a drought, figuring "Everyone else
     will";
   • not voting in a crucial election and excusing yourself by saying "One vote
     can't make any difference";
   • not worrying about having ten children in a period of population explosion,
     leaving it to other people to curb their reproduction;
   • not devoting any time or energy to pressing global issues such as the arms
     race, famine, pollution, diminishing resources, and so on, saying "Oh, of
     course I'm very concerned-but there's nothing one person can do."

   When there are large numbers of people involved, people don't realize that their
own seemingly highly idiosyncratic decisions are likely to be quite typical and are
likely to be recreated many times over, on a grand scale; thus, what each couple feels
to be their own isolated and private decision (conscious or unconscious) about how
many children to have turns into a population explosion. Similarly, "individual"
decisions about the futility of working actively toward the good of humanity amount
to a giant trend of apathy, and this multiplied apathy translates into insanity at the
group level. In a word, apathy at the individual level translates into insanity at the
mass level.

                                    •          *   *




Irrationality Is the Square Root of All Evil                                         757

        Garrett Hardin, an evolutionary biologist, Wrote a famous article about this
type of phenomenon, called "The Tragedy of the Commons". His view was that there
are two types of rationality: one (I'll call it the "local" type) that strives for the good of
the individual, the other (the "global" type) that strives for the good of the group; and
that these two types of rationality are in an inevitable and eternal conflict. I would
agree with his assessment, provided the individuals are unaware of their joint plight
but are simply blindly carrying out their actions as if in isolation.
        However, if they are fully aware of their joint situation, and yet in the face of it
they blithely continue to act as if their situation were not a communal one, then I
maintain that they are acting totally irrationally. In other words, with an enlightened
citizenry, "local" rationality is not rational, period. It is damaging not just to the
group, but to the individual. For example, people who defected in the One-Shot
Prisoner's Dilemma situation I described in June did worse than if all had cooperated.
        This was the central point of my June column, in which I wrote about
renormalized rationality, or superrationality. Once you know you are a typical
member of a class of individuals, you must act as if your own individual actions were
to be multiplied manyfold, because they inevitably will be. In effect, to sample
yourself is to sample the field, and if you fail to do what you wish the rest would do,
you will be very disappointed by the rest as well. Thus it pays a lot to reflect carefully
about one's situation in the world before defecting, that is, jumping to do the naively
selfish act. You had better be prepared for a lot of other people copping out as well,
and offering the same flimsy excuse.
        People strongly resist seeing themselves as parts of statistical phenomena, and
understandably so, because it seems to undermine their sense of free will and
individuality. Yet how true it is that each of our "unique" thoughts is mirrored a
million times over in the minds of strangers! Nowhere was this better illustrated than
in the response to the Luring Lottery. It is hard to know precisely what constitutes the
"field", in this case. It was declared universally open, to readers and nonreaders alike.
However, we would be safe in assuming that few nonreaders ever became aware of it,
so let's start with the circulation of Scientific American, which is about a million.
Most of them, however, probably did no more than glance over my June column, if
that; and of the ones who did more than that (let's say 100,000), still only a fraction-
maybe one in ten-read it carefully from start to finish. I would thus estimate that there
were perhaps 10,000 people motivated enough to read it carefully and to ponder the
issues seriously. In any case, I'll take this figure as the population of the "field".
        In my June column, I spelled out plainly, for all to see, the superrational
argument that applies to the Platonia Dilemma, for rolling an N-sided die and entering
only if it came up on the proper side. Here, a similar argument goes through. In the
Platonia Dilemma, where more than one entry is fatal to all, the ideal die turned out to
have N faces, where N is the number of




Irrationality Is the Square Root of All Evil                                             758

players-hence, with 10,000 players, a 10,000-sided die. In the Luring Lottery, the
consequences aren't so drastic if more than one entry is submitted. Thus, the ideal
number of faces on the die turns out to be about 2/3 as many-in the case of 10,000
players, a 6,667-sided die would do admirably. Giving the die fewer than 10,000 sides
of course slightly increases each player's chance of sending in one entry. This is to
make it quite likely that at least one entry will arrive!
       With 6,667 faces on the die, each superrational player's chance of winning is
not quite 1 in 10,000, but more like 1 in 13,000; this is because there is about a 22
percent chance that no one's die will land right, so no one will send in any entry at all,
and no one will win. But if you give the die still fewer faces-say 3,000-the expected
size of the pot gets considerably smaller, since the expected number of entrants
grows. And if you give it more faces -say 20,000-then you run a considerable risk of
having no entries at all. So there's a trade-off whose ideal solution can be calculated
without too much trouble, and 6,667 faces turns out to be about optimal. With that
many faces, the expected value of the pot is maximal: nearly $520,000-not to be
sneered at.
       Now this means that had everyone followed my example in the June column, I
would probably have received a total of one or two postcards with ' 1' written on
them, and one of those lucky people would have gotten a huge sum of money! But do
you think that is what happened? Of course not! Instead, I was inundated with
postcards and letters from all over the world -over 2,000 of them. What was the
breakdown of entries? I have exhibited part of it in a table, below:

       1: 1,133
       2: 31
       3: 16
       4: 8
       5: 16
       6: 0
       7: 9
       8: 1
       9: 1
       10: 49
       100: 61
       1,000: 46
       1,000,000: 33
       1,000,000,000:      11
       602,300,000,000,000,000,000,000 (Avogadro's number): 1
       10100 (a googol): 9
       1010100 (a googolplex): 14




Irrationality Is the Square Root of All Evil                                          759

         Curiously, many if not most of the people who submitted just one entry patted
themselves on the back for being "cooperators". Hogwash! The real cooperators were
those among the 10,000 or so avid readers who calculated the proper number of faces
of the die, used a random-number table or something equivalent, and then-most
likely-rolled themselves out. A few people wrote to tell me they had rolled themselves
out in this way. I appreciated hearing from them. It is conceivable, just barely, that
among the thousand-plus entries of `1' there was one that came from a lucky
superrational cooperator-but I doubt it. The people who simply withdrew without
throwing a die I would characterize as well-meaning but a bit lazy, not true
cooperators-something like people who simply contribute money to a political cause
but then don't want to be bothered any longer about it. It's the lazy way of claiming
cooperation.
         By the way, I haven't by any means finished with my score chart. However, it
is a bit disheartening to try to relate what happened. Basically, it is this. Dozens and
dozens of readers strained their hardest to come up with inconceivably large numbers.
Some filled their whole postcard with tiny `9's, others filled their card with rows of
exclamation points, thus creating iterated factorials of gigantic sizes, and so on. A
handful of people carried this game much further, recognizing that the optimal
solution avoids all pattern (to see why, read Gregory Chaitin's article "Randomness
and Mathematical Proof"), and consists simply of a "dense pack" of definitions built
on definitions, followed by one final line in which the "fanciest" of the definitions is
applied to a relatively small number such as 2, or better yet, 9.
         I received, as I say, a few such entries. Some of them exploited such powerful
concepts of mathematical logic and set theory that to evaluate which one was the
largest, became a very serious problem, and in fact it is not even clear that I, or for
that matter anyone else, would be able to determine which is the largest integer
submitted. I was strongly reminded of the lunacy and pointlessness of the current
arms race, in which two sides vie against each other to produce arsenals so huge that
not even teams of experts can meaningfully say which one is larger-and meanwhile,
all this monumental effort is to the detriment of everyone.

                                        *   *   *

        Did I find this amusing? Somewhat, of course. But at the same time, I found it
disturbing and disappointing. Not that I hadn't expected it. Indeed, it was precisely
what I had expected, and it was one reason I was so sure the Luring Lottery would be
no risk for the magazine.
        This short-sighted race for "first place" reveals the way in which people in a
huge crowd erroneously consider their own fancies to be totally unique. I suspect that
nearly everyone who submitted a number above 1,000,000 actually believed they
were going to be the only one to do so. Many of those who submitted numbers such as
a googolplex, or a `9' followed by




Irrationality Is the Square Root of All Evil                                        760

thousands of factorial signs, explicitly indicated that they were pretty sure that they
were going to "win". And then those people who pulled out all the stops and sent in
definitions that would boggle most mathematicians were very sure they were going to
win. As it turns out, I don't know who won, and it doesn't matter, since the prize is
zero to such a good approximation that even God wouldn't know the difference.
Well, what conclusion do I draw from all this? None too serious, but I do hope that it
will give my readers pause for thought next time they face a "cooperate-or-defect"
decision, which will likely happen within minutes for each of you, since we face such
decisions many times each day. Some of them are small, but some will have
monumental repercussions. The globe's future is in your hands-and yes, I mean you
(as well as every other reader of this column).

                                        *   *   *

        And with this perhaps sobering conclusion, I would like to draw my term as a
columnist for Scientific American to a close. It has been a valuable and beneficial
opportunity for me. I have enjoyed having a platform from which to express my ideas
and concerns, I have-at least sometimes-enjoyed receiving the huge shipments of mail
forwarded to me from New York several times a month, and I have certainly been
happy to make new friends through this channel. I won't miss the monthly deadline,
but I will undoubtedly come across ideas, from time to time, that would have made
perfect "Metamagical Themas". I will be keeping them in mind, and maybe at some
future time will write a similar set of essays.
        But for now, it is time for me to move on to other territory: I look forward to a
return to my professional work, and to a more private life. Good-bye, and best wishes
to you and to all other readers of this magazine, this issue, this copy, this piece, this
page, this column, this paragraph, this sentence, and, last but not least, this "this".

Post Scriptum.
        What do you do when in a crushingly cold winter, you hear over the radio that
there is a severe natural gas shortage in your part of the country, and everyone is
requested to turn their thermostat down to 60 degrees? There's no way anyone will
know if you've complied or not. Why shouldn't you toast in your house and let all the
rest of the people cut down their consumption? After all, what you do surely can't
affect what anyone else does.
        This is a typical "tragedy of the commons" situation. A common resource has
reached the point of saturation or exhaustion, and the questions for each individual
now are: "How shall I behave? Am I typical? How does a




Irrationality Is the Square Root of All Evil                                         761

lone person's action affect the big picture?" Garrett Hardin's article "The Tragedy of
the Commons" frames the scene in terms of grazing land shared by a number of
herders. Each one is tempted to increase their own number of animals even when the
land is being used beyond its optimum capacity, because the individual gain
outweighs the individual loss, even though in the long run, that decision, multiplied
throughout the population of herders, will destroy the land totally.
        The real reason behind Hardin's article was to talk about the population
explosion and to stress the need for rational global planning-in fact, for coercive
techniques similar to parking tickets and jail sentences. His idea is that families
should be allowed to have many children (and thus to use a large share of the common
resources) but that they should be penalized by society in the same way as society
"allows" someone to rob a bank and then applies sanctions to those who have made
that choice. In an era when resources are running out in a way humanity has never had
to face heretofore, new kinds of social arrangements and expectations must be
imposed, Hardin feels, by society as a whole. He is a dire pessimist about any kind of
superrational cooperation, emphasizing that cooperators in the birth-control game will
breed themselves right out of the population. A perfect illustration of why this is so is
the man I heard about recently: he secretly had ten wives and by them had sired
something like 35 children by the time he was 30. With genes of that sort proliferating
wildly, there is little hope for the more modest breeders among us to gain the upper
hand. Hardin puts it bluntly: "Conscience is self-eliminating." He goes even further
and says:

   The argument has here been stated in the context of the population problem, but
   it applies equally well to any instance in which society appeals to an individual
   exploiting a commons to restrain himself for the general good-by means of his
   conscience. To make such an appeal is to set up a selective system that works
   toward the elimination of conscience from the race.

        An even more pessimistic vision of the future is proffered us by -one Walter
Bradford Ellis, a hypothetical speaker representing the views of his inventor, Louis
Pascal, in a hypothetical speech:

   The United States-indeed the whole earth-is fast running out of the resources it
   depends on for its existence. Well before the last of the world's supplies of oil
   and natural gas are exhausted early in the next century, shortages of these and
   other substances will have brought about the collapse of our whole economy
   and, indeed, of our whole technology. And without the wonders of modern
   technology, America will be left a grossly overpopulated, utterly impoverished,
   helpless, dying land. Thus I foresee a whole world full of wretched, starving
   people with no hope of escape, for the only countries which could have aided
   them will soon be no better off than the rest. And thus unless we are saved from
   this future by the blessing of a nuclear war or a truly lethal




Irrationality Is the Square Root of All Evil                                         762

   pestilence, I see stretching off into eternity a world of indescribable suffering
   and hopelessness. It is a vision of truly unspeakable horror mitigated only by the
   fact that try as I might I could not possibly concoct a creature more deserving of
   such a fate.

Whew! The circularity of the final thought reminds me of an idea I once had: that it
will be just as well if humanity destroys itself in a nuclear holocaust, because
civilizations that destroy themselves are barbaric and stupid, and who would want to
have one of them around, polluting the universe?
        Pascal's thoughts, expressed in his article "Human Tragedy and Natural
Selection" and in his rejoinder to an article by two critics called "The Loving Parent
Meets the Selfish Gene" (which is where Ellis' speech is printed), are strikingly
reminiscent of the thoughts of his earlier namesake Blaise, who in an unexpected use
of his own calculus of probabilities managed to convince himself that the best
possible way to spend his life was in devotion to a God who he wasn't sure (and
couldn't be sure) existed. In fact, Pascal felt, even if the chances of God's existence
were one in ,a million, faith in that God would pay off in the end, because the
potential rewards (or punishments) if Heaven and Hell exist are infinite, and all
earthly rewards and punishments, no matter how great, are still finite. The favored
behavior is to be a believer, Pascal "calculated"-regardless of what you do believe.
Thus Blaise Pascal devoted his brilliant mind to theology.
        Louis Pascal, following in his forebear's mindsteps, has opted to devote his life
to the world's population problem. And he can produce mathematical arguments to
show why you should, too. To my mind, there is no question that such arguments
have considerable force. There are always points to nitpick over, but in essence,
thinkers like Hardin and Pascal and Anne and Paul Ehrlich and many others have
recognized and internalized the novelty of the human situation at this moment in
history: the moment when humanity has to grapple with dwindling resources and
overwhelmingly huge weapons systems. Not many people are willing to wrestle with
this beast, and consequently the burden falls all the more heavily on those few who
are.

                                          *    *   *

        It has disturbed me how vehemently and staunchly my clear-headed friends
have been able to defend their decisions to defect. They seem to be able to digest my
argument about superrationality, to mull it over, to begrudge some curious kind of
validity to it, but ultimately to feel on a gut level that it is wrong, and to reject it. This
has led me to consider the notion that my faith in the superrational argument might be
similar to a self-fulfilling prophecy or self-supporting claim, something like being
absolutely convinced beyond a shadow of a doubt that the Henkin sentence "This
sentence is true" actually must be true-when, of course, it is equally defensible to
believe it to be false. The sentence is undecidable; its truth




Irrationality Is the Square Root of All Evil                                             763

value is stable, whichever way you wish it to go (in this way, it is the diametric
opposite of the Epimenides sentence "This sentence is false", whose truth value flips
faster than the tip of a happy pup's tail). One difference, though, between the
Prisoner's Dilemma and oddball self-referential sentences is that whereas your beliefs
about such sentences' truth values usually have inconsequential consequences, with
the Prisoner's Dilemma, it's quite another matter.
        I sometimes wonder whether there haven't been many civilizations Out There,
in our galaxy and beyond, that have already dealt with just these types of gigantic
social problems-Prisoner's Dilemmas, Tragedies of the Commons, and so forth. Most
likely some would have survived, some would have perished. And it occurs to me that
perhaps the ultimate difference in those societies may have been the survival of the
meme that, in effect, asserts the logical, rational validity of cooperation in a one-shot
Prisoner's Dilemma. In a way, this would be the opposite thesis to Hardin's. It would
say that lack of conscience is self-eliminating-provided you wait long enough that
natural selection can act at the level of entire societies.
        Perhaps on some planets, Type I societies have evolved, while on others, Type
II societies have evolved. By definition, members of Type I societies believe in the
rationality of lone, uncoerced, one-shot cooperation (when faced with members of
Type I societies), whereas members of Type II societies reject the rationality of lone,
uncoerced, one-shot cooperation, irrespective of who they are facing. (Notice the
tricky circularity of the definition of Type I societies. Yet it is not a vacuous
definition!) Both types of society find their respective answer to be obvious-they just
happen to find opposite answers. Who knows-we might even happen to have some
Type I societies here on earth. I cannot help but wonder how things would turn out if
my little one-shot Prisoner's Dilemma experiment were carried out in Japan instead of
the U.S. In any case, the vital question is: Which type of society survives, in the long
run?
        It could be that the one-shot Prisoner's Dilemma situations that I have
described are undecidable propositions within the logic that we humans have
developed so far, and that new axioms can be added, like the parallel postulate in
geometry, or Godel sentences (and related ones) in mathematical logic. (Take a look
at Figure 31-1, and see what kind of logic will extract those two poor devils from their
one-shot dilemma.) Those civilizations to which cooperation appears axiomatic-Type
I societieswind up surviving, I would venture to guess, whereas those to which
defection appears axiomatic-Type II societies-wind up perishing. This suggestion may
seem all wet to you, but watch those superpowers building those bombs, more and
more of them every day, helplessly trapped in a rising spiral, and think about it.
Evolution is a merciless pruner of ill logic.
        Most philosophers and logicians are convinced that truths of logic are
"analytic" and a priori; they do not like to think that such basic ideas are grounded in
mundane, arbitrary things like survival. They might admit that




Irrationality Is the Square Root of All Evil                                         764

                "The problem is how to turn loose without letting go."


FIGURE 31-1. One powerful metaphor for the absurdity we have collectively dug
ourselves into. The symmetry of the situation is acutely portrayed in this cartoon
drawn by Bill Mauldin in 1960. Note that if either person releases his rope, thus
chopping of his counterpart's head, that person's hand will go limp, thus releasing his
rope and causing the other blade to fall and chop of the head of the instigator. That
idea is a centerpiece of our current nuclear deterrence strategy: Even if we are wiped
of the globe, our trUSty missiles will still wreak divine revenge on the evil empire of
Satanic Uglies who dared do harm to US.

natural selection tends to favor good logic-but they would certainly hate the
suggestion that natural selection defines good logic! Yet truth and survival value are
all tangled together, and civilizations that survive certainly have glimpsed higher
truths than those that perish. When you argue with someone whose ideas you are sure
are wrong but who dances an infuriatingly inconsistent yet self-consistent verbal
dance in front of you, your one solace is that something in life may yet change this
person's mind, even though your own best logic is helpless to do so. Ultimately,
beliefs have to be grounded in experience, whether that experience is the organism's
or its ancestors' or its peer group's. (That s what Chapter 5, particularly its




Irrationality Is the Square Root of All Evil                                       765

P.S., was all about.) My feeling is that the concept of superrationality is one whose
truth will come to dominate among intelligent beings in the universe simply because
its adherents will survive certain kinds of situations where its opponents will perish.
Let's wait a few spins of the galaxy and see. After all, healthy logic is whatever
remains after evolution's merciless pruning.

                                        *      *   *

       I was describing the Copycat project (Chapter 24) to physicist Victor
Weisskopf, and I gave him our canonical example: "If abc goes to abd, what does xyz
go to?" After we had discussed various possible answers and settled on wyz as the
most compelling for reasons of symmetry, he surprised me by saying this: "You
know, the root of the world's deepest problems is the tragic inability on the part of the
world's leaders to see such basic symmetries. For instance, that the U.S. is to the S.U.
what the S.U. is to the U.S.-that is too much for them to accept." Oh, but how could
Weisskopf be so silly? After all, we're not trying to export communism to the entire
world!
       Logician Raymond Smullyan, who first heard about the Prisoner's Dilemma
from me and who was absolutely delighted by it, also surprised me, but in a different
way: He vehemently insisted on the correctness of defection in a one-shot situation no
matter who might be on the other side, including his twin or his clone! (He did waver
about his mirror image.) But just as I was giving up on him as a lost cause, he
conceded this much to me: "I suspect, Doug, that this problem is a lot knottier than
you or I suspect." Indeed, I suspect so, Raymond.




Irrationality Is the Square Root of All Evil                                         766


                          The Tale of Happiton
                                      June, 1983


HAPPITON was a happy little town. It had 20,000 inhabitants, give or take 7,
and they were productive citizens who mowed their lawns quite regularly. Folks in
Happiton were pretty healthy. They had a life expectancy of 75 years or so, and lots of
them lived to ripe old ages. Down at the town square, there was a nice big courthouse
with all sorts of relics from WW II and monuments to various heroes and whatnot.
People were proud, and had the right to be proud, of Happiton.
        On the top of the courthouse, there was a big bell that boomed every hour on
the hour, and you could hear it far and wide-even as far out as Shady Oaks Drive, way
out nearly in the countryside.
        One day at noon, a few people standing near the courthouse noticed that right
after the noon bell rang, there was a funny little sound coming from up in the belfry.
And for the next few days, folks noticed that this scratching sound was occurring after
every hour. So on Wednesday, Curt Dempster climbed up into the belfry and took a
look. To his surprise, he found a crazy kind of contraption rigged up to the bell. There
was this mechanical hand, sort of a robot arm, and next to it were five weird-looking
dice that it could throw into a little pan. They all had twenty sides on them, but
instead of being numbered 1 through 20, they were just numbered 0 through 9, but
with each digit appearing on two opposite sides. There was also a TV camera that
pointed at the pan and it seemed to be attached to a microcomputer or something.
That's all Curt could figure out. But then he noticed that on top of the computer, there
was a neat little envelope marked "To the friendly folks of Happiton". Curt decided
that he'd take it downstairs and open it in the presence of his friend the mayor, Janice
Fleener. He found Janice easily enough, told her about what he'd found, and then they
opened the envelope. How neatly it was written! It said this:




The Tale of Happiton                                                                767

                                                                        Grotto 19, Hades
                                                                          June 20, 1983

Dear folks of Happiton,

        I've got some bad news and some good news for you. The bad first. You know
your bell that rings every hour on the hour? Well, I've set it up so that each time it
rings, there is exactly one chance in a hundred thousand-that is, 1/100,000-that a Very
Bad Thing will occur. The way I determine if that Bad Thing will occur is, I have this
robot arm fling its five dice and see if they all land with `7' on top. Most of the time,
they won't. But if they do-and the odds are exactly 1 in 100,000-then great clouds of
an unimaginably revoltingsmelling yellow-green gas called "Retchgoo" will come
oozing up from a dense network of underground pipes that I've recently installed
underneath Happiton, and everyone will die an awful, writhing, agonizing death.
Well, that's the bad news.
        Now the good news! You all can prevent the Bad Thing from happening, if
you send me a bunch of postcards. You see, I happen to like postcards a whole lot
(especially postcards of Happiton), but to tell the truth, it doesn't really much matter
what they're of. I just love postcards! Thing is, they have to be written personally-not
typed, and especially not computer-printed or anything phony like that. The more
cards, the better. So how about sending me some postcards-batches, bunches, boxes of
them?
        Here's the deal. I reckon a typical postcard takes you about 4 minutes to write.
Now suppose just one person in all of Happiton spends 4 minutes one day writing me,
so the next day, I get one postcard. Well, then, I'll do you all a favor: I'll slow the
courthouse clock down a bit, for a day. (I realize this is an inconvenience, since a lot
of you tell time by the clock, but believe me, it's a lot more inconvenient to die an
agonizing, writhing death from the evil-smelling, yellow-green Retchgoo.) As I was
saying, I'll slow the clock down for one day, and by how much? By a factor of 1.0000
1. Okay, I know that doesn't sound too exciting, but just think if all 20,000 of you
send me a card! For each card I get that day, I'll toss in a slow-up factor of 1.00001,
the next day. That means that by sending me 20,000 postcards a day, you all, working
together, can get the clock to slow down by a factor of 1.00001 to the 20, 000th
power, which is just a shade over 1.2, meaning it will ring every 72 minutes.
        All right, I hear you saying, "72 minutes is just barely over an hour!" So I offer
you more! Say that one day I get 160,000 postcards (heavenly!). Well then, the very
next day I'll show my gratitude by slowing your clock down, all day long, midnight to
midnight, by 1.00001 to the 160,000th power, and that ain't chickenfeed. In fact, it's
about 5, and that means the clock will ring only every 5 hours, meaning those sinister
dice will only get rolled about 5 times (instead of the usual 24). Obviously, it's better
for both of us that way. You have to bear in mind that I don't have any personal
interest in seeing that awful Retchgoo come rushing and gushing up out of those pipes
and causing every last one of you to perish in grotesque, mouth-foaming, twitching
convulsions. All I care about is getting postcards! And to send me 160,000 a day
wouldn't cost you folks that much effort, being that it's just 8 postcards a day just
about a half hour a day for each of you, the way I reckon it.




The Tale of Happiton                                                                  768

       So my deal is pretty simple. On any given day, I'll make the clock go off once
every X hours, where X is given by this simple formula:

       X =1.00001N

Here, N is the number of postcards I received the previous day. If N is 20,000, then X
will be 1.2, so the bell would ring 20 times per day, instead of 24. If N is 160,000,
then X jumps way up to about 5, so the clock would slow way down just under 5 rings
per day. If I get no postcards, then the clock will ring once an hour, just as it does
now. The formula reflects that, since if N is 0, X will be 1. You can work out other
figures yourself. Just think how much safer and securer you'd all feel knowing that
your courthouse clock was ticking away so slowly!
       I'm looking forward with great enthusiasm to hearing from you all.

                                                                           Sincerely yours,
                                                                             Demon \#3127

       The letter was signed with beautiful medieval-looking flourishes, in an
unusual shade of deep red ... ink?
       "Bunch of hogwash!" spluttered Curt. "Let's go up there and chuck the whole
mess down onto the street and see how far it bounces." While he was saying this,
Janice noticed that there was a smaller note clipped onto the back of the last sheet,
and turned it over to read it. It said this:

P. S. -It's really not advisable to try to dismantle my little set-up up there in the belfry:
I've got a hair trigger linked to the gas pipes, and if anyone tries to dismantle it,
pssssst! Sorry.

        Janice Fleener and Curt Dempster could hardly believe their eyes. What gall!
They got straight on the -phone to the Police Department, and talked to Officer
Curran. He sounded poppin' mad when they told him what they'd found, and said he'd
do something about it right quick. So he hightailed it over to the courthouse and ran
up those stairs two at a time, and when he reached the top, a-huffin' and a-puffin', he
swung open the belfry door and took a look. To tell the truth, he was a bit ginger in
his inspection, because one thing Officer Curran had learned in his many years of
police experience is that an ounce of prevention is worth a pound of cure. So he
cautiously looked over the strange contraption, and then he turned around and quite
carefully shut the door behind him and went down. He called up the town sewer
department and asked them if they could check out whether there was anything funny
going on with the pipes underground.
        Well, the long and the short of it is that they verified everything in the
Demon's letter, and by the time they had done so, the clock had struck five more times
and those five dice had rolled five more times. Janice Fleener had in fact had her
thirteen-year-old daughter Samantha go up and sit in a




The Tale of Happiton                                                                    769

wicker chair right next to the microcomputer and watch the robot arm throw those
dice. According to Samantha, an occasional 7 had turned up now and then, but never
had two 7's shown up together, let alone 7's on all five of the weird-looking dice!

                                       *       *      *

        The next day, the Happiton Eagle-Telephone came out with a front-page story
telling all about the peculiar goings-on. This caused quite a commotion. People
everywhere were talking about it, from Lidden's Burger Stop to Bixbee's Druggery. It
was truly the talk of the town.
        When Doc Hazelthorn, the best pediatrician this side of the Cornyawl River,
walked into Ernie's Barbershop, corner of Cherry and Second, the atmosphere was
more somber than usual. "Whatcha gonna do, Doc?" said big Ernie, the jovial barber,
as he was clipping the few remaining hairs on old Doc's pate. Doc (who was also head
of the Happiton City Council) said the news had come as quite a shock to him and his
family. Red Dulkins, sitting in the next chair over from Doc, said he felt the same
way. And then the two gentlemen waiting to get their hair cut both added their words
of agreement. Ernie, summing it up, said the whole town seemed quite upset. As Ernie
removed the white smock from Doc's lap and shook the hairs off it, Doc said that he
had just decided to bring the matter up first thing at the next City Council meeting,
Tuesday evening. "Sounds like a good idea, Doc!" said Ernie. Then Doc told Ernie he
couldn't make the usual golf date this weekend, because some friends of his had
invited him to go fishing out at Lazy Lake, and Doc just couldn't resist.
        Two days after the Demon's note, the Eagle-Telephone ran a feature article in
which many residents of Happiton, some prominent, some not so prominent, voiced
their opinions. For instance, eleven-year-old Wally Thurston said he'd gone out and
bought up the whole supply of picture postcards at the 88-Cent Store, $14.22 worth of
postcards, and he'd already started writing a few. Andrea McKenzie, sophomore at
Happiton High, said she was really worried and had had nightmares about the gas, but
her parents told her not to worry, things had a way of working out. Andrea said maybe
her parents weren't taking it so seriously because they were a generation older and
didn't have as long to look forward to anyway. She said she was spending an hour
each day writing postcards. That came to 15 or 16 cards each day. Hank Hoople, a
janitor at Happiton High, sounded rather glum: "It's all fate. If the bullet has your
name on it, it's going to happen, whether you like it or not." Many other citizens
voiced concern and even alarm about the recent developments.
        But some voiced rather different feelings. Ned Furdy, who as far as anyone
could tell didn't do much other than hang around Simpson's bar all day (and most of
the night) and buttonhole anyone he could, said, "Yeah, it's a problem, all right, but I
don't know nothin' about gas and statistics and such.




The Tale of Happiton                                                                770

It should all be left to the mayor and the Town Council, to take care of. They know
what they're doin'. Meanwhile, eat, drink, and be merry!" And Lulu Smyth, 77-year-
old. proprietor of Lulu's Thread 'N Needles Shop, said "I think it's all a ruckus in a
teapot, in my opinion. Far as I'm concerned, I'm gonna keep on sellin' thread 'n
needles, and playin' gin rummy every third Wednesday."

                                        *   *   *

        When Doc Hazelthorn came back from his fishing weekend at Lazy Lake, he
had some surprising news to report. "Seems there's a demon left a similar set-up in the
church steeple down in Dwaynesville", he said. (Dwaynesville was the next town
down the road, and the arch-rival of Happiton High in football.) "The Dwaynesville
demon isn't threatening them with gas, but with radioactive water. Takes a little
longer to die, but it's just as bad. And I hear tell there's a demon with a subterranean
volcano up at New Athens." (New Athens was the larger town twenty miles up the
Cornyawl from Dwaynesville, and the regional center of commerce.)
        A lot of people were clearly quite alarmed by all this, and there was plenty of
arguing on the streets about how it had all happened without anyone knowing. One
thing that was pretty universally agreed on was that a commission should be set up as
soon as possible, charged from here on out with keeping close tabs on all subterranean
activity within the city limits, so that this sort of outrage could never happen again. It
appeared probable that Curt Dempster, who was the moving force behind this idea,
would be appointed its first head.
        Ed Thurston (Wally's father) proposed to the Jaycees (of which he was a
member in good standing) that they donate $1,000 to support a postcard-writing
campaign by town kids. But Enoch Swale, owner of Swale's Pharmacy and the
Sleepgood Motel, protested. He had never liked Ed much, and said Ed was proposing
it simply because his son would gain status that way. (It was true that Wally had
recruited a few kids and that they spent an hour each afternoon after school writing
cards. There had been a small article in the paper about it once.) After considerable
debate, Ed's motion was narrowly defeated. Enoch had a lot of friends on the City
Council.
        Nellie Doobar, the math teacher at High, was about the only one who checked
out the Demon's math. "Seems right to me", she said to the reporter who called her
about it. But this set her to thinking about a few things. In an hour or two, she called
back the paper and said, "I figured something out. Right now, the clock is still ringing
very close to once every hour. Now there are about 720 hours per month, and so that
means there are 720 chances each month for the gas to get out. Since each chance is 1
in 100,000, it turns out that each month, there's a bit less than a 1-in-100 chance that
Happiton will get gassed. At that rate, there's about 11 chances in 12 that Happiton
will make it through each year. That may sound pretty




The Tale of Happiton                                                                  771

good, but the chances we'll make it through any 8-year period are almost exactly 50-
50, exactly the same as tossing a coin. So we can't really count on very many years ..."
         This made big headlines in the next afternoon's Eagle-Telephone-in fact, even
bigger than the plans for the County Fair! Some folks started calling up Mrs. Doobar
anonymously and telling her she'd better watch out what she was saying if she didn't
want to wind up with a puffy face or a fat lip. Seems like they couldn't quite keep it
straight that Mrs. Doobar wasn't the one who'd set the thing up in the first place.
         After a few days, though, the nasty calls died down pretty much. Then Mrs.
Doobar called up the paper again and told the reporter, "I've been calculating a bit
more here, and I've come up with the following, and they're facts every last one of
them. If all 20,000 of us were to spend half an hour a day writing postcards to the
Demon, that would amount to 160,000 postcards a day, and just as the Demon said,
the bell would ring pretty near every five hours instead of every hour, and that would
mean that the chances of us getting wiped out each month would go down
considerable. In fact, there would only be about 1 chance in 700 that we'd go down
the tubes in any given month, and only about a chance in 60 that we'd get zapped each
year. Now I'd say that's a darn sight better than 1 chance in 12 per year, which is what
it is if we don't write any postcards (as is more or less the case now, except for Wally
Thurston and Andrea McKenzie and a few other kids I heard of). And for every 8-
year period, we'd only be running a 13 percent risk instead of a 50 percent risk."
         "That sounds pretty good", said the reporter cheerfully.
         "Well," replied Mrs. Doobar, "it's not too bad, but we can get a whole lot
better by doublin' the number of postcards."
         "How's that, Mrs. Doobar?" asked the reporter. "Wouldn't it just get twice as
good?"
         "No, you see, it's an exponential curve," said Mrs. Doobar, "which means that
if you double N, you square X. "
         "That's Greek to me", quipped the reporter.
         "N is the number of postcards and X is the time between rings", she replied
quite patiently. "If we all write a half hour a day, X is 5 hours. But that means that if
we all write a whole hour a day, like Andrea McKenzie in my algebra class, X jumps
up to 25 hours, meaning that the clock would ring only about once a day, and
obviously, that would reduce the danger a lot. Chances are, hundreds of years would
pass before five 7's would turn up together on those infernal dice. Seems to me that
under those circumstances, we could pretty much live our lives without worrying
about the gas at all. And that's for writing about an hour a day, each one of us."
         The reporter wanted some more figures detailing how much different amounts
of postcard-writing by the populace would pay off, so Mrs. Doobar obliged by going
back and doing some more figuring. She figured out that if 10,000 people-half the
population of Happiton-did 2 hours a day for




The Tale of Happiton                                                                 772

the year, they could get the same result-one ring ever)' 25 hours. If only 5,000 people
spent 2 hours a day, or if 10,000 people spent one hour a day, then it would go back to
one ring every 5 hours (still a lot safer than one every hour). Or, still another way of
looking at it, if just 1250 of them worked full-time (8 hours a day), they could achieve
the same thing.
        "What about if we all pitch in and do 4 minutes a day, Mrs. Doobar?" asked
the reporter.
        "Fact is, 'twouldn't be worth a damn thing! (Pardon my French.)" she replied.
"N is 20,000 that way, and even though that sounds pretty big, X works out to be just
1.2, meaning one ring every 1.2 hours, or 72 minutes. That way, we still have about a
chance of 1 in 166 every month of getting wiped out, and 1 in 14 every year of getting
it. Now that's real scary, in my book. Writing cards only starts making a noticeable
difference at about 15 minutes a day per person."

                                       *   *   *

        By this time, several weeks had passed, and summer was getting into full
swing. The County Fair was buzzing with activity, and each evening after folks came
home, they could see loads of fireflies flickering around the trees in their yards.
Evenings were peaceful and relaxed. Doc'Hazelthorn was playing golf every
weekend, and his scores were getting down into the low 90's. He was feeling pretty
good. Once in a while he remembered the Demon, especially when he walked
downtown and passed the courthouse tower, and every so often he would shudder.
But he wasn't sure what he and the City Council could do about it.
        The Demon and the gas still made for interesting talk, but were`no longer such
big news. Mrs. Doobar's latest revelations made the paper, but. were. relegated this
time to the second section, two pages before the comics, right next to the daily
horoscope column. Andrea McKenzie read the article avidly, and showed it to a lot of
her school friends, but to her surprise, it didn't seem to stir up much interest in them.
At first, her best friend, Kathi Hamilton, a very bright girl who had plans to go to
State and major in history, enthusiastically joined Andrea and wrote quite a few cards
each day. But after a few days, Kathi's enthusiasm began to wane.
        "What's the point, Andrea?" Kathi asked. "A handful of postcards from me
isn't going to make. the slightest bit of difference. Didn't you read Mrs. Doobar's
article? There have got to be 160,000 a day to make a big difference."
        "That's just the point, Kath!" replied Andrea exasperatedly. "If you and
everyone else will just do your part, we'll reach that number-but you can't cop out!"
Kathi didn't see the logic, and spent most of her time doing her homework for the
summer school course in World History she was taking. After all, how could she get
into State if she flunked World History?
        Andrea just couldn't figure out how come Kathi, of all people, so




The Tale of Happiton                                                                 773

interested in history and the flow of time and world-events, could not see her own life
being touched by such factors, so she asked Kathi, "How do you know there will be
any you-left to go to State, if you don't write postcards? Each year, there's a 1-in-12
chance of you and me and all of us being wiped out! Don't you even want to work
against that? If people would just care, they could change things! An hour a day! Half
an hour a day! Fifteen minutes a day!"
         "Oh, come on, Andrea!" said Kathi annoyedly, "Be realistic."
         "Darn it all, I'm the one who's being realistic", said Andrea. "If you don't help
out, you're adding to the burden of someone else."
         "For Pete's sake, Andrea", Kathi protested angrily, "I'm not adding to anyone
else's burden. Everyone can help out as much as they want, and no one's obliged to do
anything at all. Sure, I'd like it if everyone were helping, but you can see for yourself,
practically nobody is. So I'm not going to waste my time. I need to pass World
History."
         And sure enough, Andrea had to do no more than listen each hour, right on the
hour, to hear that bell ring to realize that nobody was doing much. It once had
sounded so pleasant and reassuring, and now it sounded creepy and ominous to her,
just like the fireflies and the barbecues. Those fireflies and barbecues really bugged
Andrea, because they seemed so normal, so much like any other summer-only this
summer was not like any other summer. Yet nobody seemed to realize that. Or, '
rather, there was an undercurrent that things were not quite as they should be, but
nothing was being done ...
         . One Saturday, Mr. Hobbs, the electrician, came around to fix a broken
refrigerator at the McKenzies' house. Andrea talked to him about writing postcards to
the Demon. Mr. Hobbs said to her, "No time, no time! Too busy fixin' air
conditioners! In this heat wave, they been breakin' down all over town. I-work a 10-
hour day as it is, and now it's up to 11, 12 hours a day, includin' weekends. I got no
time for postcards, Andrea." And Andrea .saw that for Mr. Hobbs, it was true. He had
a big family and his children went to parochial school, and he had to pay for them all,
and ...
         Andrea's older sister's boyfriend, Wayne, was a star halfback at Happiton
High. One evening he was over and teased Andrea about her postcards. She asked
him, "Why don't you write any, Wayne?".
         "I'm out lifeguardin' every day, and the rest of the time I got scrimmages - for
the fall season."
         "But you could take some time out just 15 minutes a day-and write a few
postcards!" she argued. He just laughed and looked a little fidgety. "I don't know,
Andrea",-he said. "Anyway, me 'n Ellen have got better things to do-huh, Ellen?"
Ellen giggled and blushed a little. Then they ran out of the house and jumped into
Wayne's sports car to go bowling at the Happi-Bowl.

                                        *   *   *




The Tale of Happiton                                                                  774

        Andrea was puzzled by all her friends' attitudes. She couldn't understand why
everyone had started out so concerned but then their concern had fizzled,, as if the
problem had gone away. One day when she was walking home from school, she saw
old Granny Sparks out watering her garden. Granny, as everyone called her, lived
kitty-corner from the McKenzies and was always chatty, so Andrea stopped and asked
Granny Sparks what she thought of all this. "Pshaw! Fiddlesticks!" said Granny
indignantly. "Now Andrea, don't you go around believin' all that malarkey they print
in the newspapers! Things are the same here as they always been. I oughta know -I've
been livin' here nigh on 85 years!
        Indeed, that was what bothered Andrea. Everything seemed so annoyingly
normal. The teenagers with their cruising cars and loud motorcycles. The usual boring
horror movies at the Key Theater down on the square across from the courthouse. The
band in the park. The parades. And especially, the damn fireflies! Practically nobody
seemed moved or affected by what to her seemed the most overwhelming news she'd
ever heard. The only other truly sane person she could think of was little Wally
Thurston, that elevenyear-old from across town. What a ridiculous irony, that an
eleven-year-old was saner than all the adults!
        Long about August 1, there was an editorial in the paper that gave Andrea a
real lift. It came from out of the blue. It was written by the paper's chief editor,
"Buttons" Brown. He was an old-time journalist from St. Jo, Missouri. His editorial
was real short. It went like this:

The Disobedi-Ant

          The story of the Disobedi-Ant is very short. It refused to believe that its
  powerful impulses to play instead of work were anything but unique
  expressions of its very unique self, and it went its merry way, singing, "What I
  choose to do has nothing to do with what any-ant else chooses to do! What
  could be more self-evident?"
          Coincidentally enough, so went the reasoning of all its colony-mates. In
  fact, , the same refrain was independently invented by every last ant in the
  colony, and each ant thought it original. It echoed throughout the colony, even
  with the same melody.
          The colony perished.

Andrea thought this was a terrific allegory, and showed it to all her friends. They
mostly liked it, but to her surprise, not one of them started writing postcards.
       All in all, folks were pretty much back to daily life. After all, nothing much -
seemed really to have changed. The weather had turned real hot, and folks
congregated around the various swimming pools in town. There were lots of
barbecues in the evenings, and, every once in a while somebody'd make a joke or two
about the Demon and the postcards. Folks would chuckle and




The Tale of Happiton                                                                775

then change the topic. Mostly, people spent their time doing what they'd always done,
and enjoying the blue skies. And mowing their lawns regularly, since they wanted the
town to look nice.

Post Scriptum

                                              The atomic bomb has changed everything
                                                 except our way of thinking. And so we
                                          drift helplessly towards unparalleled disaster.
                                                                       -Albert Einstein


        People of every era always feel that their era has the severest problems that
people have ever faced. At first this sounds silly. How can every era be the toughest?
But it's not silly. Things can be getting constantly more dangerous and frightful, and
that would mean that each new generation truly is facing unprecedentedly serious
problems. As for us, we have the problem of extinction on our hands.
        Someone once said that our current situation vis-a-vis the Soviet Union is like
two people standing knee-deep in a room filled with gasoline. Both hold open
matchbooks in their hands. One person is jeering at the other:.. "Ha ha ha! My
matchbook is full, and yours is only half full! Ha ha ha!"
        The reality of our situation is about that simple. The vast majority of people,
however, refuse to let this reality seep into their systems and change their day-to-day
behaviors. And thus the validity of Einstein's gloomy utterance.

                                      *    *   *

        I remember many years ago reading an estimate that the famous geneticist
George Wald had made about nuclear war. He said he figured there was a two percent
chance per year of a nuclear war taking place. This amounts to throwing one 50-sided
die (or a couple of seven-sided dice) once a year, and hoping that it doesn't come up
on the bad side. How Wald arrived at his figure of two percent per year, I don't know.
But it was vivid. The figure has stuck with me for a couple of decades. I tend to think
that the chances are greater nowadays than they were back then: maybe about five
percent per year. But who can say?
        The Bulletin of the Atomic Scientists features a clock on its cover. This clock
doesn't tick, it just hovers. It hovers near midnight, sometimes getting closer,
sometimes receding a bit. Right now, it's at three minutes to midnight. Back at the
signing of SALT I, it was at twelve minutes before midnight. The closest it ever came
was two minutes before midnight, and I think that was at the time of the Cuban
missile crisis.
        The purpose of the clock is to symbolize the current danger of a nuclear




The Tale of Happiton                                                                 776

holocaust. It's a little like those "Danger of Eire Today" signs that Smokey the Bear
holds up for you as you enter a national forest in the summer. It is a subjective
estimate, made by the magazine's board of directors. Now what is the meaning of
"danger", if not probability of disaster per unit time? Surely, the more dangerous a
place or situation, the faster you want to get out of it, for, just that reason. Therefore,
it seemed to me that the Bulletin's number of minutes before midnight, B, was really a
coded way of expressing a Wald number, W-a probability of nuclear war per year.
And so I decided to make a subjective table, matching up the values of B that I knew
about with my own best estimates of W. After a bit of experimentation, I came up
with the following table:

       Bulletin Clock                  Wald's percentage
       (minutes before midnight)       (probability per year)

               1 min .............     20 percent
               2 mins .............    10 percent
               3 rains .............   7 percent
               4 mins .............    5 percent
               5 mins .............    4 percent
               7 mins .............    3 percent
               10 mins             :   2 percent
               12 mins .............   1.5 percent
               20 mins .............   1 percent


A fairly accurate summary of this subjective correspondence is given by the following
simple equation:

       W = 20/B

This estimates for you the holocaust danger per orbit of the earth, as a function of the
current setting of the Bulletin's clock.
        W and B may not be estimable in any truly scientific way, but there is a
definite reality behind them, even if not. so simple as that of N and X in Happiton.
Obviously it is not a "random'." dicelike process that will determine whether nuclear
war erupts in any given year. Nonetheless, it makes good sense to think of it in terms
of a probability per year, since what actually does determine history is a lot of things
that are in effect random, from the point of view of any less-than-omniscient being.
What other people (or countries) do is unpredictable and uncontrollable: it might as
well be random.
        If tensions get unbearably high in the Middle East or in Central America, that
is not something that we could have predicted or forestalled. If some terrorist group
manufactures and uses or threatens to use -a nuclear bomb, that is essentially a
"random" event. If overpopulation in Asia or starvation




The Tale of Happiton                                                                   777

In Africa or crop failures in the Soviet Union or oil gluts or shortages create huge
tensions between nations, that is like a random variable, like a throw of dice. Who
could have predicted the crazy flareup between Britain and Argentina over the silly
Falkland Islands? Who knows where the next hot spot will turn out to be? The global
temperature can change as swiftly and capriciously as a bright summer day can turn
sultry and menacing-even in Happiton.

                                        *       *       *

        It is the vivid imagery behind the Wald number and the Bulletin clock that
first got me thinking in terms of the Happiton metaphor. The story was pretty easy to
write, once the metaphor had been concocted. I had to work out the mathematics as I
went along, but otherwise it flowed easily. It was crucial to me that the numbers in.
the allegory seem realistic. The most important numbers were: (1) the chance of
devastation per year, which came out about right, as I see it; and (2) the amount of
time per day that I think would begin to make a significant difference if devoted by a
typical person to some sort of activity geared toward the right ends. In Happiton, that
threshold turned out to be about fifteen minutes per day per person. Fifteen minutes a
day is just about the amount of time that I think would begin to make a real difference
in the real world, but there are two ways that one might draw a distinction between
the situation in Happiton and the actual case.
        Firstly, some people say that the situation in Happiton is much simpler than
that of global competition and potential nuclear war. In Happiton, it's obvious that
writing postcards will do some good, whereas it's not so obvious (they claim) what
kind of action will do any good in the real world. Working hard for a freeze or for a
reduction of US-SU tensions might even be harmful, they claim! The situation is so
complex that nothing corresponds to the simplistic and sure-fire recipe of writing
postcards.
        Ah, but there is a big fallacy here. Writing postcards in Happiton is not sure-
fire. The gas could still come oozing up at any time. All that changes is the odds. Now
in the real world, we must follow our own best estimates, in the absence of perfect
information, as to what actions are likely to be positive andd what ones to be negative.
You can only follow your nose. You can never be sure that any action, no matter how
well intended, is going to improve. the situation. That's just the way life is.
        I happen to believe that the odds of a holocaust will be reduced (perhaps by a
factor of 1.0000001-) by writing to my representatives and senators fairly regularly,
by attending local freeze meetings, by contributing to various organizations, by giving
lectures here and there on the topic, and by writing articles like this. How can I know
that it will do any good? I can't, of course. And it's no different in Happiton. The best
of intentions can backfire for totally unforeseeable reasons. It might turn out that little
Wally Thurston, by moving his pencil in a certain graceful curlicue motion one




The Tale of Happiton                                                                   778

afternoon while writing his 1,000th postcard to the Demon, stirs up certain air
molecules which, by bouncing and jouncing against other ones helter-skelter, wind up
giving that tiny last push to the caroming icosahedral dice atop the belfry, and bang!
They all come up '7'! Wally, oh Wally, why such folly? Why did you ever write those
postcards?
         Those who would caution people that it might be counter-productive to work
against the arms race-unless they believe one should work for the arms race-are in
effect counseling paralysis. But would they do so in other areas of life? You never
know if that car trip to the grocery store won't be the last thing you do in your life. All
life is a gamble.
         The second distinction between Happiton and reality is this. In Happiton, for
fifteen minutes a day to make a noticeable dent, it would have had to be donated by
all 20,000 citizens, adults and children. Obviously I do not think that is realistic in our
country. The fifteen minutes a day per person that I would like to see spent by real
people in this country is limited to adults (or at least people of high-school age), and I
don't even include most adults in this. I cannot realistically hope that everyone will be
motivated to become politically active. Perhaps a highly active minority of five
percent would be enough. It is amazing how visible and influential an articulate and
vocal minority of,that size can be! So, being realistic, I limit 'my desires to an average
of fifteen minutes of activity per day for five percent of the adult American
population. I sincerely believe that with about this much work, a kind of turning point
would be reached-and that at 30 minutes or 60 minutes per day (exactly as in
Happiton), truly significant changes in the national mood (and hence in the global
danger level) could be effected.

                                        *   *   *

        I think I have explained what Happiton was written for. Trigger activity it may
not. I'm growing a little more realistic, and I don't expect much of anything. But I
would like to understand human nature. better, to understand what it is that makes us
so much like stupid gnats dully buzzing above a freeway, unable to see the onrushing
truck, 100 yards down the road, against whose windshield we are about to be
smashed.
        One last thought: Although to me it seems that nuclear war is the gravest
threat before us, I would grant that to other people it might appear otherwise. I don't
care so much what kinds of efforts people invest their time in, as long as they do
something. The exact thing that corresponds to the threat to Happiton doesn't much
matter. It could be nuclear weapons, chemical or biological weapons, the population
explosion, the U.S.'s ever-deepening involvement in Central America, or even
something more contained, like the environmental devastation inside the U.S. What it
seems to me is needed is a healthy dose of indignation: a spark, a flame, a fire inside.
Until that happens, that courthouse clock'll be tickin' away, once every hour, on the
hour, until ...




The Tale of Happiton                                                                   779

Post Post Scriptum.
        Two magazines are devoted to the prevention of nuclear war. They are: the
Bulletin of the Atomic Scientists and Nuclear Times. The Bulletin, founded in 1945,
aims to forestall nuclear holocaust by promoting awareness and understanding of the
issues involved. It describes itself as "a magazine of science and world affairs". Its
address is: 5801 South Kenwood Avenue, Chicago, Illinois 60637.
        Nuclear Times is a more recent arrival, and calls itself "the news magazine of
the antinuclear weapons movement". Its articles are shorter and lighter than those of
the Bulletin, but it keeps you up to date on what's happening all over the country and
the world. Its address is: Room 512, 298 Fifth Avenue, New York, New York 10001.
        The following organizations are effective and important forces in the attempt
to slow down the arms race and to reduce global tensions. Most of them put out
excellent literature, which is available in, large quantities at low prices (sometimes
free) for distribution. Needless to say, they can always use more members and more
funding. Many have local chapters.

The Council for a Livable World
11 Beacon Street Boston,
Massachusetts 02108

SANE
711 G Street, S.E. Washington,
D.C. 20003

Nuclear Weapons Freeze Campaign
4144 Lindell Boulevard, Suite 404
St. Louis, Missouri 63108

The Center for Defense Information
303 Capital Gallery West
600 Maryland Avenue, S.W.
Washington, D.C. 20024

Physicians for Social Responsibility,
639 Massachusetts Avenue
Cambridge, Massachusetts 02139

International Physicians for the Prevention of Nuclear War
225 Longwood Avenue, Room 200
Boston, Massachusetts 02115

Union of Concerned Scientists
1384 Massachusetts Avenue
Cambridge, Massachusetts 02238




The Tale of Happiton                                                              780


                    The Tumult of Inner Voices,
                     or, What Is the Meaning
                         of the Word "I"?
                   Grace Adams Tanner Lecture in Human Values
                           Southern Utah State College
                                Cedar City, Utah
                                   May, 1982


I pushed a button, and in an instant, Utah ceased to exist. Totally obliterated, beyond
recall. There was nothing I could do now, no matter how much I wished I hadn't
pushed that button, no matter how recent the action was. A mere second after it had
been pushed, there was no way to undo my action. A miscalculation with dire
consequences.
         Utah, a sandy state, full of deserts and strange, barren scenery. A beautiful
state, a place I had passed through many times, always with a sense of wonder. Eerie,
esonant names like "Uintah", "Wasatch", "Moab", "Koosharem", "Shivwits",
"Tavaputs", "Panguitch" .. .
         Now, all those names had been destroyed, leaving not a trace. All those names
would have to be retyped. But that was not the bad part. The bad part was that all my
ideas, the inspired ones I'd had a few days ago, had gone down the drain as well. Of
course, the first thing I'd checked was if there were any backup copies. There had
been two. But I'd destroyed both of them as well. Just moments ago there had been
three, yes, three, copies of Utah on my directory, and now they were all gone. The
disk space had been released, and perhaps for a few seconds my file's bit-patterns had
continued to spin around, no longer protected, no longer shielded, yet still intact. But
then, inevitably, somebody had wanted to write a file, and mercilessly, the operating
system handed my space over to them. Now somebody else's. bit-patterns had
overwritten mine.
         I was desperate. I hoped against hope that there was still some way to get




The Tumult of Inner Voices                                                          781

Utah back. They checked for me if there were any versions on tape, and sure enough,
it appeared there was one-a day old. Whew! At least something remained of my work.
But within minutes we discovered that it was only four lines long. It had been taped
before I'd had those good ideas. So every last shred of hope was destroyed.
        I knew I would be unable to reconstruct my ideas. I would just have to start
again from scratch. It was a horrible feeling, to have deleted my file called "Utah".
And just one single bad keystroke had done it. And yet, here I am again, typing away,
creating a new file by the same name, trying to construct new ideas to replace the old
ones. Or rather, here I am, standing before you, reading these week-old words to you,
words that came off that same spinning disk where my others had spun before them.
        Which is right? Am I really here at my terminal, typing, or am I really here,
standing before you? I can't decide. ,
        And that is one aspect of the question that I wish to confront today, in this
Grace Adams Tanner Lecture on Human Values. The \_question is an intellectual,
philosophical one: What is the meaning of the word "I"? Yet the question is also a
pragmatic, real-life, soul-ripping issue: Which one of the many people who I am, the
many inner voices inside me, will dominate? Who, or how, will I be? Which part of
me decides? And can that part in turn have inner conflicts about how to decide which
version of me it wants to let dominate?

                                      *   *   *

        Such were the questions I was confronting when I was typing that earlier
version of Utah that I spoke about at the beginning. But I have to confide that I made
a slight distortion. I spoke about those place names like "Uintah" and "Wasatch"
being destroyed when I pushed the button. Actually, those place names were not in
the original version of my talk. Only after the irreversible calamity of deleting that
original version of Utah the file had taken place did the metaphor of destroying Utah
the state by pushing a button lead me to-imagine the places in it that would be
obliterated. Thus, only in my new file called "Utah" did I actually type those place
names one by one. And that is this file, upon which I am now typing. Thus, only if I
make the same mistake once again-and I certainly hope I won't be that careless -will
those names be lost.
        However, it is no accident that the analogy between destroying Utah the state
and Utah the computer file arose in my mind. Far from an accident. In fact, what, was
that old file of inspired ideas about? It was about an inner fight between two voices
inside of me, two competing selves, two major facets of the person Doug Hofstadter,
vying against each other. 7t was, in fact, a hypothetical dialogue taking place before
you, a dialogue between two persons both of whom are inside me, both of whom are
genuinely myself, but who are at odds, in some sense, with each other.
        Those of you who are familiar with my book Gödel, Escher, Bach will




The Tumult of Inner Voices                                                        782

member the characters of Achilles and the Tortoise (as welt as a few others), whom I
used in my Dialogues. I wrote another dialogue more recently, with Achilles and the
Tortoise once again, in which they discuss the soul-searching question "Who shoves
whom around inside the careenium?" In that dialogue [Chapter 25 of this book], I
deal-or should say, they deal-with the question of what governs the soul. How can we
understand the nature of our selves when we are composed of so many myriads of
parts, none of which we understand? How are those parts put together? How does the
total add up to, a self, a soul, a you or a me?
        That is the metaphor of the "careenium"-a kind of enormous arena in which
billions of marbles careen around, bashing into each other unconsciously, and yet
which gives rise, when one stands back and looks at the whole, to a vast
consciousness. The title question, "Who shoves whom around inside the careenium?"
concerns how to look at such a system, one that has various levels of description. The
central issue is whether the marbles shove the total system around, or whether the
desires of the system as a whole shove the marbles around. It's a slippery issue, and
pinning down the nature of free will-referred to also as free won't in the dialogue!-is
the name of the game.
        Part of me was intending to read that dialogue to you tonight. One inner voice
spoke for it. But another, more urgent inner voice spoke eloquently against it. And the
transcript of the debate between those inner selves was what was deleted by my
careless finger. That was the original file called "Utah".
        Who was this other Doug Hofstadter that was so rudely intruding? And what
did he want to say? Why was he fighting for control of my top level? What in him
insisted that the story of the careenium was not appropriate? Well, I know of no -
better way to explain this than to let him speak.for himself. So without further ado,
here he is.

                                         *   *   *

        Thank you, but I feel a little awkward about this. After all, I don't really feel as
if I am a different person from you. That is, from the person who graciously
consented to let me have a few words. I am really the same person, am I not?
Sometimes I am not sure. After all, who is it that was invited to give the Grace Adams
Tanner Lecture on Human Values? Was it Douglas Hofstadter the author, or Douglas
Hofstadter the person? The former won a Pulitzer Prize, writes a monthly column, is
publicly visible, and hence seems a likely candidate for being invited to give a lecture.
And the latter is unknown, except to friends. Yet all the energy of the former comes
from an invisible, inner person, a hot fiery core of combatting ideas and hopes and
goals. And so in essence, all that was just done is to symbolically yield the floor to the
real person in whose mind all those ideas seethe and tangle with each other.




The Tumult of Inner Voices                                                             783

        And since I am now the real person, not the image, I can feel free to talk about
what it is that is gripping me these days, and splitting me into several people (of
which I am one). The thing that is gripping me, the thing that is splitting me into
subselves, that thing is an ever-increasing sense of the reality of that other ludicrously
simple act of button-pushing, the act that will destroy not a computer file but a real
state, in fact all our states, all our towns. It is the "like-it-or-not-it's-real" issue of
nuclear war. It's what has occasionally been called "the unthinkable".
        I have always liked that name for it. "The unthinkable"-it carries with it the
notion that it truly will never happen. That it is so awful that no one could ever
conceivably start a nuclear war. That nuclear war is synonymous with the end of the
world, with Armageddon. But it seems that such a view has faded, over the years. It
seems that "the unthinkable" is being contemplated more and more by the
governments of the nuclear powers..

                                        *       *       *

        I live in Bloomington, Indiana, a town of somewhere around 60,000
inhabitants. Somewhere in the Soviet Union, there is a missile silo with a missile
inside it destined specially for Bloomington, Indiana. After all, there are several
thousands of nuclear-tipped missiles (I love the delicacy of that word "nuclear-
tipped"!-one can almost visualize a cute little hood ornament on the tip of a gracefully
streamlined rocket, an artistic flourish added merely as an extravagant but stylish
afterthought)-several thousands of these that will form a first-strike force. There are
workers who have never heard of Bloomington, Indiana who daily do routine checks
to make sure this missile will hit its mark and do its duty. Some of them may know
where it is going. They pronounce it "Bluminktoan, Eendianna", and maybe then they
chuckle.
        Somewhere else, perhaps in Utah, perhaps not, there are American workers,
people very similar to those in the Soviet Union, who are taking loving care of very
similar missiles aimed at Gorky, Novosibirsk, Omsk, towns in the Soviet Union. They
pronounce those names with accents as atrocious as those of their Soviet counterparts.
        These people on opposite continents are very similar to each other, and bear
each other no ill will. Yet each missile is there' for no other purpose than to put out
the soul-flames of hundreds of thousands of human beings in. a distant town,
instantaneously or over several weeks.; The horror is not imagined.

                                        *   *   *

        The American people are known the world over for their generosity and
warmth. America has given the world many wonderful things. There is the notion of
freedom, the image of the western frontier, the Hollywood movies, jazz of many
types, a vision of how technology and science can make life increasingly pleasurable,
a looseness of language and dressing style that




The Tumult of Inner Voices                                                             784

have inspired people' the' world over, an informality and palsiness that are deep
products of the American character. I am proud to write for the magazine Scientific
American, to be part of that wonderfully characteristic American institution, and to
have my writings represent my country. I feel mentality is a uniquely American
product, and I am proud of it. This country has contributed greatly to humanity's self-
image; it has become a special symbol in the minds of people all over the world.
Similarly, the Russian people have contributed monumentally to world \_,culture,
through their novelists like Dostoyevsky and Tolstoy, through the noble music of such
composers as Rachmaninoff and Prokofiev and Scriabin ana Shostakovich (some of
my favorites), through their scientists and their mathematicians and philosophers,
through their rich and sad culture, so full of torment and resignation. In the town of
Gorky lives the towering spirit of Andrei Sakharov, often known as the father of the
Russian H-bomb, now a dissident struggling to turn around his government. In
ordinary towns all over the Soviet Union live ordinary individuals like you and me,
who desire nothing more than simple lives without a Sword of Damocles hanging
over them.
       But we all have this sword hanging over us, and the thread by which it is
hanging is getting thinner every day. The "unthinkable" has not only been thought
about, it has been planned in infinite detail, on both sides. Dozens of nuclear war
scenarios have been considered and weighed in the balance, on both sides. Dozens of
versions of Armageddon have been played out -either on computers or in the minds of
war planners. These people plant colored pins on maps and think only about numbers,
not about lives.
       Not only have plans-software-been drawn up, but of course, all the hardware,
the materiel, is there. It is in place. But as if that were not enough, each day, on the
surface of the earth, somewhere between three and five new atomic warheads are-
brought into being. We all know that only two atomic bombs have ever been dropped
on people, each one killing somewhere around sixty thousand people. And those two
were just tiny, "cute" bombs compared to the ones we are making every day.

                                       *   *   *

        There are some people who, ostrich-like, wish to think that nothing has
changed radically since the days of World War II. A friend of mine told me recently
of a lunchtime conversation he had with his boss. My friend had brought up the
horrors of nuclear war, and his boss-a kindly man with good intentions-protested that
in reality, a nuclear war wouldn't be much different from any previous war. "How do
you see that?" my friend asked. "What's the smallest nuclear weapon?" said his boss.
"Oh, about 500 tons of TNT", replied my friend, up on his statistics. "And the largest
conventional bomb?" asked his boss. "Maybe around 20 tons of TNT", said my friend.
"There-you see?" said the boss. "Only a factor of 25."




The Tumult of Inner Voices                                                          785

FIGURE 33-1. One common posture to take, in light of the distressing news
constantly bombarding us. [Drawing by David Moser. ]

         When my friend repeated this story to me, I couldn't believe my ears. My
friend's' boss felt comforted by the thought that the smallest nuclear bombs are "only"
25 times bigger than the largest conventional bombs. He was willing to completely
neglect a factor of 25-in itself preposterous. But 25 is not the right number at all. A
typical nuclear bomb is more like one megaton, which is to say more like 50,000
times bigger than the largest of conventional bombs. Not 25 times, but 50,000 times
larger. If you compare it to a typical, rather than the largest, conventional bomb, the
ratio becomes something on the order of a million. One million. The figure is
staggering,, literally incomprehensible. To make it more graphic, though certainly no
more conceivable, I can quote the statistic that a single nuclear warhead of moderate
size-2.2 megatons-carries more destructive power than all the bombs dropped on
Germany during all of World War II. The bomb destined for Bloomington, Indiana, is
in all likelihood about half that big. Maybe knowing that it's only half that big would
comfort my friend's boss. Yet somehow I doubt it. (See Figure 33-1.)

                                       *   *   *

       But why am I speaking to you about these ghoulish issues? Did I come here to
speak about nuclear war, or about the nature of the soul? Why am




The Tumult of Inner Voices                                                         786

I talking so passionately about issues that terrify rather than issues that fascinate and
enchant? Because I am more than one person. Because I keep this inner voice quiet.
Because something in me has woken up being dormant for many years.
         Now if I turn my attention to this strange notion of "inner voices" and dormant
ideas", I will have come right back to the other topic-the nature the soul. So my own
internal split forces me to oscillate back and forth on a topic I hate to contemplate and
one I love to contemplate. I am victim of my own mind, a mind that has woken up to
the nightmare of today's realities. Or seen on another level, my mind is "shoved
around" by internal agents of my brain that have become activated after a long
dormancy. And where my brain shoves, so must my mind follow. How can I be
saying these strange things? What is the meaning of "mind as slave to rain"?
         To confront this, let me shift gears, and talk about brains. A brain is a
collection of many, many parts-some 100 billion neurons. These parts are ked up to
each other in fantastically complex ways. Most neurons, for instance, are connected
up with several thousand other neurons, often quite far removed ones. Neurons come
in a few types, but basically, they are all quite similar to one another. Systems with
large numbers of identical parts have been studied for years in that branch of physics
known as statistical mechanics. For example, think of a lattice-a three-dimensional
lattice, like an enormous three-dimensional checkerboard-in which a particle can sit
with its "spin" pointing either up or down. Each particle directly affects only its
immediate neighbors through the magnetic field created by its spinning. So now each
particle must decide if it "wishes" to be pointing up or pointing down. Its "decision" is
deterministic, being governed by the states of its neighbors. But the strange thing is
that their states are in turn governed by their neighbors, and so the whole thing turns
out to be interconnected in a vast interlocked way, despite the fact that any given
particle "feels" only its immediate neighbors.
         As a, consequence, the behavior of such a system has some striking properties
that go under the general heading of "collective phenomena". This is an elusive
notion. An example is the common phenomenon of the traffic jam: You won't locate a
traffic jam if you restrict your search to the insides of a single taxi. After failing to
find it inside the glove compartment, you move on to the trunk, and then you look
inside the hood, opening up the battery, the radiator, then moving on to the gas tank
.... A ludicrous image! A traffic jam is just not on the level of an individual car. It is a
pattern composed of cars, a pattern that moreover has deep repercussions on the cars
it is composed of.
         The nature of collective phenomena is that they are patterns composed of
parts, and they in turn exert powerful influences on their parts, acting to keep them in
line. Think of hurricanes, life, intelligence. You don't see a hurricane by looking
inside an atom of oxygen. You don't see life by looking at an amino acid. Nor do you
see intelligence by looking at one isolated




The Tumult of Inner Voices                                                             787

teeny part of its substrate, whether that part is a synapse, as in a brain, o a. binary
digit, as in a machine.

*   *   *

        What is the dynamics of a collective phenomenon? How does a thought get
generated and propagate inside a brain? Nobody really knows. All we can do is resort
to simpler analogies. For instance, imagine a school of fish swimming through the
sea, when the lead fish encounter a danger: a strange shadow passing overhead, or a
sudden movement ahead. In a split second the whole school has wheeled about in
unison, and is hightailing it away at top speed. Such collective actions, such
cooperative effects, take place when all the components are in phase with each other.
They act as a unit. In a similar way, the spins in a magnetic substance act in concert.
If certain pivotal spins flip, a whole "school" of spins-a magnetic domain-flips in
unison.
        These so-called magnetic domains are a little like countries, and their
individual spins are like people. Of course, people are much more complicated than
spins or particles are. They do not just have two states, up or down! However, in
times of crisis, people are often placed in just such positions, of having to see a
situation in very black-and-white terms. I remember vividly how Iran seemed to
switch overnight from a country favorably inclined toward the United States-and I
mean at Iran's grass-roots level, not just at its governmental level-to a country full of
bitter hatred toward this country. Just all at once it became polarized-a term,
incidentally, applying to magnetic substances as well, in which all the spins (or a
substantial majority of them) point in one direction.
        And external events can swing these polarizations the other way in dizzyingly
short times. We all know how countries can seemingly "flip" in very short periods of
time. (Think of Argentina's amazingly sudden near-total unification behind a junta
that, up until a few weeks ago, was bitterly hated!) But such a thing can happen only
if the country is polarized or otherwise coalesces so as to act as one large collective
entity. During much of any country's existence, it does not act as a single, polarized,
black-and-white entity, but rather as a much more complex entity composed of dozens
and dozens of smaller entities-domains-all of which are pursuing their own goals and
coexisting in one way or other with each other.
        A brain is much like this. A brain, with its billions of neurons, resembles a
community made up of smaller communities, each in turn made up of smaller ones,
and so on. The highest-level communities just below the level of the whole are what I
like to call "subselves" or "inner voices". I mean the latter not in the sense of a
schizophrenic who hallucinates inner voices, but rather, in the sense of competing
facets that try to commandeer the whole, something like hijackers, although often
benevolent hijackers. Perhaps it is more like a vote of the passengers on a plane where
they want to go, after it is in the air!




The Tumult of Inner Voices                                                           788

         To give an example of such commandeering, one that I pursue in=so depth in
that Achilles-Tortoise dialogue on the careenium, consider what happens when I pick
up a Rubik's Cube. What makes me pick one up, in the first place? I see it sitting there
on the shelf, glittering just as always-but once in a while its glittering catches my eye
a little more than usual.. Somehow, I am "primed" for a cubing session. So I pick it.
up. I may have a million other more pressing things to do than solve the cube, but
somehow I cannot resist picking up that cube and scrambling it, and then twisting its
colorful faces one way and another, very enjoyably, back to order. And having done it
once, I do it over again. Then I may say to myself, "Really! You have better things to
do! Just once more!" And so I do it once more, but somehow that is not satisfying
enough, and with a slightly guilty feeling, I do it again. And again.
         As the old potato-chip ad said, "Betcha can't eat just one!" I am also reminded
of the Academie Francaise, which tries to keep a tight hold on the French language, to
prevent the people from following certain natural trends that are apparent to everyone.
But the Academic tries to exert what I call "top-down pressure" to prevent things that
obviously cannot be prevented. There is too much "bottom-up" momentum to put the
lid on. You cannot keep the lid on a boiling pressure cooker. You cannot keep the lid
on an angry populace, such as the Iranian people or the Salvadoran people or the
Chilean people. No matter how tightly you press down, the pressure from inside will
in the end overwhelm you.

                                       *   *   *

        And thus it is within the brain. Competing subselves cannot be held in check
indefinitely. They cannot be clamped down, forbidden to act. For each "inner voice"
is in actuality composed of millions of smaller parts, each of which is active, and
under the proper circumstances, those small activities will someday all "point in the
same direction", and at that moment the inner voice will crystallize, will undergo what
is called a phase transition, will emerge from obscurity and proclaim itself an active
member of the community of selves. And if it is powerful enough, it will try to exert
pressure and to be recognized. It will attempt to seize power. It will not want to
relinquish power, once it has it. That's what I mean by "commandeering the soul".
        I attempted to illustrate this with the Rubik's Cube example, but I see it
happening in me all the time. I have a "piano-playing subself", who, once he is given
the floor, refuses to relinquish it for hours on end-until, say, my back-his back?-grows
achy, or until he gets sleepy. Or until the phone rings or my watch beeps at me, telling
me that some other facet of life must be attended to.
        And somehow, in such circumstances, there is a governing personality who
can grab control away from the "hijacker". In fact, it is not at all hard to dislodge the
piano-playing hijacker, or the cubing hijacker, or any other subself, when a phone call
comes. Isn't that a sign of the times? This kind




The Tumult of Inner Voices                                                           789

of selective interruptibility is one of the most critical characteristics of a hierarchically
organized system such as a brain, which has evolved to deal with a world in which
there are events of various priorities that have to be dealt with sequentially. Choices
must be made, so there must be a highest-level body whose purpose is to make
choices rapidly and reliably, -one that sorts out the priority of subselves and allows
only the one deemed most important to take charge.
        An interesting problem arises when even this deciding agent must be
preempted by some sort of extreme emergency situation that arises. The decider can
be in the midst of trying to sort out some ordinary conflict of subselves when—

                                         *   *   *

          Believe it or not, I was just interrupted by a phone call from a friend. Just as
well, I think. For it gives me the opportunity to review where I am going and to
resume some of the ideas that I had left off in the middle. In particular, it allows me to
make the analogy that I wished to make between a country that is dormant and a
subself that is dormant. In each of us right now-and I am quite confident of this-there
are competing inner voices, perhaps one of them dormant, but still present, in some
implicit sense-that say opposite things about; the prospect of nuclear war. One of-
them says something like "Nuclear war certainly would be bad, in fact, unthinkable-
and everyone knows that, including all the military people and all-so there will never
really be a nuclear war, it's all just a way of maneuvering and bluffing and so on. So
I'll just go on with everyday life." This -is the voice in me that has been dominant for
years and years! And then there is the other voice, the one that tries more seriously to
envision the true nature not only of the devastation that such a war would cause, but
also that tries more seriously to evaluate what is going on in the proliferation of
nuclear arms.
          Certainly it is hard to keep shutting one's eyes to the fact that all over the
world, people are thinking about "the unthinkable". Some people are talking about
how unthinkable it is, but others are talking about how it might be thinkable, after all.
There are those who are telling you how to build a shelter, how to stockpile food, how
to keep a gun to ward off neighbors, friends, and strangers when they. try to burst into
your safe little haven. There are those who are talking about evacuating whole cities
into rural areas-as if we'd have the time to do such an incredible thing, or as if people
would stomach the idea. These people-people involved in civil defensehave a vested
interest in reassuring us all that nuclear war is indeed conceivable, survivable-not all
that bad, in fact! What an incredible kind of job to have.
          But worst of all, there are those who seem blind to the idea that nuclear war
would truly spell the end of the world as we know it.. There are millions of ordinary
citizens who are somehow relieved when they see a map of their




The Tumult of Inner Voices                                                              790

city with the various circles drawn around the downtown square, because they notice
that where they live is outside the "90 percent killed" circle, even outside the "50
percent killed" circle. So no need to worry. They'll survive., And that's as far as they
choose to think about it! Some vague apprehension, maybe, about fallout, or difficulty
of getting gas for the car, but that's about all.
        Now I shouldn't really be accusing some people of thinking this way. The
strange thing is that we all tend to think this way-or at least parts of us do. (At least I
can speak for myself, and I think I am a very typical person.) For we are dealing with
something that not only is very vague and unknowable, but also something that is
unimaginably catastrophic, something the likes of which has never happened on this
planet. So we are not equipped to imagine it (but see Figure 33-2 for some help). And
so we turn off. And this turning-off happens to some extent in each and every one of
us. Certainly it has been the dominant mode in me for many years. One develops and
encourages a sense of security in the ridiculousness of nuclear war.
        But the stockpiles are increasing every day. The dangers are increasing every
day. The warmongering' talk is increasing every day. The number of flashpoints
around the world is increasing every day. The mistrust and suspicion and polarization
of peoples is increasing every day. The only thing we have on our side is the hope that
apathy is not increasing. The hope that a country as large as our own can itself
undergo a "phase transition", an awakening, a realization of the insanity of the course
on which we are embarked.

                                        *       *       *

        Phase transitions take place in simple physical systems, schools of fish,
individual brains, and in countries as well, when there are sufficiently strong and
numerous interactions between the components of the, system, and when those
interactions add up in such a way as to make for large-scale correlations, or, put
another way, long-distance effects despite the short-range nature of the direct
interaction. When such long-distance effects occur, then a new kind of entity springs
up, an entity on a higher level of organization than its constituents, and that entity
obeys certain laws of its own.
        Performers are highly aware of this collective aspect of crowds, for instance. A
singer will speak of the interaction between herself and the crowd, of how she senses
the mood of the crowd as a whole. Yet how can this be? Isn't a crowd composed of
individuals who are totally unknown to each other, individuals; with nothing in
common? Yes; however, they do have one thing in common: they are all there,
physically, listening- to the same performer, and so they are influencing each other
whenever they laugh at her jokes or applaud her, or encourage her or seem impatient
in any way. Such collective modes tend to lock in very quickly, to create self-
reinforcing loops of interaction between performer and audience.




The Tumult of Inner Voices                                                             791

FIGURE 33-2. A realistic view of the world armaments situation. The chart shows the
world's current firepower in terms of the firepower of World War II. The dot in the
center square represents all the firepower of World War H (including the atomic
bombs dropped on Hiroshima and Nagasaki): three megatons. The other dots
represent the world's present nuclear weaponry. This comes to 18, 000 megatons,
which equals 6, 000 World War H's. The U. S' (and allies) and the S. U., (and allies)
share this firepower approximately equally.
        The top lefthand circle enclosing nine megatons represents the weapons in just
one Pas submarine. This is " l to the firepower of three World War II's and is enough
to destroy over 200 of the Soviet Union's largest cities. We have 31 suck submarines
and ten similar- Polaris submarines, The bottom lefthand circle enclosing 24
megatons represents one new Trident submarine with the power of eight World War
H's: enough to destroy every major city in the Northern hemisphere. The. Soviets have
similar levels of destructive power.
        Just two squares on this chart (300 megatons) represent enough firepower to
destroy all the large and medium-size cities in the entire world. [Designed by Jim
Geier and Sharyl Green in 1981. )




The Tumult of Inner Voices                                                        792

        Such self-reinforcing "loops' ate of the essence `in phase transitions collective
modes, for they are what tend to keep the whole thing going. And thus it is with the
collective mode of my neurons, the one that has somehow gotten triggered into
activity after many years of dormancy. This new inner voice is one that I am not yet
entirely comfortable living with. But it is one that haunts me and will not leave me
alone. It has seized some power inside me and it will not let go. And the
"government" that it has to some extent usurped is not entirely displeased with the
state of affairs.
        Let me now try to return control to the more dispassionate and objective "top-
level" self who began this talk.

                                       *   *   *

        Thank you. It has been interesting to me to observe the flipping back and forth
that has taken place as the previous subself tried to express himself. One thing that it
clearly shows is that there are no clear boundary lines to be drawn between "that
subself", "this subself" and any other subselves of Doug Hofstadter the person. All of
them are fictions, because the only real thing is the sum total, the integrated person.
And that integrated person is clearly not the same person he was a few months ago,
when he was blithely ignoring the notion of nuclear horror, somehow unwilling to
face the possibility squarely.
        This phase transition has not been an entirely pleasant thing to undergo; no
more than any coup d estat would be. Not that it was so revolutionary. It, all arose
peacefully, nonviolently, from within. There were no provocateurs from without. Or
perhaps, I should say, there was one-a 92-year-old lady who was my neighbor last
year, and whom I befriended. I would visit her on occasion, and we would have
wonderful conversations that rambled from the music of Chopin to the pangs of sad
romances to the secrets of the mind and-once in a while-to politics.
        One day as I was leaving after one of these discussions, Hildegarde said to me,
in a very gentle way, "One thing I'd like to ask you someday is how it is that with
your very alert mind, you don't seem to feel the need to do something-or to try to do
something-about nuclear war." It was a very gentle nudge, really only a passing
remark indicating her puzzlement about me. But it did set me to wondering how it
was that I could systematically ignore the biggest thing in all our lives, day after day
after day.
        Partially, the reason that I gave to myself was that it was just too big. There
was no use in worrying about it. But that rang phony to me. It didn't sound like me!
So actually, I had no answer, and that realization began to eat at me. It began to feel
like either pure "ostrichism", or pure egotism. Either way, I didn't like it. But a sense
of shame or guilt is never the way to bring about a phase transition. It's got to come
from somewhere far deeper than that. And fortunately, there were seething, churning
forces down deep inside me




The Tumult of Inner Voices                                                           793

slowly aligned, slo ly started to bring about that collective'mode that crystallized in
the inner voice that you heard.
        This "waking up" of an individual has its parallel in the collective waking up
of a nation. It will happen when enough citizens band together, seeing some common
interest, sensing some common goal. There is a sort of S,'critical point" when that
number reaches a threshold and suddenly, there is a turnaround at a national level. But
just how or when that will happen is very tricky to say.

                                       *   *   *

       In a recent book entitled Science: Good, Bad, and Bogus, Martin Gardner
quoted a beautiful passage, written by psychologist William James roughly a century
ago, about the act of waking up and rising in the morning. This passage captures for
me something very deep about the way the soul of a person arises from a myriad
smaller actions that are completely unknowable and yet that are somehow
coordinated. I would like to quote that passage to you:

          We know what it is to get out of bed on a freezing morning in a room
  without a fire, and how the very vital. principle within us protests against the
  ordeal. Probably most persons have lain on certain mornings for an hour at a
  time unable to brace themselves to the resolve. We think how late we shall be,
  how the duties of the day will suffer; we say, "I must get up, this is
  ignominious," etc.; but still the warm couch feels too delicious, the cold outside
  too cruel, and resolution faints away and postpones itself again and again just as
  it seemed on the verge of bursting the resistance and passing over into the
  decisive act. Now how do we ever get up under such circumstances? If I may
  generalize from my own experience, we more often than not get up without any
  struggle or decision at all. We suddenly find that we have got up. A fortunate
  lapse of consciousness occurs; we forget both the warmth and the cold; we fall
  into some revery connected with the day's life, in the course of which the idea
  flashes across us, "Hollo! I must lie here no longer! "-an idea which at that
  lucky instant awakens no contradictory or paralyzing suggestions, and
  consequently produces immediately its appropriate motor effects. It was our
  acute consciousness of both the warmth and the cold during the period' of
  struggle, which paralyzed our activity then and kept our idea of rising in the
  condition of wish and not of will. The moment these inhibitory ideas ceased, the
  original idea exerted its effects.

        I find this to be a remarkably perceptive passage, so accurate in its
understanding of the way people really work. In the "Careenium dialogue, l expressed
some similar ideas of my own, with which perhaps it would be appropriate to
conclude my talk tonight.
        Achilles says to the Tortoise, "In, emotionally wrenching cases, you can
hardly decide what you will feel. Something just happens inside you. Subtle forces
shift deep inside you, hidden, subterranean. It's quite scary, in a way,




The Tumult of Inner Voices                                                         794

because in real crises like that; instead of being able to decide how you'll a you find
out what sort of stuff you're made of. It's more passive than active -or more accurately
put, the action is on levels of yourself that are far lower -far more microscopic-than
you have direct control over."
        The Tortoise replies, "Correct. You and your neurons are not on speaking
terms, any more than a country could be on speaking terms with its citizens. There is,
in both cases, a kind of collective action of a myriad tiny elements on low levels that
tips the balance-exactly as in a country that `decides' to go to war or not. It will flip or
not, depending on the polarization of its citizens. And they seem to align in larger and
larger groups, aided by communication channels and rumors and so on. All of a
sudden, a country that seemed undecided will just `swing' in a way that surprises
everyone."
        Achilles continues, "Or, to shift imagery again, it's like an avalanche caused
by the collective outcome of the way that billions upon billions of snow crystals are
poised. -One tiny event can get amplified into stupendous proportions-a chain
reaction. But the crystals have to be poised in the right way, otherwise nothing will
happen."
        And Mr. Tortoise takes over: "In cases of judgment, whether it be of one
musical composer over another, one potential title or subtitle for a book over another,
or whatever, the top level pretty much has to wait for decisions to percolate up from
the bottom level. The masses down below are where ; the decision really gets made,
in a time of brooding and rumination. Then the top level may struggle to articulate the
seething activity down below, but those verbalized reasons it comes up with arc
always a posteriori. Words alone are never rich enough to explain the subtlety of a
difficult choice. Reasons may sound plausible but they are never the essence of a
decision. The verbalized reason is just the tip of an iceberg. Or, to change images,
conflicts of ideas are like, wars, in which every reason has its army. When reasons
collide, the real battleground is not at the verbal level (although some people would
love to believe so); it's really a battle between opposing armies of neural firings,
bringing in their heavy artillery of connotations, imagery, analogies, memories,
residual atavistic fears, and ancient biological realities."
        Finally, Achilles exclaims, "My goodness, it sounds terrifying! You make the
battlefield of the mind sound like a • vast mined battlefield! Or a treacherous ice field
on a steep mountain face. I never realized that a mechanistic explanation of thinking
could sound so organic and living. It's sort of awful and yet it's sort of awe-inspiring
as well."

                                         *       *       *

        Achilles' remarks hit the nail on the head, for me. Life, when you contemplate
its basis in biology, is in many ways terrifying; yet there is a kind of majesty to the
depth and complexity of it all. The same holds for humanity as a whole. In many
ways, we are a shocking bunch, doing the most




The Tumult of Inner Voices                                                             795

terrible things to each other and to other living-beings-,'yet there is also an element of
the sacred in humanity, something sacred in spite of the profane in each one of us.
        The pile of contradictions that each one of us is still often adds up to
something beautiful and cherishable. To preserve that sacred and beautiful facet from
the menace created by the profane and awful facet is worth every effort that we can
muster, drawing on the power of the many subselves and inner voices that resonate
within us and make us what we are.




The Tumult of Inner Voices                                                            796

                                      Epilogue
        After writing such a long book, I have a long list of people to whom I owe
genuine thanks for many different reasons. I find it very hard to draw the line between
people who have contributed directly to this book, and people whose contributions,
though real, are indirect. Yet I must attempt to do so, for the sphere of indebtedness,
extends out hazily to encompass\_ practically everyone I know. In what follows, I
shall do my best.

        To begin with, I should like to thank Dennis Flanagan and Gerard Piel for
offering me the opportunity to write for their distinguished magazine. Each month, I
worked with Dennis on the microscopic level of the columns, and I , thank him for his
good judgment. Though we had our share of disagreements, we developed a warm
friendship that I value.
        Martin Gardner suggested that I might be his successor. To 'be recommended
by someone of Martin's honesty, wit, and insight is a very high compliment. Thank
you, Martin, for that and for all the wonderful things you have written and continue to
write.
        This book was written in many places. The first few columns were written
when I was a john Simon Guggenheim Fellow visiting Stanford University's
Computer Science Department, and I would like to thank the Guggenheim Foundation
for its support. The majority of the columns were written in Bloomington, Indiana,
where for seven years I have been on the faculty of the Computer Science Department
of Indiana University. Some new material was written while I visited the Institute for
Cognitive Science at the University of California at San Diego in early 1984, and the
rest at MIT's Artificial Intelligence Laboratory in Cambridge, Massachusetts, where I
spent most of my sabbatical year.
        I would particularly like to thank two of my hosts. Donald Norman made me
feel most welcome at UCSD's Institute for Cognitive Science. I enjoyed not only all
the facilities there, but also a couple of runs along the beautiful Del Mar ocean front
with Don. At MIT, I was truly lucky to have the interest and support of Marvin
Minsky, who went out of his way to make my stay especially comfortable. He even
supported two people working .with me, something I will never forget.
        To Indiana University, however, I owe the most. IU offered me a job in





1977 when I had little to show by way of achievements in cognitive science. Since
then, my department has been an extremely supportive and friendly environment. I
would like to thank several close colleagues, whose friendship I cherish: Dan
Friedman, John O'Donnell, Frank Prosser, Cindy Brown, Mitch Wand, Dave Wise,
Paul Purdom, Ed Robertson, Stan Kwasny, Bob Filman, Will Clinger, George
Epstein, and the three JB's: Jim Burns, John Buck, and John Barnden. I have
exchanged idea$ with all of them, and together, they have markedly influenced this
book. The Computer Science staff has also been a joy to work with over the years. I
would like to single out Kathy Thompson, whose spunk and wacky humor have
brightened many a dismal day.

        In Bloomington, I have made friends too numerous to mention. As happens in
any university town, many of them have left, but they have all made Bloomington a
special and wonderful place to live.
        The tremendous interchange of ideas I've had with Don Byrd over these past
seven years is reflected on all scales of this book, and the generous companionship he
has offered is reflected on all scales of its author's life.
        Two friends whose intellectual influence on me has been profound are Gray
Clossman and Marsha Meredith. But even if their intellectual influence had been nil,
they have been friends in need, friends in deed. For that I have to thank them deeply.
        Ann Trail, with her sparkling sense of humor, her optimism and generosity,
and especially her sense of mortality, has deeply and permanently enriched my life.
This book reflects her style in so many ways.
        Other Bloomington friends have madee such a difference as well. John and
Joanie Woodcock have long been close friends, and have always been warm and
lively conversation partners. With Scott and Ruth Sanders I have shared political
hopes and disappointments, and many exuberant discussions. The Leake family-Roy
and Alice, David and Patsy-have been true friends, full of interest and empathy. Ruth
Sonneborn and her late husband Tracy were among my very first Bloomington
friends, and I will never forget evenings spent at their house engaged in delightfully
passionate philosophical arguments. I have relished many consonances and
dissonances over musical matters with Al and Helga Winold. Over stimulating
lunches with Mike Dunn and others, I gained a new kind of respect for philosophers.
It has been my privilege to know University Chancellor Herman B Wells, who, if it
could be said of anyone, is the soul of Indiana University.
        I have shared, enthusiasms for all sorts of ideas-often over meals or coffee-
with other Bloomington friends of then and now: Ann McMillan, Sue Wintsch, Vahe
Sarkissian, Adrienne Gnidec, Al and Linda David, Tulle Hazelrigg, Jimmy my and
Gilan Tocco, Judy Mahy and Rich Shiffrin, Marlene Mannella and Evan Smith, Tom
Ernst, John Goldsmith, Enrico Predazzi, Marion O'Connor, Sujan Yang, Vicky
Grossack, and Anneke Campbelland the list goes on. I wish I could say something
about each one of them,





but that would take a whole book--and' then I would have to write ih
acknowledgments to that book, which could get to be a problem.

        When I was at Stanford in 1980-81, I benefited from contact with many
people. Scott Kim, as usual, was full of inspiration and creative energy. I shared
Mexican food and long talks with Pentti and Dianne Kanerva. Conversations with
Scott Buresh, Marcia Bianchi, Debbie Schweninger, Louis Mendelowitz, Louella
Kates, Liz Powers, Allen Wheelis, Larry Breed, Margie and Sia Khosrovi, Eric
Hamburg, Debbie Starbuck, Fanya Montalvo, and Stan Isaacs made my Guggenheim
year at Stanford a lively and memorable one.
        My upstairs neighbor that year was a most remarkable woman in her 90's,
named Hildegarde Kneeland. An economist by profession, Hildegarde taught many of
today's most influential economists. Even today, she is passionately concerned about
the fate of humanity. Hildegarde touched off the fire in me concerning nuclear
madness. I wish I could do for others what Hildegarde did for me.
        Several of my family's oldest and dearest friends have died in these past few
years: Dan Mendelowitz, George Feigen, and Felix Bloch. I had known them all since
I was very small, and each has left indelible tracks in my soul. Traces of Dan, George,
and Felix lurk throughout this book.

        My love for writing and alphabets was heightened by my grandmother, Mary
Givan, who shared her love for letterforms with her grandson by showing him
wondrous and eye-opening alphabetic books. A little later, in 1960, my uncle and
aunt, Albert and Manya Hof§tadter, introduced their 15-year-old nephew to abstract
art at New York's Guggenheim Museum and Museum of Modern Art, and though I
protested how silly it all was, they taught me things that changed my way of looking
at visual forms. These experiences started me down many of the artistic and scientific
pathways described herein.
        I have known Ernest and Edith Nagel for almost 25 years, and they remain a
beacon of sanity in this crazy world. Many echoes of fondly remembered
conversations are found in this book.

        During my enjoyable stay in San Diego, I fruitfully exchanged memes with
Paul Smolensky, Don Norman, Dave Rumelhart, Larry West, Karen Pickens, Wendie
Maurer, Liam Bannon, Larry McGilvery, and many others. I would also like to thank
the ICS staff for making things work smoothly while I was there. It's a superlative
place to do cognitive science.
        While in the Boston-Cambridge area, I was overwhelmed by the number of
good friends I have there. Gloria Minsky is one of the world's warmest ,people, and it
is such a pleasure to enjoy a "Min Chin Din" with her and Marvin and whoever
happens by. Betty Dexter is much more than a fantastic secretary-she is a great friend.
Her laughing presence on the seventh floor





Of the Al Lab was a delight. Dan Dennett, as always, was-bubbling over with ideas
and enthusiasm. I just wish I had gotten to see more of Dan and his wife Susan.
        Working with me at the Lab were Marek Lugowski and Melanie Mitchell,
graduate students, and David Rogers (also known affectionately as "Dr. it Ogers"),
post-doc. It is so much fun to bat about research ideas with people who are as
intrigued with them as I am!
        Some other people who made the Boston atmosphere so exciting are Henry
Lieberman, Bernie Greenberg, Chris Morgan, Greg Huber, David Levitt, John
Amuedo, Marek Holytiski, Margaret Minsky, Joe Shipman, Russell Brand, Randy
Davis, Fanya Montalvo, Carl Hewitt, Jay McClelland, and Dedre Gentner.

       Scattered around the globe are numerous other friends who have contributed to
my thoughts and moods over the past few years. I'd like to mention Charles Brenner,
David Policansky, Mary Adele and Norman Mather, Elwyn and Darlene Wolcott,
Pete Rimbey, Francisco Claro, Inga Karliner, Marek Demianski, Maria Nosowska,
Zamir Bavel, Peter Suber, Phil and Sarah Taylor, Bob Wolf, Len Shar, Dorothy and
Peter Denning, Betty and David Hamburg, and Piet Hoenderdos.
       Of special interest are three translators of Godel, Escher, Bach into other
languages: Bob French and Jacqueline Henry (into French, of course), and
       Ronald Jonkers (into Dutch). In the summer of 1983 in Paris, we thrashed out
many tricky translation problems, and those stimulating sessions contributed to my
understanding of the connection between translation and analogy.
       For lending their expertise in specific columns or articles, I would like to
thank Bob Axelrod, Bill Huff, Dave Martin, Merald Wrolsted, David Singmaster,
Mitch Feigenbaum, Dan Mauldin, and Paul Stein. At Scientific American, I was
helped numerous times by Adele Premice, Brian Hayes, Sally Jenks, Mary Knight,
and Sam Howard.
       I am most pleased to be doing this, my third book, with Basic Books once
again. Martin Kessler, president, was interested in such a book from the moment he
heard about the column,. and has been very supportive. My day-to-day contact has
been with Maureen Bischoff, just as with my earlier books, and I must say, I get a
great kick out of our lengthy telephone calls. Maureen's rapier wit takes the tension
out of many difficult situations. Other
       people who have helped a great deal at Basic include Vincent Torre, Ellen
       Prior, Elizabeth ,Werter, Linda Carbone, Ann Rudick, Thalia Doukas, Ruth
Elwell, and Jeremy Orgel. Thanks also to Debra Manette, Sandra Dohls, Sabrina
Soares, Michael Wilde, David Graf, John Masur, Kathi Lee, Donna Singer, Lisa
Adams, and John McAusland.

       As always, my family looms large in my life. My parents, Robert and Nancy
Hofstadter, have constantly served as' critics and supporters. My love for them truly
cannot be expressed. My sister Laura is a wonderful person, full





of ideas and humor. Her company is cherished. My sister Molly; for unknown reasons
unable to talk or understand, has all my love. Her sad plight was a deep mystery that
made me start to wonder, many years ago, about mind, brain, and soul.
        Finally I come to two special people without whom I would be a very different
person today. David Moser has not only been my closest consultant on this book; he
has also been like a brother: empathetic, generous, and caring' 'Carol Brush has been
the light in my life over the past two years. We have gained so much from each other,
and deeply changed each other's lives.

        Taking leave is so hard, and I am going to leave Bloomington soon, for other
pastures. A unique constellation of colleagues at the University of Michigan hoped
that they could find a way to have me join them. Largely through their efforts, I was-
offered the Walgreen Chair in Human understanding-a rare opportunity. After
brooding on it for a while, I concluded that this was just too good to pass up, and so,
with feelings of excitement tempered by sadness, I accepted the offer.     -
        As any reader can tell, what I have written here is much more than simply
acknowledgments for this book: I wanted to thank the many people who have "been
there". I also wanted to express special feelings of affection for Bloomington, Indiana,
a town where I have flourished in every aspect of life over the past few years. These
acknowledgments are a sentimental "thank you" and "farewell to Bloomington. I will
certainly miss Howard's, the Spoon, the Grind, the Horn, and the Harmonica, but it is
time for me to move on to Ann Arbor, and; as Hildegard Kneeland once said so
memorably to me, to "welcome the future".

                                                                           -D.R.H.
                                                        Bloomington, November 1984.





                                        Bibliography
Anderson, Alan Ross; ed. Minds and Machines. Englewood Cliffs, NJ.: Prentice-Hall, 1964.
 A classic collection of stimulating articles on the mind-body problem, including Alan
 Turing's important paper "Computing Machinery and Intelligence", and J.R. Lucas'
 provocative "Minds, Machines, and Gödel".
Applewhite, Philip. Molecular Gods: How Molecules Determine Our Behavior. Englewood
 Cliffs, NJ.: Prentice-Hall, 1981. A book whose subtitle helped spark my "Careenium"
 dialogue (Chapter 25).
Atlan, Henri. Entre le cristal et la fumee: Essai sur l organisation du vivant. Paris: Editions du
 Seuil, 1979. A biologist who views life as the emergence of ordered complexity out of
 randomness here explains his philosophy.
Axelrod, Robert. The Evolution of Cooperation. New York: Basic Books, 1984. A beautifully
 written account of how mutually beneficial behavior-that is, cooperation-can emerge among
 purely egoistic organisms that share an environment over time.
Ayala, Franciscp Jose and Theodosius Dobzhansky, eds. Studies in the Philosophy of
 Biology: Reduction and Related Problems. Berkeley: University of California Press, 1974.
 The proceedings of one of the most fascinating conferences I have ever heard about. Many
 great biological thinkers were present, and discussed the most fundamental questions about
 how life and minds can be reconciled with physical law. I wish I had been therel
Bandelow, Christoph. Inside Rubik's Cube and Beyond. Boston: Birkhfiuser, 1982- A clear
 and mathematically oriented book about the Cube and its successors. Perhaps the best of all
 the Cube books in English.
Barr, Avron. "Artificial Intelligence: Cognition as Computation". In The Study of
 Information: Interdisciplinary Messages, edited by Fritz hlachlup and Una Mansfield. New
 York: Wiley-Interscience, 1983. A paper putting forward the orthodox Al dogma. that
 mental activity is "information processing"-more specifically, that the manipulation of
 "symbols" (representational data structures) by suitable computer programs is no more and
 no less than what minds do.
Beck, Anatole and David Fowler. "A Pandora's Box of Non-Games". In Seven Years of
 Manifold, . edited by Ian Stewart and John Jaworski. Nantwich, England: Shiva
 Publications, 1981. An amusing collection of silly pseudo-games with more of a moral than
 might first meet the eye.
Beckett, Samuel. Waiting for Godot. New York: Evergreen, 1954. The classic existential
 drama in which meaninglessness pervades-and at whose core is the famous nonsensical
 verbal vomit emitted from the mouth of the sad sack named "Lucky".
Benton, William. Normal Meanings. Paducah, Ky.: Deer Crossing Press, 1978. A book of
 strangely evocative poems, a bit comprehensible in parts, totally incomprehensible in
 others.
Bernstein, Leonard. The Joy of Music. New York: Simon \& Schuster, 1959. An exciting
 medley of ideas by this exuberantly articulate thinker and musician. His dialogues are
 especially enjoyable.
Biggs, John R. Letterforms and Lettering. Poole, England: Blandford Press, 1977. One of the
 best books on the fluidity of letterforms I have run across. Includes short sections on other
 languages, such as Hebrew, Chinese, and Arabic.
Blesser, Barry et al. "Character Recognition based on Phenomenological Attributes". Visible
 Language 7, no. 3 (Summer 1973). An early article by researchers who clearly had come to
 appreciate the depth of the letter-recognition problem.
Bloch, Arthur. Murphy's Law; Murphy's Law, Book Two; Murphy's Law, Book Thres. Los
 Angeles: Price/Stern/Sloan, 1977, 1980, 1982. Humorous and cynical observations about
 the human condition, featuring many self-referential or self-undermining aphorisms.
Boeke, Kees. Cosmic View: The Universe in 40 jumps. New York: John Day, 1957. A book
 to instill humility and awe in anyone, as well as a vivid sense of the meaning of the term
 "astronomical number".

Bibliography                                                                                 802

'Bombaugh, Charles Carroll. Oddities and Curiosities of Words And literature, edited
  annotated by Martin Gardner. New York: Dover, 1961. An antique collection of
  palindromes, acrostics, pangrams, and so on, for wordmongers and people who love the
  bizarre fringes of language. (Gardner's footnotes at the back are the best part of the book.)
Boole, George. The Laws of Thought. New York: Dover, 1961. A reprint of the old classic
  from the mid-1850's. The hubris of the title is quite remarkable, especially in light of what
  130 years' progress has revealedl
Brams, Steven, Morton D. Davis, and Philip Strafin, Jr. "The Geometry of the Arms Race".
  International Studies Quarterly 23, no. 4 (December 1979): 567-88. Looking at international
  behavior in terms of the iterated Prisoner's Dilemma and similar payoff matrices.
Brilliant, Ashleigh. I May Not Be Perfect, but Parts of Me Are Excellent: I Have Abandoned
  My Search for Truth, and Am Now Looking for a Good Fantasy; Appreciate Me Now and
  Avoid the Rush ; I Feel Much Better, Now That I've Given Up Hope. Santa Barbara, Calif.:
  Woodbridge Press, 1979-1984. Four books containing many incisive epigrams about life,
  death,, love, relationships, greed, egotism, loneliness, fear, and so on. None is longer than
  seventeen words.
Bush, Donald J. The Streamlined Decade. New York: George Braziller, 1975. Showing how
  the style of an era permeates its creations in all media. Full of elegant photos. Compare with
  the books by Loeb and McCall (see below).
Byrd, Donald. "Music Notation by Computer". Ph.D. thesis, Indiana University Computer -
  Science Department. Bloomington, 1984. About the problems of developing a computer
  program that will have some "understanding" of the subtleties of music notation. The
  program, SMUT, produced the musical examples in this book.
Chaitin, Gregory. "Randomness and Mathematical Proof". Scientific American 232, no. 5
  (May 1975): 47-52. An enlightening way of defining the meaning of "random pattern", and
  the unexpected and deep resonances with GSdel's incompleteness theorem and other
  metamathematicat results.
Charniak, Eugene, C. K. Riesbeck, and Drew V. McDermott. Artifictal Intelligence
  Programming. Hillsdale, NJ.: Lawrence Erlbaum Associates, 1980. Sophisticated ways of
  using Lisp and Lap-like languages in artificial-intelligence research.
Collet, Pierre and Jean-Pierre Eckmann. Iterated Maps on the Interval as Dynamical Systems.
  Boston: Birkhauser, 1980. An in-depth study of the iteration of simple smooth functions on
  the interval [0,1], and the resulting roads to turbulent behavior as parameters are varied.
Compugraphic Corporation. A Portfolio of Text and Display Type. Wilmington, Mass.:
  Compugraphic Corporation, 1982. A collection of many typefaces (mostly book faces),
  including several extensions to other alphabets of faces originally conceived only for our
  alphabet.
Conway, John Horton, Elwyn Berlekamp, and Richard K`Guy. Winning Ways (for your
  mathematical plays).' New York: Academic Press, 1982. A two-volume set of remarkable
  games, some analyzed, some unanalyzed, filled with humorous pictures and an oddball type
  of creative wordplay that is Conway's hallmark. Included in Volume 2 are discussions of the
  Cube and Conway's game of Life.
Coueignoux, Philippe. "La reconnaissance des caractcres". La Recherche 12, no. 126
  (October, 1981): 1094-1103. An interesting article about the workings of some computer
  systems that can read text in a variety of typefaces. The ideas derive from work by Blesser
  and colleagues (see above). This work is rather practically oriented, and does not come
  close to giving computers an understanding of letterforms in their full generality.
Csanyi, Vilmos. General Theory of Evolution. Budapest: Akademiai Kiad6, 1982. A thorough
  investigation of the process of simultaneous evolution on both genetic and "memetic"
  fronts.
Davies, Paul. God and the New Physics. New York: Simon \& Schuster, 1983. In this book,
  Davies tackles the biggies of metaphysics: creation, free will, religion, souls, and more.
  Since he is a good scientist as well as a good writer, his musings are articulate and
  penetrating.

Bibliography                                                                               803

-----Other Worlds New York Simon \& Schuster, 1980 A popular account , by a highly
  reliable professional, of the mysteries at the base of quantum mechanics.
Davis, Morton D. Game Theory: A Nontechnical Introduction. New York: Basic Books,
  1983. An excellent overview of all the main ideas of game theory, including many
  unresolved issues such as the Prisoner's Dilemma.
Dawkins, Richard. The Se\#ish Gene. New York: Oxford University Press, 1976. A book that
  views organisms as by-products of a purely molecular-level competition for efficient self-
  replication. A topsy-turvy, disorienting, yet powerfully revealing viewpoint.
DeLong, Howard. A Profile of Mathematical Logic. Reading, Mass.: Addison-Wesley, 1970.
  A wonderful book on the philosophical and technical issues of logic by someone who
  knows how to achieve an artistic balance between formalisms and ideas that appeal to the
  intuition.
Dennett, Daniel C. Brainstorms: Philosophic Essays on Mind and Psychology. Cambridge,
  Mass.: Bradford Books, MIT Press, 1978. A collection of penetrating analyses of problems
  of mind, brain, and computer models of thought, perception, and sensation. Excellent
  rebuttals of such figures as B. F. Skinner and J. R. Lucas, and a wonderful "dessert" at the
  end. Dennett's graceful style and comparative lack of in-references and jargon make this
  book much more engaging than the average philosophy book.
-----"Can Machines Think?" Unpublished manuscript. Afresh look at the power of the Turing
  Test, completely in sympathy with my view that the test as originally posed is, as valid as
  ever, despite the doubts of many.
-----Cognitive Wheels: The Frame Problem of Artificial Intelligence". In Minds, Machines,
  and Evolution. Edited by C. Hookway. New York: Cambridge University Press, 1985. An
  article about why artificial intelligence is so far from achieving common sense.
-----Elbow Room: The varieties of Free Will Worth Wanting. Cambridge, Mass.: Bradford
  Books, MIT Press, 1984. One provocative metaphor after another, all building up towards
  an image of "self-made selves" that enjoy as much free will as it is reasonable to hanker
  after. Although I resist some of this book's conclusions, I find it the best writing on the
  subject that I know.
-----The Logical Geography of Computational Approaches (A View from the East Pole)".
  Unpublished manuscript. Dennett wittily divides the world of Al approaches into two
  camps: the orthodox one (High Church Computtionalism, centered on the "East Pole"
  (located at MIT), and the unorthodox one ("New Connectionism"), scattered hither and yon.
-----The Myth of the Computer: An Exchange". New York Review of Books. (June 24 1982):
  a civil reply to John Searle´s biting review of The Mind's I.
-----The Self as a Center of Narrative Gravity". In Self and Consciousness, edited by P. M.
  Cole et al. New York: Praeger, 1985. A wonderful new metaphor for thinking about
  abstractions such as "I".
Dewdney, A. K. "Computer Recreations: A computational garden sprouting anagrams,
  pangrams, and few weeds Scientific American 251, no. 4 (October 1984): 20-27. Includes a
  description of Lee Sallows' pangram machine, concluding with a public challenge to
  discover a computer-generated pangram.
DeWitt, Bryce S. and Neill Graham, eds. The Many- Worlds Interpretation of Quantum
  Mechanics Princeton NJ. Princton University Press 1973. A well-rounded presentation of
  one of the most disorienting yet irrefutable ways of thinking about reality yet invented.
Dyer, Michael. In-Depth Understanding. Cambridge, Mass.: MIT Press, 1983. A book
  summarizing a Ph.D. project in language understanding that puts together many recent Al
  ideas about how memory is organized and how flexible control structures fit into the
  picture,,
Dylan, Bob. Tarantula. New York: Macmillan, 1971. The poet-singer lets loose some stream-
  of-consciousness musings that are sometimes reasonably intelligible and sometimes totally
  wacko.
Edson, Russell. The Clam Theater. Middletown, Conn.: Wesleyan University Press, 1973
  Strange and surrealistically written fantasies permeated by a tragic vision of life. To


Bibliography                                                                             804

  me, the most haunting passage is a description of a human head: "this teetering bulb of
  dread and dream ...".
Eidswick, Jack. "How to Solve the n x n X n Cube". Mathematics and Statistics Department,
  University of Nebraska, Lincoln, 1982. A useful manual for those beset by generalized
  cubic frustration.
Endl, Kurt. Rubik's Cube Made Simple; The Pyramid; Pyraminx Cube; Impossiball;
  Megaminx; Rubik's Master Cube. Giessen, Germany: WUrfel-Verlag GmbH, 1982. These
  booklets, together with the previous entry, will give you a good reference shelf on
  Cubology.
Erman, Lee D. et al. "The Hearsay-II Speech-Understanding System: Integrating Knowledge
  to Resolve Uncertainty". ACM Computing Surveys 2, no. 2 (June 1980): 213-53. A review
  article describing, with a good deal of hindsight, the architecture and performance of what I
  consider to be the most inspiring work yet done in artificial intelligence.
Evans, Thomas G. "A Program for the Solution of a Class of Geometric-Analogy Intelligence
  Test Questions". In Semantic Information Processing, edited by Marvin
Minsky. Cambridge, Mass.: MIT Press (1968): 271-353. Perhaps the first truly large system
  ever written in Lisp, this impressive project is often pointed to as having "solved" the
  problem of analogies. How far from the truth that is!
Ewing, John and Czes Koiniowski. Puzzle It Out: Cubes, Groups, and Puzzles. New York:
  Cambridge University Press, 1982. Another lively and mathematically interesting book
  about the Cube and its variants.
Fahlman, Scott E., Geoffrey Hinton, and Terrence J. Sejnowski. "Massively Parallel
  Architectures for Al: NETL, Thistle, and Boltzmann Machines". In "Proceedings of the
  AAAI-83 Conference", Washington, D. C., August, 1983. A comparison of some recent
  ideas for AI architectures, including ones inspired by physics, in which statistical
  emergence plays a central role.
Falletta, Nicholas. The Paradoxicon. New York: Doubleday, 1983. An excellent collection of
  paradoxes of all sorts, ranging from Zeno and Epimenides to modern-day voting paradoxes,
  optical illusions, and antinomies in the philosophies of science and mathematics.
Fauconnier, Gilles. Espaces Mentaux: Aspects de la construction du seas dans les langues
  naturelles. Paris: Les Editions de Minuit, 1984. A probing inquiry into how we make sense
  of frame-crossing and counterfactual statements of this sort: "If Clark Gable had been a
  woman, Scarlett O'Hara would have been a man". Full of ideas about reference, identity,
  slippability, and essence.
Feigenbaum, Mitchell. "Universal Behavior in Nonlinear Systems". Los Alamos Science 1,
  no. 1 (Summer 1981): 4-27. An excellent introduction to the studies of chaos that come
  from iterating simple functions. Excellent graphics, some of which are reproduced in
  Chapter 16.
Feldman, Jerome and Dana Ballard. "Connectionist Models and Their Properties". Cognitive
  Science 6, no. 3 (July-September 1982): 205-54. One of several intriguing approaches to
  computational collectivism currently being investigated in cognitive science. Like most
  such models, this one is strongly influenced by ideas about human perception.
Feynman, Richard P. The Character of Physical Law. Cambridge, Mass.: MIT Press, 1967.
  The transcripts of five marvelous lectures delivered at Cornell University in 1964 by a
  physicist who is not only sharp as a nail but also sparklingly witty.
Feynman, Richard P., Robert B. Leighton, and Matthew Sands. The Feynman Lectures in
  Physics. Reading, Mass.: Addison-Wesley, 1965. Lectures given to beginning students,
  often more suitable for advanced students, filled with the excitement of science and the
  peppery observations of Feynman.
Franke, H. W. Computer Graphics-Computer Art. New York: Phaidon, 1971. By far the best
  collection of examples and discussion of computers and art I have yet seen, even though it
  is quite old by now.
Frey, Alexander H., Jr. and David Singmaster. Handbook of Cubik Math. Hillside, NJ.:
  Enslow, 1982. The group-theoretical ideas behind the cube are developed systematically,
  for possible use in an advanced high-school or college-level course.

Bibliography                                                                              805

Friedman, Daniel P. The Little LISPer. Chicago: Science Research Associates, 1974. an
  engaging and elegant introduction to the recursive ideas of Lisp, by one of its most ardent
  and accomplished practitioners.
Fromkin, Victoria A, ed. Errors in Linguistic Performance: Slips of the Tongue, Ear, Pen, and
  Hand. New York: Academic Press, 1980. A compendium of articles on how we can infer
  general mechanisms of thought from observing everyday surface-level quirks that show up
  ubiquitously in speech, typing, handwriting, listening, and so on.
Frutiger, Adrian. Type Sign Symbol. Zurich: Editions ABC, 1980. This book by the creator of
  Frutiger and Univers, two of today's most popular and elegant sans-serif typefaces, deals
  with the interaction of technical constraints and artistic creation, and features many
  beautiful letterforms and styles.
Gablik, Suzi. Progress in Art. New York: Rizzoli International, 1976. A heretical theory
  suggesting that there is a reason why art is where it is today, and that art follows a
  comprehensible trajectory in some abstract space, even if it is explicable only after the fact.
Gardner, Martin. Fads and Fallacies. New York: Dover, 1952. A book that comes about as
  close as one could hope to being a course on common sense. Many chapters of this classic
  bunk-puncturer should be part of the cultural education of everyone.
Gardner, Martin. "Mathematical Games: White and brown music, fractal curves, and one-
  over-f fluctuations". Scientific American 238, no. 4 (April 1978): 16-32. A solid
  introduction to the concepts of fractals and multi-level statistical structures.
Gardner, Martin. Science: Good, Bad, and Bogus. Buffalo: Prometheus Press, 1981.
  Following in the footsteps of Fads and Fallacies, this book sharp-swordedly decapitates one
  dragon of hokum after another-and yet sadly, like a hydra's many heads, they always seem
  to pop back up again.
Gardner, Martin. Wheels, Life, and Other Mathematical Amusements. New York: W. H.
  Freeman, 1983. Gardner's final collection of his sparkling Scientific American columns,
  showing why mathematics is the ultimate metamagic.
Gebstadter, Egbert B. Thetamagical Memas: Seeking the Whence of Letter and Spirit. Perth:
  Acidic Books, 1985. A curious pot-pourri, bloated and muddled-yet remarkably similar to
  the present work. This is a collection of Gebstadter's monthly rows in Literary Australian
  together with a few other articles, all with prescripts. Gebstadter is well known for his love
  of twisty analogies, such as this one (unfortunately not found in his book): "Egbert
  Gebstadter is the Egbert Gebstadter of indirect self-reference."
Geisel, T. and J. Nierwetberg. "Universal Fine Structure of the Chaotic Region in Period-
  Doubling Systems". Physical Review Letters 47, no. 14 (October 5, 1981): 975-78. An
  investigation turning up remarkable order deep in the heart of chaos.
Gentner, Dedre. "Structure-Mapping: A Theoretical Framework for Analogy". Cognitive
  Science 7 (1983): 155-70. One in a series of papers attempting to draw guidelines that will
  help distinguish good analogies from bad.
Gödel, Kurt. On Formally Undecidable Propositions in "Principia Mathematica" and Related
  Systems. Translated by Bernard Meltzer, edited and with an introduction by R. B.
  Braithwaite. New York: Basic Books, 1962. The fundamental work that opened up worlds
  to logicians, philosophers, computer scientists, and others.
Golden, Michael. "Don't Rewrite the Bible". Newsweek (November 7,4983):47. This piece
  sounds almost like my "Person Paper" in the way it exploits crude mockery of progressive
  new usages in order to make oppressive old usages sound good. The difference is that
  Golden is not being sarcastic.
Golomb, Solomon. "Rubik's Cube and Quarks". American Scientist 70, no. 3 (May June
  1982): 257-59. In which Golomb puts forth his analogy between twisted corners on the
  Cube and fractionally charged subatomic particles.
Gonick, Larry and Mark Wheelis. A Cartoon Guide to Genetics. New York: Barnes \& Noble,
  1983. A short and snappy-but most informative and entertaining-overview of molecular
  biology, genetics, and their human import.
Gould, Stephen Jay. Hen's Teeth and Horse's Toes. New York: W. W. Norton, 1983. Written
  in a lively and engaging way, this book (like its predecessors, Ever Since Darwin and

Bibliography                                                                                806

 The Panda's Thumb) relates with great clarity the issues of evolution as illustrated in
 nature's unbelievable variety of quirks.
Gray, J. Patrick and Linda Wolfe. "The Loving Parent Meets the Selfish Gene". Inquiry 23:
 233-42. A cogent set of arguments against the pessimistic theses of Louis Pascal, to which
 he then replies (see below).
Grossman, I. and W. Magnus. Groups and Their Graphs. New York: Random House New
 Mathematical Library, 1975. An excellent visual introduction to group theory, featuring so-
 called "Cayley diagrams", which lay bare the structure of a group in a marvelously clear
 way.
Ground Zero. Nuclear War: What's in It for You? New York: Pocket Books, 1982. A
 dispassionate, calmly reasoned discussion for "just-plain folk" of what it means for
 countries to make and threaten each other with nuclear weapons. The best primer that I have
 come across on the subject, and one that I wish everyone would read and take to heart.
Guillemin, Victor. The Story of Quantum Mechanics. New York: Scribner's, 1968. A
 thoughtful survey of where quantum mechanics came from and what it has revealed about
 the physical world, concluding with a long discussion of the philosophical implications of
 quantum mechanics for such things as free will and causality.
Haab, A,, A. Stocker, and W. Hattenschweiler. Lettera 1 and Lettera 2. Arthur Niggli, 1954
 and 1961. Some deliciously fanciful alphabets are featured in these books. Books like these
 ought to make anyone marvel at the fluidity of letters and human perception.
Hachtman, Tom. Double Takes. New York: Harmony Books, 1984. Two-in-one caricatures,
 done in a very satisfying way. It makes one wonder: could a computer ever do this kind of
 thing?
Hansel, C. E. M. ESP and Parapsychology: A Critical Re-Evaluation. Buffalo: Prometheus
 Books, 1980. A scholar looks at the claims of parapsychologists and finds them wanting.
 Hanson, Norwood Russell. Patterns of Discovery. New York: Cambridge University Press,
 1969. A philosophically rich discussion of how physical concepts progressed by fits and
 starts through the centuries, culminating in the strange world of elementary particles that we
 are now finding are not so elementary.
Hardin, Garrett. "The Tragedy of the Commons". Science 162, no. 3859 (December 13,
 1968): 1243-48. Using the metaphor of shared grazing land that gets ruined by overuse
 without any specific person being responsible, Hardin argues with great force that people's
 narrow-minded selfishness will drive us to extinction unless we band together and form
 strong organizations dedicated to global ecological goals, particularly population control.
Harel, David. "Response to Scherlis and Wolper". Communications of the ACM 23, no. 12
 (December 1980): 736-37. An amusing list of tricky self-references known to Harel,
 sporting a footnote claiming to be the first published example of "a reference to the very
 self-reference being discussed in the very letter being read!"
Hart, H. L. A. "Self-Referring Laws". In Festskrift Tillagnad Karl Olivecrona. Stockholm:
 Kungliga Boktryckeriet, P. A. Norstedt och Soner, 1964. A reply to Alf Ross (see below), in
 which Hart, perhaps today's leading philosopher of law, claims that self-amending laws are
 legally possible.
Harth, Erich. Windows on the Mind: Reflections' on the Physical Basis of Consciousness.
 New York: William Morrow, 1982. Lively discussions of brain, mind, consciousness, and
 how they all tie together, by a physicist with a philosophical bent.
Haugeland, John, ed. Mind Design: Philosophy, Psychology, and Artificial Intelligence.
 Cambridge, Mass.: Bradford Books, MIT Press, 1981. A self-proclaimed sequel to Alan
 Ross Anderson's Minds and Machines, now that artificial intelligence has had a couple of
 decades to develop. This volume is of uneven quality, but it contains many articles sure to
 provoke thought.
Heiser, Jon F. et al. "Can Psychiatrists Distinguish a Computer Program Simulation of
 Paranoia from the Real Thing? The Limitations of Turing-Like Tests as Measures of the
 Adequacy of Simulations". Journal of Psychiatric Research 15, no. 3 (1979): 149-62.
 Researchers who developed the program called "Parry" claim that the Turing Test's validity
 is cast in doubt because Parry passed a pseudo-Turing Test.

Bibliography                                                                              807

Hinton, Geoffrey and James Anderson, eds. Parallel Models of Associative Memory.
  Hillsdale, N .J.: Lawrence Erlbaum Associates, 1981. A collection of articles exploring
  different models, all of which are based on the thesis that human perception and memory
  are statistically emergent phenomena and are best modeled as such.
Hintze, Wolfgang. Der Ungarische Zauberwiirfel. Berlin: VEB Deutscher Verlag der
  Wissenschaften, 1982. This is an excellent book about the Cube, including new
  mathematical results on subgroups. Serious cubists who can read German should obtain this
  book.
Hobby, John, and Gu Guoan. "A Chinese Meta-Font". Stanford, Calif.: Stanford University
  Computer Science Department Technical Report STAN-CS-83-974, 1983. A description of
  a system for creating Chinese characters whose style is "tunable" to some extent. This
  system is very similar in some ways to Han Zi, the variable-style Chinese-character-
  generation system developed by David Leake and myself at Indiana, at about the same time.
Hodges, Andrew. Alan Turing: The Enigma. New York: Simon \& Schuster, 1983. The
  definitive biography of this extremely important figure in mathematics and computer
  science, written with empathy, insight, and charm.
Hofstadter, Douglas R. Ambigrams. Forthcoming in 1985 from the Centre d'Art
  Contemporain (Geneva, Switzerland). A collection of ambigrams (or "inversions", as Scott
  Kim prefers to call them) aptly subtitled "A Panoply of Palindromic Pinwheels". The
  foreword is by Scott (to whose book I symmetrically wrote the foreword). Also included are
  a lecture on ambigrams (in French) and an interview (in English).
-----Analogies and Metaphors to Explain Godel's Theorem". Two-Year College Mathematics
  Journal 13, no. 2 (March, 1982): 98-114. A collection of simple images and ideas that help
  to build up one's intuitions to the point where one can fully digest the intricacies of Godel's
  clever construction.
-----The Architecture of Jumbo". In "Proceedings of the Second Machine Learning
  Workshop", Monticello, Illinois, 1983. A description of the architecture of a biologically-
  inspired Al system that exploits parallelism and randomness in order to build up "well-
  chunked wholes" out of isolated parts and then to allow them to seek maximal "happiness"
  by internally reconfiguring themselves in a manner governed by a computational
  temperature.
-----The Copycat Project: An Experiment in Nondeterminism and Creative Analogies".
  Cambridge, Mass.: MIT Artificial Intelligence Laboratory Al Memo 755, April 1984. A
  description of ongoing research in my group, first at Indiana, then at MIT, and now at
  Michigan. The goal is to make a system capable of doing analogies in a tiny domain, but
  with the insight (and short sight) of humans. Describing how the jumbo architecture can be
  adapted to this task is the burden of this paper. .
-----512 Words on Recursion". Math Bulletin, Bronx High School of Science (1984): 18-19.
  A recursively structured article, which quotes a compressed version of itself (in which, of
  course, a more compressed version is quoted, etc.).
-----Gödel, Escher, Bach: an Eternal Golden Braid. New York: Basic Books, 1979. Originally
  titled "Gödel’s Theorem and the Human Brain", this book tries to build up the concepts
  necessary to show how ideas that arise in Gödel’s proof will one day be at the core of
  explanations of consciousness and the meaning of the word "I". Analogy, humor, and
  contrapuntal dialogues are essential features of GEB.
-----Gridfonts". Unpublished manuscript, 1984. A still-growing collection of (currently) about
  300 skeletal typefaces-that is, distinctive and artistically consistent ways of rendering all 26
  letters in a highly confined grid. The purpose is to provide material for discussion of the
  roots of style and the fluid nature of categories.
-----On Seeking Whence". Unpublished manuscript, 1982. A discussion of the immense
  difficulty of extrapolating linear patterns, and its relationship to such deeply human
  experiences as scientific induction, analogy-making, understanding music, and creativity.
-----Poland: A Quest for Personal Meaning". Poland (May, 1981): 42-47. A compression of a
  longer piece called "Poland: A Mythical Quest", describing my powerful emotional
  reactions to Poland on my first visit there, in 1975.

Bibliography                                                                                 808

-----In Search of Essence". Unpublished manuscript, 1983. A discussion of what the "essence"
  of a written passage is, by way of a description of problems encountered in translating the
  dialogue Controcrostipunctus (from C,6, Escher, Bach) into French.
-----Simple and Not-So-Simple Analogies in the Copycat Domain". Unpublished manuscript.
  A collection of 84 analogy problems in the Copycat domain that show how diverse the tiny
  alphabetic universe can be.
Hofstadter, Douglas R., Gray Clossman, and Marsha Meredith. "Shakespeare's Plays Weren't
  Written by Him, but by Someone Else of the Same Name". Bloomington: Indiana
  University Computer Science Department Technical Report 96, 1980. "A study of
  intensionality in frame-based representation systems" is the subtitle. The paper's main
  purpose is to discuss the relationship between slots that an entity fills, and that entity's
  identity. The notion of "Core ID" is presented and explored in an analysis of the title
  sentence. Closely related to the problems discussed by Fauconnier (see above).
Hofstadter, Douglas R. and Daniel C. Dennett, eds. The Mind's : Fantasies and Reflections on
  Self and Soul. New York: Basic Books, 1981. Our purpose was to jolt people on all sides of
  the fence in this anthology of fictional and evocative pieces on the curious fact (or illusion)
  that something we call an "I" is somehow connected to some hunk of matter floating
  somewhere and somewhen in some universe.
Holland, John et al. Induction: Processes of Inference, Learning, and Discovery. Forthcoming.
  A wide-ranging study of how learning takes place in genuine human minds, and how
  aspects of it might be modeled in computer simulations. This pioneering work is by a
  computer scientist, two' psychologists, and a philosopher, and its breadth reflects theirr
  diversity.
Hollis, Martin and Steven Lukes, eds. Rationality and Relativism. Cambridge, Mass.: MIT
  Press, 1982. Showing how scholars can get mired in the endlessly circular reasoning that
  arises when any system of belief tries to provide its own justification. Reminiscent of
  "MUnchhausens Zopf": the quandary of Baron MUnchhausen trying to pull himself out of a
  quicksand quagmire simply by tugging on his own braid.
Hopfield, J. J. "Neural networks and physical systems with emergent computational abilities".
  In Proceedings of the National Academy of Sciences 79, Washington, D.C., 1982: 2554-58.
  A system that makes use of statistics and parallelism to model familiar properties of the
  mind: reconstructing a full memory from a partial one, generalizing, and so on.
How to Learn Lettering. Hong Kong: Nam San Publisher, n.d. A collection of Chinese
  characters in modern artistic styles, mostly for the use of advertisers, but also a treasure
  trove for those interested in the elusive "sameness" we see in wildly different renderings of
  a single Platonic essence.
Huff, William. "A Catalogue of Parquet Deformations". School of Architecture, State
  University of New York at Buffalo. Approximately 35 parquet deformations created by
  students in William Hull's studio, assembled and commented on by Huff.
Hughes, Patrick and George Brecht. Vicious Circles and Infinity: An Anthology of
  Paradoxes. New York: Penguin, 1975. A' delightful and provocative collection of
  paradoxical material and choice epigrams embodying paradoxes of every conceivable sort.
  Includes lengthy discussions of the paradox of the Unexpected Examination, also known as
  the "Hangman's Paradox".
Huneker, James Gibbons. Chopin: The Man and His Music. New York: Dover, 1966.
  Originally published around the turn of the century, this romantic biography features florid
  prose so rich that it verges on meaninglessness, and yet it is wonderfully evocative.
Jaspert, W. Pincus, W. Turner Berry, and A. F. Johnson. The Encyclopaedia of Type Faces.
  Poole, England: Blandford Press, 1983. An engrossing catalogue of over 1,000 typefaces.
  One of its unique features is its thorough crediting of designers: the people whose
  exquisitely developed sense of line and space is far too often completely taken for granted.
  An index to designers and an index to typefaces are included. This way you can get a sense
  for the stylistic range of various designers.
Johnson-L'Laird, P. N. and P. C. Wason, eds. Thinking. New' York: Cambridge University
  Press, 1975. A diverse collection of pieces by top-notch authors about nearly

Bibliography                                                                                809

  all aspects of human thought, including imagery, perception, inference, categories, memory,
  language, and so on.
Kadanoff, Leo. "Roads to Chaos". Physics Today 36, no. 12 (December 1983): 46-53. A good
  summary of the connection between mathematical approaches to chaos and the physical
  phenomena to be explained.
Kamack, H. J. and T. R. Keane. "The Rubik Tesseract". Unpublished manuscript, 1982. A
  companion piece to the one by Eidswick (see above), this discussion of a 3 x 3 x 3 x 3
  "Rubik's hypercube" is a tour de force in visualization-for its readers almost as much as for
  its authors!
Kahneman, Daniel and Amos Tversky. "The Simulation Heuristic". In Judgment Under
  Uncertainty: Heuristics and Biases, edited by Daniel Kahneman, P. Slovic, and A. Tversky.
  New York: Cambridge University Press, 1982: 202-8. In what ways are people inclined to
  let situations mentally glide into alternate versions of themselves, and how do subtle
  cognitive pressures modify those tendencies? Two outstanding cognitive psychologists here
  describe their findings about this sort of slippability or "alternity" in their subjects' minds.
Kanerva, Pentti. Self-Propagating Search: A Unified Theory of Memory. Stanford, Calif.:
  Center for the Study of Language and Information, Technical Report, Stanford
University, 1984. A neurocomputational theory of memory: How to make similarity-sensitive,
  reconstructive software out of statistically robust addressing hardware. This may sound a bit
  abstruse, but in my opinion Kanerva's elegant work constitutes one of the most important
  steps toward reconciliation of fluidity and mechanism-mind and brain -that I have seen
  taken.
Kennedy, Paul E. Modern Display Alphabets. New York: Dover, 1974. An excellent
  collection of visually attractive letterforms in styles of all sorts, and exhibiting all degrees of
  wildness.
Kim, Scott E. "The Impossible Skew Quadrilateral: A' Four-Dimensional Optical Illusion". In
  Proceedings of the 1978 AAAS Symposium on Hypergraphics: Visualizing Complex
  Relationships in Art and Science. Boulder, Colo.: Westview Press, 1978. Like the work by
  Kamack and Keane (see above), this article takes a powerful visual imagination to write or
  to read. The familiar "impossible triangle" is here carried one step further, by a bold process
  of analogy. The writing style is full of typical Kimian recursive trickery and parallel
  passages.
-----Inversions. Peterborough, N.H.: Byte Books, 1981. A collection of inversions (or
  "ambigrams", as I prefer to call them) aptly subtitled "A Catalogue of Calligraphic
  Cartwheels". The foreword is by myself (to whose book Scott symmetrically wrote the
  foreword). The prose is as full of illusions and tricks of parallelism as the drawings are,
  although it is not immediately apparent.
-----Noneuclidean Harmony". In The Mathematical Gardner, edited by David A. Klarner.
  Belmont, Calif.: Wadsworth International/Prindle, Weber, and Schmidt, 1981. A serious
  treatise on atonal geometry masquerading as a piece of humor. Describes Georg Cantor's
  important results on supersonic pitches, 2 against 3 correspondence, and uncountable
  rhythms.
-----Visual Art: The Creative Cycle". Response, no. I (December, 1981). Minnesota Artists
  Exhibition Program. A short article putting forth a thesis about creative activity, at the heart
  of which is a loop similar to the one I describe in the P.S. to Chapter 12 of this book.
Kirkpatrick, S., C. D. Gelatt, Jr., and M. P. Vecchi. "Optimization by Simulated Annealing".
  Science 220, no. 4598 (May 13, 1983): 671-80. An article describing a new statistical
  technique for seeking the optimal state of a system with a huge number of degrees of
  freedom, based on local improvement modulated by occasional local degradation, the
  amount of which is determined by a "cooling schedule", analogous to the annealing of a
  metal to strengthen it.
Kleppner, Daniel, Michael G. Littman, and Myron L. Zimmerman. "Highly Excited Atoms".
  Scientific American 244, no. 5 (May 1981): 130-49. A bridge between microscopic
  quantum-mechanical phenomena and macroscopic classical phenomena is provided by the
  study of "Rydberg atoms", with which this paper is concerned.

Bibliography                                                                                    810

Knuth, Donald. "The Concept of a Meta-Font". Visible Language 16, no. 1 (Winter 1982): 3-
  27. Knuth describes his program, METAFONT, for creating an entire family of typefaces at
  one fell swoop, by parametrizing all letters of the alphabet with a consistent set of
  parameters. This is the paper that triggered my response printed as Chapter 13.
-----TEX and Metafont: New Directions in Typesetting. Bedford, Mass.: Digital Press, 1979.
  Knuth describes in detail how to use his programs that facilitate typesetting and letter
  design.
Kobylariska, Krystyna. Chopin in His Own Land. Cracow: Polish Music Publishers, 1956. A
  collection of hard-to-find Chopiniana, beautifully reproduced in large format.
Koestler, Arthur. The Act of Creation. New York: Dell, 1964. The great novelist and amateur
  psychologist delves into his own mind to come up with interesting but often wrong-headed
  speculations on humor, creativity, and insight.
-----The Roots of Coincidence: An Excursion into Parapsychology. New York: Vintage,
  1972. Koestler reveals his belief in the occult, somewhat reminiscent of Arthur Conan
  Doyle's belief in fairies.
Kolata, Gina. "Does Godel's Theorem Matter to Mathematics?" Science 218 (November 19,
  1982): 779-80. An excellent description of how the specification of what today would be
  considered "large integers" depends on abstruse concepts from mathematical logic.
Koning, H. and J. Eizenberg. "The language of the prairie: Frank Lloyd Wright's prairie
  houses". Environment and Planning B, 8 (1981): 295-323. A description of how pseudo-
  Frank-Lloyd-Wright houses can be generated from a "shape grammar". Kripke, Saul.
  Naming and Necessity. Cambridge, Mass.: Harvard University Press, 1972. This theory of
  extensions, intensions, "rigid designators", and identity is. in direct opposition to my views
  (put forward in the "Shakespeare" paper cited above) that the roots of identity come from
  patterns of embeddedness in a network.
Kuwayama, Yasaburo. Trademarks and Symbols, Volume 1: Alphabetical Designs. New
  York: Van Nostrand Reinhold, 1973. A celebration of the craziness of letterforms, this
  collection contains letters to boggle anyone's mind.
Larcher, Jean. Fantastic Alphabets. New York: Dover, 1976. The letterforms in Kuwayama's
  book are highly imaginative, but by comparison, the letterforms in Larcher's book are
  downright zany. It just goes to show that a computer program that could deal with letters in
  their full complexity would be able to deal with the entire world.
Lehninger, Albert. Biochemistry, 2d. ed. New York: Worth, 1975. A well-written treatise
  covering the huge field of biochemistry in a lucid manner.
Lennon, John. In His Own Write and A Spaniard in the Works. New York: Simon \&
  Schuster, 1964 and 1965. This Beatle had a wonderful sense of silliness, all his own.
  Stories, pictures, poems-even a little drama or two. Perhaps better than the Beatles' music, if
  I may venture a heretical opinion.
Letraset, Inc. Graphic Art Materials Reference Manual. Paramus, NJ.: Letraset, Inc., 1981.
  This is the best inexpensive typeface manual available, and is well worth the price. It
  contains nearly all the common typefaces, as well as many unusual ones. This is a great way
  to get into typefaces.
Letraset, Inc. Letraset Greek Series. Athens: A. Pallis, 1984. This to me is an amazing
  catalogue of typefaces, because each and every one of them represents a "transalphabetic
  leap". Here we see such faces as Blanchard, Futura Black, Helvetica, Korinna, Optima,
  Souvenir, University Roman, Zipper, and a number of others-all in Greek! All I can say is,
  "Wow!"
Levin, Michael. "Mathematical Logic for Computer Scientists". Cambridge, Mass.: MIT
  Laboratory for Computer Science Technical Report LCS TR 131, June, 1974. An
  unorthodox treatment of logic for those more comfortable with Lisp than with number
  theory.
Lipman, Jean and Richard Marshall. Art about Art. New York: E. P. Dutton, 1978. An
  amusing annotated collection of self-conscious art, including pieces by Roy Lichtenstein,
  Andy Warhol, Robert Rauschenberg, Jasper Johns, Robert Arneson, Tom Wesselmann,
  Larry Rivers, Mel Ramos, Peter Saul, and many others.

Bibliography                                                                                811

Loeb, Marcia. New Art Deco Alphabets. New York: Dover, 1975. Showing how somebody
  can-perfectly recreate the spirit of an era in many different ways. These original alphabets
  all look like they came straight out of the 1930's. Compare with the books by Bush (see
  above), and McCall (see below).
Lucas, J. R. "Minds, Machines, and Gbdel". Reprinted in Minds and Machines, edited by
  Alan Ross Anderson. Englewood Cliffs, NJ.: Prentice-Hall, 1964. The article that launched
  a thousand rebuttals. Although I find its conclusions totally unjustified, the issues are well
  raised and deserve to be thought about far more. Closely related to the issues discussed in
  the book by Webb (see below).
Machlup, Fritz and Una Mansfield, eds. The Study of Information: Interdisciplinary
  Messages. New York: Wiley-Interscience, 1983. A multi-faceted collection of articles by
  high-level scholars about information theory, cybernetics, artificial intelligence, cognitive
  science, librarianship, and more. Many of the pieces are statements of opinion, and conflict
  with one another. This volume includes the exchange between Allen Newell and myself,
  which is continued in the P.S. to my Chapter 26.
Mandelbrot, Benoit. The Fractal Geometry of Nature. New York: W. H. Freeman, 1982. This
  latest presentation of Mandelbrot's visions is rich in imagery and ideas. Fractal shapes
  constitute one of the twentieth century's most fertile mathematical playgrounds.
Marek, George R. and Maria Gordon-Smith. Chopin. New York: Harper \& Row, 1978. One
  of the many biographies of Chopin in English. I find it to be balanced and well written.
Marx, George, Eva GajzSgo, and Peter Gnadig. "The universe of Rubik's cube". European
  journal of Physics 3 (1982): 34-43. The implications of extending the concept of entropy to
  the surface of the Cube are reported on, with the results of several statistical studies done on
  computer.
Max, Nelson. Space-Filling Curves. Chicago: International Film Bureau, Inc., 1974. A
  marvelous movie done with computer graphics showing the fascination of some of the
  earliest-discovered fractal curves. The computer-produced music adds extra appeal and
  eerieness.
May, Robert. "Simple Mathematical Models with Very Complicated Dynamics". Nature 261,
  no. 5560 (June 10, 1976): 459-67. An early review article about the chaos one is led to
  when one iterates simple functions on the interval [0,1 ]. Very informative, except for the
  reversal of two figures!
Maynard-Smith, John. Evolution and the Theory of Games. New York: Cambridge University
  Press, 1982. Showing how new depths of understanding of the mechanisms of evolution can
  come from modeling the interactions of organisms in a common environment in terms of
  mathematical game theory.
McCall, Bruce. Zany Afternoons. New York: Alfred A. Knopf, 1982. A genius for capturing
  the flavor of things turns all his talents to recreating the twenties, thirties, forties, and fifties-
  not only in America but also abroad. This very funny book is an amazing achievement.
  Compare with the books by Bush and Loeb (see above).
McCarthy, John. "Attributing Mental Qualities to Machines". In Philosophical Perspectives
  on Artificial Intelligence, edited by Martin Ringle, 161-95. Atlantic Highlands, N. J.:
  Humanities Press, 1979. Under what circumstances does it make sense to speak of machines
  having desires or beliefs? When to adopt this "intentional stance" is here discussed by one
  of the pioneers of artificial intelligence.
-----History of Lisp". In History of Programming Languages, edited by Richard L.
  Wexelblatt, 217-23. New York: Academic Press, 1980. Lisp's genesis as recalled by its
  inventor.
McCarty, L. T. and N. S. Sridharan. "A Computational Theory of Legal Argument". New
  Brunswick, NJ.: Rutgers University Computer Science Department Technical Report LPR-
  TR-13, January 1982. A professor of law and a computer scientist work together to
  implement "prototype deformation", which they consider the key to computer modeling of
  argumentation by precedent.
McClelland, James L., David E. Rumelhart, and Geoffrey E. Hinton, eds. Interactive
  Activation: A Framework for Information Processing. Forthcoming. A collection of articles

Bibliography                                                                                       812

  on statistically emergent parallel computation featuring "temperature" as a regulator of
  randomness.
McKay, Michael D. and Michael S. Waterman. "Self-Descriptive Strings". Mathematical
  Gazette 66, no. 435 (1982): 1-4. A mathematical study of a simple class of sentences that
  attempt to inventory themselves.
Meehan, James R. "Tale-Spin, an Interactive Program that Writes Stories". In Proceedings of
  the 5th International Conference on Artificial Intelligence-1977, Vol. 1: 91-98. Pittsburgh:
  Department of Computer Science, Carnegie-Mellon University. One of the funniest (and
  therefore most sensible) articles ever written on computers and their "understanding" of
  language. Meehan explains how, in one "mis-spun tale", his program made gravity fall into
  a river and drown. Meehan's discussions of mis-spun tales in general are among the most
  illuminating passages I know on that perennial vexation: the fact that AI programs
  stubbornly resist acquiring common sense, no matter how hard or how often they are
  spanked.
Miller, Casey and Kate Swift. The Handbook of Nonsexist Writing for Writers, Editors, and
  Speakers. New York: Barnes and Noble, 1980. A first-class book. It is a scandal that this
  book is not found on every newsman's desk. Thinking about these issues is not only socially
  important, but challenging and fascinating.
Words and Women. New York: Doubleday, Anchor Press, 1977. A magnificent and
  devastating job of showing the absurdity of thinking that women have been dealt a fair hand
  by our society. If language is the thermometer of society, then this book reveals that
  collectively, we are gravely ill.
Mills, George. "Gddel's Theorem and the Existence of Large Numbers". Northfield,
  Minnesota: Mathematics Department, Carleton College, October 1981. Showing how the
  descriptions of some stupendously large numbers can be obtained via recursive definitions
  of functions and the process of diagonalization. Repeated jootsing leads to some remarkable
  destinations!
Minsky, Marvin. "A Framework for Representing Knowledge". In The Psychology of
  Computer Vision, edited by P. H. Winston. New York: McGraw-Hill, 1975. The paper that
  defined such now-central notions to Al as frames, slots, and default assumptions.
-----Matter, Mind, and Models". In Semantic Information Processing, edited by Marvin
  Minsky. Cambridge, Mass.: MIT Press, 1968. Some cogent speculations about the meaning
  of the word "I" in computational terms.
-----Why People Think Computers Can't". AI Magazine (Fall 1982): 3-15. With his usual
  unusual insight, Minsky tears people apart for not understanding how to think about
  thinking (or about machines, for that matter).
Mondrian, Piet. Tout l'tEuvre Peint de Piet Mondrian. Paris: Flammarion, 1976. For a grand
  overview of the evolution of the style of one painter, I know of no better book than this,
  which traces Mondrian from his earliest representational paintings to his most abstract and
  geometrical ones, revealing the sweep to be continuous and logical, but no less dramatic for
  that.
Morrison, Philip and Phylis, and the Office of Charles and Ray Eames. Powers of Ten. New
  York: Scientific American Books, 1983. Kees Boeke's inspirational Cosmic View (see
  above) revisited some thirty years later, with considerably more commentary. A charming
  and . - worthy successor.
Myhill, John. "Some Philosophical Implications of Mathematical Logic: Three Classes of
  Ideas". Review of Metaphysics 6, no. 2 (December 1952): 165-98. In 1982 I met Myhill and
  asked him about this paper. He told me he considered it to be a piece of junk. That
  astonished me, since I consider it to be a thoughtful and important piece of philosophizing
  making use of mathematical metaphors, something that hardly anyone dares to do. You
  never know how someone will evaluate their own work 30 years later!
Nagel, Ernest, and J. R. Newman. Gödel’s Proof. New York: New York University Press,
  1958. A gracious and highly accessible introduction to the twists in Godel's reasoning, as
  well as to the philosophical issues surrounding his work.


Bibliography                                                                             813

Nakanishi, Akira. Writing Systems of the World. Rutland, Vt.: Tuttle, 1980. For a sampling
  of the many different spirits residing in letterforms, see this book. It features reproductions
  of newspaper pages showing each writing system in several different styles.
Newell, Allen. "Physical Symbol Systems". Cognitive Science 4, no. 2 (April-June 1980):
  135-83. A lengthy article putting forth the orthodox dogma on which artificial intelligence
  has traditionally considered itself founded: the idea that a universal computer embodies all
  the prerequisites for intelligent behavior.
Norman, Donald. "Categorization of Action Slips". Psychology Review 88, no. 1 (January
  1981): 1-15. A delightful compendium of types of error that people commit in performing
  everyday actions such as putting water on to boil, answering the telephone, unbuckling
  one's seatbelt, driving home from work, and so on. There is remarkable regularity behind
  the seeming chaos, and Norman's purpose is to chart and exploit that regularity in order to
  reveal hidden mechanisms of thought.
Nozick, Robert. Philosophical Explanations. Cambridge, Mass.: Harvard University, Belknap
  Press, 1981. A very wide-ranging book on philosophy, intended for lay readers as well as
  for professionals. It covers matters from personal identity and the meaning of reference to
  free will and morality, and only occasionally lapses into brief spasms of absolutely
  inscrutable jargon.
Parfit, Derek. Reasons and Persons. Oxford: Clarendon Press, 1984. A very thoughtful
  treatise on moral dilemmas and ethical behavior, rooted in close consideration of the
  deepest roots of caring: why do I care about my present and future self? Parfit knows that to
  make any serious attempt to answer this riddle, one must look long and hard at the meaning
  of the word "I" in the real world and in many counterfactual ones, which he does with skill
  and insight.
Pascal, Louis. "Human Tragedy and Natural Selection". Inquiry 21: 443-60. Taking up where
  Garrett Hardin left off (see above), Pascal paints a gloomy picture of a population explosion
  as the natural outcome of selection itself, something built into the nature of societies just as
  deeply as the sexual drive is built into individuals.
"Rejoinder to Gray and Wolfe". Inquiry 23: 242-51. A bitter indictment of the human race's
  apathy before visibly onrushing catastrophe.
Peattie, Lisa. "Normalizing the Unthinkable". Bulletin of the Atomic Scientists 40, no. 3
  (March 1984): 32-36. Like Pascal, Peattie is concerned with people's apparent ability to turn
  off their sensitivities and to focus on the very local to such an extent that great tragedies are
  allowed to ensue. In a striking analogy, she likens the current public apathy about the arms
  madness to the inhumanity of collaborators in Hitler's concentration camps.
Perec, Georges. La Disparition. Paris: Editions Denoel, 1969. If anything, writing without 'e's
  is harder in French than in English, yet here is an entire novel in that bizarre dialect.
  Naturally enough, its subject is the mysterious disappearance of item number five in a
  collection of twenty-six objects. It was probably inspired by the 'e'-less novel Gadsby,
  written in English by Ernest Vincent Wright in the late 1930's.
Perfect, Christopher and Gordon Rookledge. Rookledge's International Typefinder. New
  York: Frederick C. Beil, 1983. A wonderful (though expensive) compendium of typefaces,
  indexed in such a way that you can look a typeface up by its features, thus allowing you to
  home in quickly on an unknown specimen instead of spending hours leafing through
  catalogues. Lovers of my "horizontal and vertical problems" will delight in this book.
Phillips, Tom. A Humument: A Treated Victorian Novel. New York: Thames \& Hudson,
  1980. Like a child who has covered a wall with colorful crayon drawings, Tom Phillips has
  completely obliterated the pages of an old novel (A Human Document by W. H. Mallock)
  with his colorful scribblings, and only here and there do traces of the original show through.
  A droll stunt!
Poincare, Henri. "On Mathematical Creation". In The World of Mathematics, vol. 4, edited by
  James R. Newman, 2041-50. New York: Simon \& Schuster, 1956..A lecture presented to
  the Psychological Society in Paris in the early part of this century. Anticipating
  developments in cognitive science some eighty years later, Poincare speculates on the
  nature of the events taking place inside his skull as he makes mathematical discoveries.

Bibliography                                                                                  814

Po1ya, George. How to Solve It. New York: Doubleday, 1957. P61ya, like Poincart- a
 mathematician fascinated by thought processes, attempts here to give recipes for how to
 attack mathematical problems. The problem with this is that there is-and can be-no failsafe
 recipe. Even trying to give guidelines is probably futile. The nose that smells the right route
 is simply rare, and there are no two ways about it.
Post, Emil. "Absolutely Unsolvable Problems and Relatively Undecidable Propositions:
 Account of an Anticipation". In The Undecidable, edited by Martin Davis, 338-443. New
 York: Raven, 1965. In this paper, Post concludes that mathematicians' thought processes are
 essentially creative and non-mechanizable. Related to the article by Myhill (see above).
Poundstone, William. The Recursive Universe: Cosmic Complexity and the Limits of
 Scientific Knowledge. New York: William Morrow, 1985. A superlative account of the
 reductionist miracle: fantastically complex entities-living, self-reproducing organisms-turn
 out to be vast arrays of very simply interacting parts. Von Neumann's trick of self-
 reproduction without infinite regress (adapted from Gödel) is one of the main topics
 explored here, and with the help of Conway's absorbing game of Life, which plays the
 starring role in his book, Poundstone does a masterful job of explaining how it comes about.
Racter. The Policeman's Beard Is Half Constructed. New York: Warner Books, 1984.
 "Racter" is a program written by Bill Chamberlain and Thomas Etter. The Policeman's
 Beard is a book written by Racter. It is all somewhat tongue-in-cheek, because Racter does
 not really know much about half-constructed beards, but what is lovely about Racier's prose
 is the way it skirts the fringes of meaning, weaving drunkenly across the boundary between
 sense and senselessness.
Rapoport, Anatol. Two-Person Game Theory. Ann Arbor: University of Michigan Press, Ann
 Arbor Science Library, 1966. A sound treatment of game theory, featuring a most
 .interesting personal discussion on opinions about the meaning of "rationality" in Prisoner's-
 Dilemma-like situations.
Reps, Paul. Zen Flesh, Zen Bones. New York: Doubleday, Anchor Press, n.d. An easily
 available collection of Zen koans, highly amusing and, perhaps, even enlighteningproviding
 you take it all with sufficiently many grains of salt.
Rogers, Hartley. Theory of Recursive Functions and Effective Computability. New York:
 McGraw-Hill, 1967. A standard reference work on many advanced concepts in
 metamathematics, including such concepts as "productive" and "creative" sets, referred to in
 my Chapter 13.
Ross, Alf. "On Self-Reference and a Puzzle in Constitutional Law". Mind 78, no. 309
 (January 1969): 1-24. A condundrum in the philosophy of law: Can laws modify
 themselves, or is that paradoxical? Ross' view is that logical inconsistency is unacceptable
 in law, and therefore that self-amendment is impossible. For a response, see the paper by
 Hart (above).
Rucker, Rudy. Infinity and the Mind. Boston: Birkhauser, 1982. A book that had to be
 written. Presenting the most abstruse concoctions of the mind in language that is not
 abstruse, and connecting it with thoughts about consciousness and the mystery of existence
 -this is what Rucker excels in.
Ruelle, David. "Les Attracteurs Etranges". La Recherche 11, no. 108 (February 1980): 131-
 44. A good article relating these wispy mathematical clouds to the physics they are
 supposed to explain, featuring a number of excellent illustrations.
Rumelhart, David E. and Donald A. Norman. "Simulating a Skilled Typist: A Study of
 Skilled Cognitive-Motor Performance". Cognitive Science 6, no. 3 (JulySeptember 1982):
 1-36. Anticipating the current "parallel distributed processing" project at the University of
 California at San Diego, the research described in this article is among the most interesting
 work on modeling human performance that I have encountered.
Russett, Bruce. The Prisoners of Insecurity. New York: W.H. Freeman, 1983. The arms race
 as an iterated Prisoner's Dilemma, and how we might be able to break out of the deadlock.
 The last section-"Responsibility"--concludes with this admonition (and how I wish people
 would take it to heart!): "In a democracy, silence about nuclear issues carries an implication
 not just of indifference but of acceptance. If we stand silent in the face of an arms race-and

Bibliography                                                                               815

  the war to which it may lead us-we must share responsibility for the outcome. `Silence
  gives consent.' "
Ryder, Frederick, and Company. Ryder Types, 2 vols. (with periodic supplements). Chicago:
  Frederick Ryder and Company. The best catalogue of typefaces I have run across-but it is
  expensive. Some of the oddest faces I have ever seen are found in the four supplements I
  own.
Sagan, Carl, ed. Communication with Extraterrestrial intelligence. Cambridge, Mass.: MIT
  Press, 1973.. An entertaining transcript of an international meeting in the days when
  Russians and Americans spoke to each other. Sagan puts forth his notion of various types of
  earth-life-based "chauvinisms". A hopeful and intellectually refreshing book-the kind one
  wishes there were hundreds of.
Sampson, Geoffrey. "Is Roman Type an Open-Ended System? A Response to ,Douglas
  Hofstadter". Visible Language 17, no. 4 (Autumn 1983): 410-12. Sampson tlisputes my
  claim that it is impossible to capture the fluid spirit of letters of the alphabet in parametrized
  computer subroutines; he suggests that it is possible, if you limit your goals to capturing the
  spirit of letters suitable for printing serious books in.
Schank, Roger. Dynamic Memory: A Theory of Reminding and Learning in Computers and
  People. New York: Cambridge University Press, 1982. Although I disagree with some of
  his theorizing, I agree fully with Schank's focus on the types of problems that cognitive
  science ought to be most concerned with, and I like the examples he uses.
Scherlis, William L. and Pierre L. Wolper. "Self-Referenced Referenced, and Self-
  Referenced". Communications of the ACM 23, no. 12 (December 1980): 736. A humorous
  short note on papers that cite themselves.
Schrodinger, Ernst. What Is Life? and Mind and Matter. New York: Cambridge University
  Press, 1967 (reprint of 1944 edition). This philosophically-minded physicist prophetically
  speculates on the nature of the hereditary message, before the days when DNA's structure or
  function were known. Also venturing into the deep waters of consciousness, he comes up
  with this mystical conclusion: "The over-all number of minds is just one."
Schwenk, Theodor. Sensitive Chaos. New York: Schocken, 1976. A book about fluids in the
  wild and in the laboratory, filled with striking patterns and fantastic photographs strongly
  suggesting this mystical conclusion: "The over-all number of fluids is just one."
Searle, John. "Minds, Brains, and Programs". The Behavioral and Brain Sciences 3
  (September 1980): 417-57. Like the Lucas article (see above), this one launched a thousand
  rebuttals (including two by me). In this, its original setting, it was followed by nearly thirty
  rebuttals and counter-rebuttals. It is amusing and educational to read them all. I view this
  article as a litmus test, in the sense that someone convinced by Searle's imagery is almost
  sure to have a very negative opinion of AI.
-----The Myth of the Computer: An Exchange". New York Review of Books (June 24, 1982):
  56-57. Searle portrays Dennett and me as blind advocates of "strong AI"-the notion that "the
  appropriately programmed computer literally has a mind". He resents the fact that such
  views are "well financed and backed by prestigious teams of research workers", and tells
  how he is constantly working toward "the relentless exposure of its preposterousness".
-----The Myth of the Computer". New York Review of Books (April 29, 1982): 3-6. A rather
  negative review of The Mind's I by someone roundly criticized in the book. The one good
  thing Searle does in this review is to stress the central epistemological problem facing AI: to
  explain the nature of the fluid reference, or semanticity, that mental activity exhibits.
Serafini, Luigi. Codex Seraphinianus. Milan: Franco Maria Ricci, 1981. A complete pseudo-
  encyclopedia, in two volumes. The writing system, the page numbering, the weird
  diagrams, and especially the wonderful color illustrations are all products of the cryptic
  mind of Serafini, an Italian architect. Also available in one volume, and more
  inexpensively, from Abbeville Press in New York, N.Y.
Seuss, Dr. On Beyond Zebra. New York: Random House, 1955. Humorous ways of extending
  the alphabet beyond 'z': a metaphor for jumping out of the system (' jootsing").
Simon, Herbert A. "Cognitive Science: The Newest Science of the Artificial". Cognitive
  Science 4, no. 2 (April June 1980): 33-46. Simon's confident assertion that computers

Bibliography                                                                                   816

  already have what it takes to possess full intelligence: symbol manipulation. In his
  conclusion, he claims: "Wherever the boundary is drawn, there exists today a science of
  intelligent systems that extends beyond the limits of any single species."
-----Studying Human Intelligence by Creating Artificial Intelligence". American Scientist 69,
  no. 3 (May June 1981): 300-309. This is the lecture in which Simon spoke of the all-
  important 100-millisecond barrier, below which it is of no interest to cognitive scientists to
  know what happens (see my Chapter 26). In print he changed "100 millisecondsto "ten
  milliseconds", although the point remains the same either way.
Singmaster, David. Notes on Rubik's "Magic Cube". Hillside, NJ.: Enslow, 1981. On its
  cover, this book proudly boasts: "'The definitive treatise.'-Scientific American. " Well, if
  Scientific American says so, who am I to dispute it?
Sloman, Aaron. The Computer Revolution in Philosophy: Philosophy, Science, and Models of
  Mind. Atlantic Highlands, NJ.: Humanities Press, 1978. A book that warns philosophers
  that they will miss the boat if they don't jump on the Al bandwagon. This a good though
  somewhat tendentious discussion of the philosophical import of Al.
Smith, Brian C. "Reflection and Semantics in Lisp". Palo Alto, Calif.: Xerox Palo Alto
  Research Center Report, 1983. A boiling-down of a 500-page Ph.D. thesis into a dozen
  pages or so, concerning a system capable of reasoning about itself (and reasoning about
  such reasoning, etc.). This work exemplifies the "meta-meta" style of Al research (see my
  Chapter 23-especially its P.S. ).
Smith, Stephen B. The Great Mental Calculators: The Psychology, Methods, and Lives of
  Calculating Prodigies. New York: Columbia University Press, 1983. A set of portraits of
  very strange yet very human beings whose aberrant minds let them do calculating tasks that
  we ordinary people-no matter how number-loving-could not conceivably do.
Smolensky, Paul. "Harmony Theory: A Mathematical Framework for Stochastic Parallel
  Processing". University of California at San Diego Institute for Cognitive Science Technical
  Report ICS No. 8306. Based on the ideas of statistical mechanics, this project utilizes
  stochastic parallelism, regulated by a "temperature" that gradually drops to zero, to search
  most efficiently for the optimal global state of a system.
Smullyan, Raymond. This Book Needs No Title: A Budget of Living Paradoxes. Englewood
  Cliffs, NJ.: Prentice-Hall, 1980. A very humorous collection of observations about the
  constant intermingling of life and paradox, by a logician whose awareness of paradox is
  especially keen.
Solo, Dan X. Sans Serif Display Alphabets. New York: Dover, 1979. A collection of elegant
  letterforms whose subdued flair resides exclusively in their gentle curves, stroke taperings,
  line endings, and the interplay between positive and negative space.
-----Special-Effects and Topical Alphabets. New York: Dover, 1978. After you've looked at a
  book like this, you know why the problem of letterforms is synonymous with the problem
  of full human intelligence.
Sonneborn, Tracy M. "Degeneracy of the Genetic Code: Extent, Nature, and Genetic
  Implications". In Evolving Genes and Proteins, edited by Vernon Bryson and
Henry J. Vogel. New York: Academic Press, 1965. Perhaps the first paper to suggest that
  there is evolutionary rhyme and reason to the particular match-up between codons and
  amino acids that exists in the arbitrary-seeming genetic code.
Soppeland, Mark. Words. Los Altos, Calif.: William Kaufmann, 1980. A witty collection of
  words drawn as self-referential pictures: stacks of books whose shapes spell out "books",
  and so on.
Sorrels, Bobbye. The Nonsexist Communicator. Englewood Cliffs, NJ.: Prentice-Hall, 1983.
  A sizable collection of sexist usages and nonsexist remedies; some of the remedies,
  however, are needlessly awkward.
Sperry, Roger. "Mind, Brain, and Humanist Values". In New Views on the Nature of Man,
  edited by John R. Platt. Chicago: University of Chicago Press, 1965. The article that asks,
  and tries to answer, the question "Who pushes whom around in the population of causal
  forces that occupy the cranium?" My answer is in Chapter 25.


Bibliography                                                                               817

Spinelli, Aldo. Loopings. Amsterdam: Multi-Art Points Edition, 1976. A somewhat clumsy
  realization of a good idea: a book whose pages are all self-inventorying sentences, or close
  to that. At the end, Spinelli discusses the kinds of locked-in loops and entryways into
  looping that can arise.
Stanley, H. Eugene et al. "Interpretation of the Unusual Behavior of H2O and D20 at Low
  Temperature: Are Concepts of Percolation Relevant to the 'Puzzle of Liquid Water'?"
  Physica. 106a (1981): 260-77. An article that gives a fairly recent picture of the "flickering-
  cluster" view of water.
Stein, Gertrude. How to I Write. Craftsbury Common, Vt.: Sherry Urie, 1977. (Originally
  published in 1931.) Classic volume of absurdities. It would be a nightmare to try to read the
  whole thing; each page is like a unique grain of sand, but all the pages taken together add up
  to a rather bland beach. Better to savor it in small bits.
Steiner, George. After Babel: Aspects of Language and Translation. New York: Oxford
  University Press, 1975. This fascinating exploration of the depths of personal meaning is
  full of obscurities, but it is still one of the best statements on the subtlety of language that 'I
  know.
Strich, Christian. Fellini's Faces: 418 Photographs from the Archives of Federico Fellini. New
  York: Holt, Rinehart, and Winston, 1981. A celebration of the remarkable diversity of
  human physiognomies, as well as their amazing expressivity.
Stryer, Lubert. Biochemistry. New York: W. H. Freeman, 1975. An elegantly produced text
  of biochemistry, on a slightly lower level than Lehninger. The figures-many in color -are
  especially appealing.
Suber, Peter. "A Bibliography of Works on Reflexivity". Unpublished manuscript, 1984. An
  extensive collection of pointers to the vast universe of writings on self-modifying laws, self-
  referring literature, self-fulfilling prophecies, self-replicating machines, self-monitoring
  computer programs, and on and on.
-----The Paradox of Self-Amendment: A Study of Logic, Law, Omnipotence, and Change.
  Forthcoming. This is the first (and only) book-length study of logical paradoxes in law. It
  focuses on one family of paradoxes (rules that authorize change being applied to
  themselves), but eventually tries to address the general problem of how legal reasoning
  copes with logical paradox. The game of Nomic is featured in an appendix.
Thomas, Dylan. Collected Poems. New York: New Directions, 1953.1 am the first to admit I
  don't understand poetry well, but this beguiling collection, with its combination of beautiful
  language and bewildering opacity, leaves me especially disturbed.
Thomas, Lewis. The Medusa and the Snail. New York: Viking, 1979. Short essays on science
  and life, some of which are very insightful, others of which I find puzzling or just plain
  silly.
Turing, Alan. "Computing Machinery and Intelligence". Reprinted in Minds and Machines,
  edited by Alan Ross Anderson. Englewood Cliffs, NJ.: Prentice-Hall,
1964. In which the now-infamous "Turing Test" for machine thought is proposed. Lucid and
  straightforward, this article contains plenty of provocative ideas, even today.
Turner-Smith, Ronald. The Amazing Pyraminx. Hong Kong: Meffert Novelties, 1981. A
  simple guide to the notation, group theory, and solution of the Pyraminx puzzle.
Ulam, Stanislaw. Adventures of a Mathematician. New York: Scribners, 1976. A fascinating
  account of the intellectual life of a highly innovative and fun-loving mathematician,
  including some speculations on mind and consciousness.
Vetterling-Braggin, Mary, ed. Sexist Language: A Modern Philosophical Analysis. Totowa,
  NJ.: Littlefield, Adams \& Co., 1981. A valuable collection of articles from many
  perspectives on sexism. Includes one section on the comparison between sexism and
  racism.
von Neumann, John. Theory of Self-Reproducing Automata, edited and completed by
  Arthur W. Burks. Urbana, Illinois: University of Illinois Press, 1966. Here von Neumann
  describes how a machine could construct replicas of itself, or even machines more complex
  than itself, out of raw materials. At the core of this work by von Neumann, however, is


Bibliography                                                                                    818

  Gödel’s method of achieving mathematical self-reference-a fact that many people seem to
  overlook or downplay.
Walker, Alan, ed. The Chopin Companion. New York: W. W. Norton, 1966. An excellent
  collection of articles by composers, performers, and musicologists on Chopin, his music,
  and its influence.
Watzlawick, Paul. Change. New York: W. W. Norton, 1974. A theory of psychological
  disorder as caused by encounters with paradox in daily living. The suggested therapy is to
  fight paradox with paradox (if that isn't paradoxical!).
Webb, Judson. Mechanism, Mentalism, and Metamathematics. Hingham, Mass.: D. Reidel,
  1980. A penetrating scholarly analysis of the import of the most important
  metamathematical work by Gddel, Church, Turing, Kleene, Tarski, Post, and others. In
  particular, it seeks to elucidate the relationship of the famous "limitative theorems"
  applicable to formal systems with attempts to mechanize human mental activity. In essence,
  its conclusion is that those results strengthen, rather than diminish, the notion that mental
  activity can in principle be mechanized.
Wells, Carolyn. A Nonsense Anthology. New York: Dover Press, 1958. A wonderful
  collection of strange poems and silly snippets of pseudo-sense.
Wheelis, Allen. The Scheme of Things. New York: Harcourt Brace Jovanovich, 1980. A
  moving story of human and animal attachments, and the emotional resonances of places and
  scenes. The story hints at many self-references, most notably to another novel called The
  Way Things Are, written by Wheelis' protagonist Oliver Thompson who, like Wheelis
  himself, is an unorthodox psychiatrist living in San Francisco. Wheelis within Wheelis!
Williams, J. D. The Compleat Strategyst. New York: McGraw-Hill, 1954. One of the earliest
  and clearest books ever written on game theory. It is a very elementary book, with
  humorous drawings and droll stories to illustrate its points.
Wilson, E. O. The Insect Societies. Cambridge, Mass.: Harvard University, Belknap Press,
  1971. It is heavy going to read this book, but parts of it are engrossing, as they tell how
  order at a high level comes from the independent activity of low-level organisms acting
  without knowledge of one another.
Winston, Patrick Henry. The Psychology of Computer Vision. New York: McGraw-Hill,
  1975. Although fairly old by AI standards, this collection of six articles is still of
  considerable interest. One is Minsky's article on frames; another is an article by Winston on
  learning and recognition; and another is an article by Waltz on economical strategies for
  vision in the "blocks world".
Winston, Patrick Henry and Berthold Horn. Lisp. Reading, Mass.: Addison-Wesley, . 1981. A
  solid textbook on Lisp, starting right at' the beginning and going up to AI programming
  techniques.
Yaguello, Marina. Alice au pays du langage. Paris: Editions du Seuil, 1981. A French linguist
  examines the phonetics, syntax, semantics, and gradual shifts of her own language to
  illuminate the strangeness of human thought. Many of the examples in the book are taken
  from popular language, slang, or humor, which makes it particularly unstuffy.
-----Les mots et les femmes. Paris: Petite Bibliotheque Payot, 1978. On fighting the battle
  against sexist language in French-but unfortunately, because of the two-gender system of
  romance languages, eradication of sexism will prove far more difficult in those languages
  than in English (where it is already ferociously hard).
Zapf, Hermann. About Alphabets. Cambridge, Mass.: MIT Press, 1960. A famous
  contemporary type designer describes his adventures among the alphabets of the world.




Bibliography                                                                              819

